% Production d’écrits utilitaires : la requête — Français Tle Tech — Cours signé OnBuch+
% Programme officiel MINESEC. Compiler avec : tectonic N11-production-ecrits-utilitaires-requete.tex
\newcommand{\DOCMATIERE}{Français}
\newcommand{\DOCNIVEAU}{Tle Tech}
\newcommand{\DOCMODULE}{Français – classe de Terminale technique}
\newcommand{\DOCLECON}{N11}
\newcommand{\DOCTITRE}{Production d’écrits utilitaires : la requête}
% OnBuch+ — préambule « cours signés » ENRICHI (compilé avec tectonic / XeLaTeX)
% Système visuel « Pop » d'OnBuch+ : fond crème (jamais blanc), bordures encre
% épaisses, ombres « dures » décalées, titres Archivo Black en capitales.
% Version MATHÉMATIQUES : graphes (pgfplots/tikz), boîtes pédagogiques
% (définition, propriété, méthode, attention, le savais-tu, exercice résolu,
% exercice, TP, bilan), page de garde et signature OnBuch+ ancrée en bas.
%
% Macros attendues dans main.tex AVANT \input{preamble} :
%   \DOCMATIERE  \DOCNIVEAU  \DOCMODULE  \DOCLECON (numéro)  \DOCTITRE
\documentclass[11pt]{article}

\usepackage{fontspec}
\usepackage[french]{babel}
\usepackage[a4paper,margin=2cm,top=3.4cm,bottom=2.8cm,headheight=1.4cm,footskip=1.3cm]{geometry}
\usepackage{amsmath,amssymb,amsfonts}
\usepackage{mathtools} % \xmapsto, \coloneqq... souvent utilisés par les modèles
\usepackage{cancel} % simplifications barrées (bilans de forces)
% \abs{z} (module, valeur absolue) et \norm{u} (norme), utilisés par les modèles
\providecommand{\abs}{}\let\abs\relax\DeclarePairedDelimiter{\abs}{\lvert}{\rvert}
\providecommand{\norm}{}\let\norm\relax\DeclarePairedDelimiter{\norm}{\lVert}{\rVert}
\usepackage[table]{xcolor} % [table] : fournit aussi \rowcolors (alternance de lignes)
\usepackage{pagecolor}
\usepackage{tikz}
\usetikzlibrary{arrows.meta,positioning,calc,shapes.geometric,shapes.misc,shapes.arrows,decorations.pathmorphing,decorations.pathreplacing,decorations.markings,patterns,shadows,mindmap,trees,fit,backgrounds,matrix,intersections,angles,quotes}
\usepackage{pgfplots}
\usepackage[version=4]{mhchem} % équations nucléaires \ce{^{235}_{92}U}
\usepackage[european,siunitx]{circuitikz} % schémas électriques
\pgfplotsset{compat=1.18}
\usepgfplotslibrary{fillbetween,groupplots} % aires entre courbes (calcul intégral)
\usepackage{enumitem}
\usepackage{fancyhdr}
% [most] charge toutes les bibliothèques tcolorbox (listings, minted, raster,
% documentation...) et gonfle inutilement la mémoire TeX sur un document long
% (~300 boîtes) : on ne charge que ce qui est réellement utilisé.
\usepackage[skins,breakable]{tcolorbox}
\usepackage{graphicx}
\usepackage{colortbl}
\usepackage{booktabs}
\usepackage{tabularx}
\usepackage{multirow}
\usepackage{diagbox} % case d'en-tête diagonale des tableaux à double entrée
% Colonnes extensibles centrée / gauche / droite (C, L, R), souvent utilisées par les modèles
\newcolumntype{C}{>{\centering\arraybackslash}X}
\newcolumntype{L}{>{\raggedright\arraybackslash}X}
\newcolumntype{R}{>{\raggedleft\arraybackslash}X}
\usepackage{array}
\usepackage{multicol}
\usepackage{siunitx}
% parse-numbers=false : accepte aussi \SI{2,5 \times 10^{-3}}{...}
\sisetup{locale=FR,output-decimal-marker={,},group-minimum-digits=4,parse-numbers=false}
\DeclareSIUnit\ohm{\ensuremath{\symup{\Omega}}} % U+2126 absent de la police
\DeclareSIUnit\division{div} % réglages d'oscilloscope (ms/div, V/div)
\DeclareSIUnit\tr{tr} % tours (vitesse de rotation, ex. \SI{1500}{\tr\per\minute})
\DeclareSIUnit\molar{M} % concentration molaire (ex. \SI{3}{\molar}, \SI{1}{\micro\molar})
\DeclareSIUnit\an{an} % année (le modèle confond parfois avec \year, un compteur TeX) — ex. \SI{5}{\kilo\meter\per\an}
\DeclareSIUnit\FCFA{FCFA} % franc CFA (ex. \SI{500}{\FCFA\per\kilo\gram})
\DeclareSIUnit\degreeBrix{\textdegree Brix} % teneur en sucre d'un sirop/jus (réfractométrie)
\DeclareSIUnit\month{mois} % le modèle utilise parfois \month, un compteur TeX, comme unité de durée
\DeclareSIUnit\microliter{\micro\liter} % le modèle écrit parfois \microliter en un seul token
\DeclareSIUnit\million{millions} % dénombrement (ex. \SI{15}{\million\per\mL} spermatozoïdes)
\usepackage{qrcode}
% \degree : commande fréquemment utilisée par les modèles pour un angle ou une
% température en dehors de \SI{}{} (non fournie par les paquets chargés).
\providecommand{\degree}{^{\circ}}
% \qed (fin de démonstration), utilisé par les modèles sans amsthm
\providecommand{\qed}{\hfill$\square$}
% \barre : double barre des valeurs interdites dans un tableau de variations (tkz-tab)
\providecommand{\barre}{\ensuremath{\|}}
% environnement proof (amsthm non chargé) utilisé par les modèles
\ifdefined\proof\else\newenvironment{proof}[1][Démonstration]{\par\noindent\textit{#1.}\ }{\qed\par}\fi
\usepackage{lastpage}
\usepackage{hyperref}
\hypersetup{hidelinks}

% ============================================================
% Palette « Pop » OnBuch+ — mêmes hex que la classe P (plus_ui.dart)
% ============================================================
\definecolor{poporange}{HTML}{F59321}
\definecolor{popdark}{HTML}{E07A0C}
\definecolor{popcream}{HTML}{FFF6E8}
\definecolor{popink}{HTML}{1C1714}
\definecolor{poppurple}{HTML}{7A5AE0}
\definecolor{popgreen}{HTML}{2E9E6A}
\definecolor{poppink}{HTML}{E0457B}
\definecolor{popblue}{HTML}{2F7DE1}
\definecolor{popgold}{HTML}{C98A12}
\definecolor{popmuted}{HTML}{8A7F76}
\definecolor{popline}{HTML}{EADFD2}
\definecolor{popgray}{HTML}{8A7F76}
\colorlet{popgrey}{popgray}
% Alias tolérants (le modèle nomme parfois les couleurs différemment)
\colorlet{popwhite}{white}
\colorlet{popblack}{popink}
\colorlet{popred}{poppink}
\colorlet{popyellow}{popgold}
% Teintes claires (fonds de tableaux/encadrés) — le modèle nomme parfois
% les couleurs avec un suffixe L (ex. popblueL) sans les avoir définies.
\colorlet{poporangeL}{poporange!20}
\colorlet{popdarkL}{popdark!20}
\colorlet{poppurpleL}{poppurple!20}
\colorlet{popgreenL}{popgreen!20}
\colorlet{poppinkL}{poppink!20}
\colorlet{popblueL}{popblue!20}
\colorlet{popgoldL}{popgold!20}
\colorlet{popredL}{popred!20}
\colorlet{popgrayL}{popgray!20}
% Teintes claires pour fonds de tableaux / zones
\colorlet{popblueL}{popblue!10!white}
\colorlet{poporangeL}{poporange!14!white}
\colorlet{popgreenL}{popgreen!12!white}
\colorlet{poppurpleL}{poppurple!12!white}
\colorlet{poppinkL}{poppink!12!white}
\colorlet{popgoldL}{popgold!14!white}

\pagecolor{popcream}

% ============================================================
% Typographie
% ============================================================
\defaultfontfeatures{Ligatures=TeX}
\setmainfont{PlusJakartaSans-Medium.ttf}[Path=../../../fonts/,BoldFont=PlusJakartaSans-Bold.ttf]
% Maths en sans-serif (Fira Math) avec chiffres et lettres droites de Plus
% Jakarta, pour que les formules se fondent dans le texte.
\usepackage[math-style=ISO,bold-style=ISO]{unicode-math}
\setmathfont{FiraMath-Regular.otf}[Path=../../../fonts/]
\setmathfont{PlusJakartaSans-Medium.ttf}[Path=../../../fonts/,range={up/num,up/{Latin,latin},\mathup}]
\usepackage{icomma} % virgule décimale sans espace en mode math
% Caractères mathématiques Unicode absents des polices de texte (Plus Jakarta,
% Archivo Black) : rendus par leur équivalent mathématique, en texte comme en
% formule — sinon ils s'affichent en carré vide (ex. « Arithmétique dans ℤ »).
\usepackage{newunicodechar}
\newunicodechar{ℝ}{\ensuremath{\mathbb{R}}}
\newunicodechar{ℤ}{\ensuremath{\mathbb{Z}}}
\newunicodechar{ℂ}{\ensuremath{\mathbb{C}}}
\newunicodechar{ℕ}{\ensuremath{\mathbb{N}}}
\newunicodechar{ℚ}{\ensuremath{\mathbb{Q}}}
\newunicodechar{𝔻}{\ensuremath{\mathbb{D}}}
\newunicodechar{α}{\ensuremath{\alpha}}
\newunicodechar{β}{\ensuremath{\beta}}
\newunicodechar{γ}{\ensuremath{\gamma}}
\newunicodechar{δ}{\ensuremath{\delta}}
\newunicodechar{ε}{\ensuremath{\varepsilon}}
\newunicodechar{θ}{\ensuremath{\theta}}
\newunicodechar{λ}{\ensuremath{\lambda}}
\newunicodechar{μ}{\ensuremath{\mu}}
\newunicodechar{σ}{\ensuremath{\sigma}}
\newunicodechar{τ}{\ensuremath{\tau}}
\newunicodechar{φ}{\ensuremath{\varphi}}
\newunicodechar{ψ}{\ensuremath{\psi}}
\newunicodechar{ω}{\ensuremath{\omega}}
\newunicodechar{Σ}{\ensuremath{\Sigma}}
% majuscules produites par \MakeUppercase dans les titres (α-aminés -> Α)
\newunicodechar{Α}{\ensuremath{\alpha}}
\newunicodechar{Β}{\ensuremath{\beta}}
\newunicodechar{Γ}{\ensuremath{\Gamma}}
\newunicodechar{Δ}{\ensuremath{\Delta}}
\newunicodechar{Ε}{\ensuremath{\varepsilon}}
\newunicodechar{Θ}{\ensuremath{\Theta}}
\newunicodechar{Λ}{\ensuremath{\Lambda}}
\newunicodechar{Μ}{\ensuremath{\mu}}
\newunicodechar{Τ}{\ensuremath{\tau}}
\newunicodechar{Φ}{\ensuremath{\Phi}}
\newunicodechar{Ψ}{\ensuremath{\Psi}}
\newunicodechar{Ω}{\ensuremath{\Omega}}
\newunicodechar{↦}{\ensuremath{\mapsto}}
\newunicodechar{∈}{\ensuremath{\in}}
\newunicodechar{∉}{\ensuremath{\notin}}
\newunicodechar{∧}{\ensuremath{\wedge}}
\newunicodechar{⇔}{\ensuremath{\Leftrightarrow}}
\newunicodechar{⟺}{\ensuremath{\Longleftrightarrow}}
\newunicodechar{⇒}{\ensuremath{\Rightarrow}}
\newunicodechar{⟹}{\ensuremath{\Longrightarrow}}
\newunicodechar{∘}{\ensuremath{\circ}}
\newunicodechar{∪}{\ensuremath{\cup}}
\newunicodechar{∩}{\ensuremath{\cap}}
\newunicodechar{≡}{\ensuremath{\equiv}}
\newunicodechar{⊂}{\ensuremath{\subset}}
\newunicodechar{⊥}{\ensuremath{\perp}}
\newunicodechar{∃}{\ensuremath{\exists}}
\newunicodechar{∀}{\ensuremath{\forall}}
\newunicodechar{∛}{\ensuremath{\sqrt[3]{\,}}}
\newunicodechar{∜}{\ensuremath{\sqrt[4]{\,}}}
\newunicodechar{⃗}{\ensuremath{{}^{\to}}}
\newunicodechar{ⁿ}{\ifmmode^{n}\else\textsuperscript{\ensuremath{n}}\fi}
\newunicodechar{ˣ}{\ifmmode^{x}\else\textsuperscript{\ensuremath{x}}\fi}
\newunicodechar{ᵇ}{\ifmmode^{b}\else\textsuperscript{\ensuremath{b}}\fi}
\newunicodechar{ʳ}{\ifmmode^{r}\else\textsuperscript{\ensuremath{r}}\fi}
\newunicodechar{ᵘ}{\ifmmode^{u}\else\textsuperscript{\ensuremath{u}}\fi}
\newunicodechar{ʸ}{\ifmmode^{y}\else\textsuperscript{\ensuremath{y}}\fi}
\newunicodechar{ᵃ}{\ifmmode^{a}\else\textsuperscript{\ensuremath{a}}\fi}
\newunicodechar{ᵉ}{\ifmmode^{e}\else\textsuperscript{\ensuremath{e}}\fi}
\newunicodechar{ᵏ}{\ifmmode^{k}\else\textsuperscript{\ensuremath{k}}\fi}
\newunicodechar{ᵖ}{\ifmmode^{p}\else\textsuperscript{\ensuremath{p}}\fi}
\newunicodechar{ᴺ}{\ifmmode^{N}\else\textsuperscript{\ensuremath{N}}\fi}
\newunicodechar{ᶜ}{\ifmmode^{c}\else\textsuperscript{\ensuremath{c}}\fi}
\newunicodechar{ʰ}{\ifmmode^{h}\else\textsuperscript{\ensuremath{h}}\fi}
\newunicodechar{ᵠ}{\ifmmode^{q}\else\textsuperscript{\ensuremath{q}}\fi}
\newunicodechar{ₙ}{\ifmmode_{n}\else\textsubscript{\ensuremath{n}}\fi}
\newunicodechar{ₐ}{\ifmmode_{a}\else\textsubscript{\ensuremath{a}}\fi}
\newunicodechar{ᵢ}{\ifmmode_{i}\else\textsubscript{\ensuremath{i}}\fi}
\newunicodechar{ᵣ}{\ifmmode_{r}\else\textsubscript{\ensuremath{r}}\fi}
\newunicodechar{ₖ}{\ifmmode_{k}\else\textsubscript{\ensuremath{k}}\fi}
\newunicodechar{ᵦ}{\ifmmode_{\beta}\else\textsubscript{\ensuremath{\beta}}\fi}
\newunicodechar{ₓ}{\ifmmode_{x}\else\textsubscript{\ensuremath{x}}\fi}
\newfontfamily\popdisplay{ArchivoBlack-Regular.ttf}[Path=../../../fonts/]
\newfontfamily\poplogo{SpaceGrotesk-Bold.ttf}[Path=../../../fonts/]
\newfontfamily\popsemi{PlusJakartaSans-SemiBold.ttf}[Path=../../../fonts/]
\newfontfamily\popxbold{PlusJakartaSans-ExtraBold.ttf}[Path=../../../fonts/]
\newfontfamily\popxboldls{PlusJakartaSans-ExtraBold.ttf}[Path=../../../fonts/,LetterSpace=3.0]

\setlist[enumerate]{leftmargin=1.7em,itemsep=5pt,topsep=4pt}
\setlist[itemize]{leftmargin=1.7em,itemsep=4pt,topsep=4pt,label={\color{poporange}\textbullet}}
\setlength{\parindent}{0pt}
\setlength{\parskip}{5pt}
\renewcommand{\arraystretch}{1.25}

% pgfplots : style Pop par défaut
\pgfplotsset{
  every axis/.append style={axis line style={popink,line width=1pt},tick style={popink},
    grid style={popline,line width=0.5pt},label style={font=\small},tick label style={font=\footnotesize}},
  popcurve/.style={poporange,line width=1.8pt,no markers,smooth},
}
% Même style disponible dans un \draw TikZ simple (hors pgfplots)
\tikzset{popcurve/.style={poporange,line width=1.8pt}}
\tikzset{popgrid/.style={popline,very thin}}
\colorlet{popcurve}{poporange} % le nom du style est parfois utilisé comme couleur

% ============================================================
% Pill / squiggle
% ============================================================
\newcommand{\poppill}[3]{%
  \tikz[baseline=-0.55ex]\node[draw=popink,line width=1.2pt,rounded corners=2.6pt,%
    fill=#2,inner xsep=6.5pt,inner ysep=3pt]{%
    \popxboldls\fontsize{7}{7}\selectfont\color{#3}\MakeUppercase{#1}};%
}
\newcommand{\popsquiggle}[1]{%
  \tikz[baseline=-0.4ex]{
    \draw[#1,line width=2.6pt,line cap=round]
      plot [smooth] coordinates {(0,0.09) (0.34,0.01) (0.68,0.21) (1.02,0.01) (1.36,0.10)};
  }%
}
% Mot-clé surligné (marqueur orange sous le texte)
\newcommand{\cle}[1]{{\popxbold\color{popdark}#1}}

% ============================================================
% En-tête / pied
% ============================================================
\pagestyle{fancy}
\fancyhf{}
\renewcommand{\headrulewidth}{0pt}
\renewcommand{\footrulewidth}{0pt}
\lhead{\raisebox{-0.15ex}{\poplogo\fontsize{13}{13}\selectfont\color{popink}OnBuch\color{poporange}+}}
\rhead{\poppill{\DOCMATIERE\ \textbullet\ \DOCNIVEAU\ \textbullet\ Leçon \DOCLECON}{white}{popink}}
\lfoot{\popxbold\fontsize{7.6}{7.6}\selectfont\color{popmuted}\MakeUppercase{Cours signé OnBuch+}\ \textbullet\ {\popsemi\DOCTITRE}}
\rfoot{\tikz[baseline=-0.6ex]\node[rounded corners=8.5pt,fill=popink,minimum height=17pt,minimum width=17pt,inner xsep=5pt,inner sep=0]{\popxbold\fontsize{8}{8}\selectfont\color{popcream}\thepage\,/\,\pageref*{LastPage}};}

% ============================================================
% Titres
% ============================================================
\newcommand{\chaptitre}[2]{%
  \begin{tcolorbox}[enhanced,colback=poporange,colframe=popink,boxrule=2.6pt,arc=11pt,
      drop shadow=popink,left=15pt,right=15pt,top=12pt,bottom=11pt,before skip=3pt,after skip=18pt]
    \poppill{#2}{popink}{popcream}\par\vspace{9pt}
    {\raggedright\hyphenpenalty=10000\exhyphenpenalty=10000\popdisplay\fontsize{22}{24}\selectfont\color{popcream}\MakeUppercase{#1}\par}
  \end{tcolorbox}
}
% Section numérotée : pastille orange avec le numéro + titre Archivo
\newcounter{popsec}
\newcommand{\coursec}[1]{%
  \stepcounter{popsec}\setcounter{popsub}{0}%
  \par\vspace{16pt}\noindent
  \begin{minipage}{\linewidth}
  \tikz[baseline=-0.7ex]\node[draw=popink,line width=1.6pt,fill=poporange,rounded corners=4pt,minimum size=22pt,inner sep=0]{\popdisplay\fontsize{12}{12}\selectfont\color{popcream}\thepopsec};%
  \hspace{8pt}{\popdisplay\fontsize{15}{17}\selectfont\color{popink}\MakeUppercase{#1}}\par
  \vspace{4pt}\noindent{\color{poporange}\rule{36pt}{3.6pt}}
  \end{minipage}\par\nopagebreak\vspace{8pt}
}
\newcounter{popsub}
\newcommand{\courssub}[1]{%
  \stepcounter{popsub}%
  \par\vspace{9pt}\noindent{\popxbold\fontsize{11.5}{13}\selectfont\color{popink}{\color{poporange}\thepopsec.\thepopsub}\hspace{6pt}#1}\par\nopagebreak\vspace{4pt}
}

% ============================================================
% Boîtes pédagogiques — même recette : fond blanc, bordure épaisse colorée,
% ombre dure encre, badge plein. Titre optionnel [#1].
% ============================================================
\tcbset{popbase/.style={breakable,enhanced,colback=white,boxrule=2.3pt,arc=9pt,
  shadow={3pt}{-3pt}{0pt}{popink},left=14pt,right=14pt,top=11pt,bottom=11pt,before skip=11pt,after skip=15pt,
  fontupper=\popsemi\fontsize{10.6}{14.5}\selectfont\color{popink},
  fontlower=\popsemi\fontsize{10.6}{14.5}\selectfont\color{popink}}}
% Coche dessinée (le glyphe ✓ n'existe pas dans les polices du document)
\newcommand{\popcheck}[1]{\tikz[baseline=-0.5ex]\draw[#1,line width=1.6pt,line cap=round,line join=round](-0.13,0.0)--(-0.04,-0.09)--(0.13,0.1);}
\providecommand{\coche}{\popcheck{popgreen}}
\providecommand{\resultat}[1]{\par\smallskip\noindent\fcolorbox{popgreen}{popgreenL}{\parbox{\dimexpr\linewidth-2\fboxsep-2\fboxrule\relax}{\textbf{Résultat.} #1}}\par\smallskip}
\newcommand{\popboxhead}[3]{\poppill{#1}{#2}{white}\ifx\relax#3\relax\else\hspace{6pt}{\popxbold\fontsize{10.6}{12}\selectfont\color{popink}#3}\fi\par\vspace{7pt}}

\newtcolorbox{aretenir}[1][]{popbase,colframe=popgreen,before upper={\popboxhead{À retenir}{popgreen}{#1}}}
\newtcolorbox{exemplebox}[1][]{popbase,colframe=poporange,before upper={\popboxhead{Exemple}{poporange}{#1}}}
\newtcolorbox{definition}[1][]{popbase,colframe=popblue,before upper={\popboxhead{Définition}{popblue}{#1}}}
\newtcolorbox{propriete}[1][]{popbase,colframe=poppurple,before upper={\popboxhead{Propriété}{poppurple}{#1}}}
\newtcolorbox{methode}[1][]{popbase,colframe=popgold,before upper={\popboxhead{Méthode}{popgold}{#1}}}
\newtcolorbox{attention}[1][]{popbase,colframe=poppink,before upper={\popboxhead{Attention}{poppink}{#1}}}
\newtcolorbox{experience}[1][]{popbase,colframe=popblue,colback=popblueL,before upper={\popboxhead{Expérience}{popblue}{#1}}}
% « Le savais-tu ? » : carte violette pleine (contexte, culture, Cameroun)
\newtcolorbox{savaistu}[1][]{popbase,colback=poppurple,colframe=popink,
  fontupper=\popsemi\fontsize{10.4}{14.2}\selectfont\color{white},
  before upper={\poppill{Le savais-tu\,?}{white}{poppurple}\ifx\relax#1\relax\else\hspace{6pt}{\popxbold\color{white}#1}\fi\par\vspace{7pt}}}
% Exercice résolu : énoncé en haut, \tcblower puis la solution sur fond crème
\newcounter{popexo}\newcounter{popexr}
\newtcolorbox{exoresolu}[1][]{popbase,colframe=poporange,colbacklower=popcream,
  segmentation style={popink,line width=1.2pt,dashed},
  before upper={\stepcounter{popexr}\popboxhead{Exercice résolu \thepopexr}{poporange}{#1}},
  before lower={\poppill{Solution}{popink}{popcream}\par\vspace{6pt}}}
% Exercice (non résolu en place) — la correction est dans la partie Corrigés
\newtcolorbox{exercice}[1][]{popbase,colframe=popink,
  before upper={\stepcounter{popexo}\popboxhead{Exercice \thepopexo}{popink}{#1}}}
% Corrigé
\newtcolorbox{corrige}[1][]{popbase,colframe=popgreen,colback=popgreenL,
  before upper={\popboxhead{Corrigé}{popgreen}{#1}}}
% Objectifs / prérequis / bilan
\newtcolorbox{objectifs}{popbase,colframe=popink,colback=poporangeL,
  before upper={\popboxhead{Objectifs de la leçon}{popink}{}}}
\newtcolorbox{prerequis}{popbase,colframe=popink,colback=popblueL,
  before upper={\popboxhead{Prérequis}{popblue}{}}}
\newtcolorbox{bilan}{popbase,colframe=popink,colback=white,boxrule=2.8pt,
  before upper={\popboxhead{Fiche bilan}{poporange}{L'essentiel à connaître pour l'examen}}}
% Figure encadrée avec légende
\newtcolorbox{popfigure}[1][]{enhanced,breakable=false,colback=white,colframe=popline,boxrule=1.4pt,arc=8pt,
  left=10pt,right=10pt,top=10pt,bottom=8pt,before skip=10pt,after skip=12pt,halign=center,
  lower separated=false,halign lower=center,
  fontlower=\popsemi\fontsize{9}{11.5}\selectfont\color{popmuted},
  #1}
\newcommand{\legende}[1]{\par\vspace{6pt}{\popsemi\fontsize{9}{11.5}\selectfont\color{popmuted}#1}}

% ============================================================
% Signature OnBuch+
% ============================================================
\newcommand{\poplogoinline}[1]{{\poplogo\fontsize{#1}{#1}\selectfont\color{popink}OnBuch\color{poporange}+}}
\newcommand{\signatureonbuch}{%
  \begin{tcolorbox}[enhanced,colback=popink,colframe=popink,boxrule=2.2pt,arc=10pt,
      drop shadow=poporange,left=14pt,right=14pt,top=10pt,bottom=10pt,before skip=0pt,after skip=0pt]
    \tikz[baseline=-0.7ex]\node[circle,fill=popgreen,draw=popcream,line width=1.3pt,minimum size=22pt,inner sep=0]{\popcheck{white}};%
    \hspace{9pt}\begin{minipage}[c]{0.62\linewidth}
      {\popxbold\fontsize{12}{13}\selectfont\color{popcream}Signé\ \textbullet\ {\poplogo OnBuch\color{poporange}+}}\par\vspace{2pt}
      {\raggedright\popsemi\fontsize{8}{10.5}\selectfont\color{popline}\MakeUppercase{Équipe pédagogique OnBuch+}\\\MakeUppercase{Conforme au programme officiel MINESEC}\par}
    \end{minipage}\hfill
    {\poplogo\fontsize{22}{22}\selectfont\color{popcream}OnBuch\color{poporange}+}
  \end{tcolorbox}%
}

% ============================================================
% Page de garde
% #1 titre · #2 matière · #3 niveau · #4 module · #5 n° leçon · #6 durée ·
% #7 description / objectifs (texte ou itemize)
% ============================================================
\newcommand{\pagegarde}[7]{%
  \thispagestyle{empty}
  \newgeometry{a4paper,margin=1.7cm,top=1.6cm,bottom=1.4cm}%
  \begin{tikzpicture}[remember picture,overlay]
    \fill[poporange!18] ([xshift=-3.2cm,yshift=-3.6cm]current page.north east) circle (4.6cm);
    \fill[poppurple!14] ([xshift=2.4cm,yshift=7.6cm]current page.south west) circle (3.8cm);
  \end{tikzpicture}%
  {\poplogo\fontsize{30}{30}\selectfont\color{popink}OnBuch\color{poporange}+}\hfill
  \raisebox{0.6ex}{\poppill{Cours signé \textbullet\ Premium}{popink}{popcream}}\par
  \vspace{1pt}\popsquiggle{poppurple}\par
  \vspace{8pt}
  \begin{tcolorbox}[enhanced,colback=poporange,colframe=popink,boxrule=2.8pt,arc=13pt,
      drop shadow=popink,left=18pt,right=18pt,top=12pt,bottom=13pt,after skip=10pt]
    \poppill{#2 \textbullet\ #3}{popink}{popcream}\hspace{5pt}\poppill{#4}{white}{popink}\par\vspace{8pt}
    {\popdisplay\fontsize{14}{14}\selectfont\color{popink}LEÇON #5}\par\vspace{4pt}
    {\raggedright\hyphenpenalty=10000\exhyphenpenalty=10000\popdisplay\fontsize{25}{27}\selectfont\color{popcream}\MakeUppercase{#1}\par}
    \vspace{8pt}
    {\popxbold\fontsize{9.5}{9.5}\selectfont\color{popink}\MakeUppercase{Durée conseillée : #6}}
  \end{tcolorbox}
  \begin{tcolorbox}[enhanced,colback=white,colframe=popink,boxrule=2.2pt,arc=10pt,
      drop shadow=popink,left=16pt,right=16pt,top=10pt,bottom=10pt,after skip=8pt]
    \poppill{Dans cette leçon}{poppurple}{white}\par\vspace{5pt}
    \popsemi\fontsize{9.3}{12.6}\selectfont\color{popink}#7
  \end{tcolorbox}
  \noindent\begin{minipage}[t]{0.487\linewidth}
    \begin{tcolorbox}[enhanced,colback=popgreen,colframe=popink,boxrule=2.2pt,arc=10pt,
        drop shadow=popink,left=11pt,right=11pt,top=10pt,bottom=10pt,height=3.1cm,valign=center]
      \tcbox[colback=white,colframe=popink,boxrule=1.3pt,arc=5pt,boxsep=0pt,nobeforeafter,
        left=3pt,right=3pt,top=3pt,bottom=3pt,valign=center]{\qrcode[height=1.9cm]{https://play.google.com/store/apps/details?id=com.luvvix.onbuch&hl=fr}}%
      \hspace{8pt}\begin{minipage}[c]{0.5\linewidth}
        {\popdisplay\fontsize{11}{12.5}\selectfont\color{white}\MakeUppercase{Télécharge OnBuch}}\par\vspace{4pt}
        {\raggedright\popsemi\fontsize{8.4}{11}\selectfont\color{white}Gratuit \textbullet\ Google Play\\Tuteur IA, annales, concours, résultats\par}
      \end{minipage}
    \end{tcolorbox}
  \end{minipage}\hfill
  \begin{minipage}[t]{0.487\linewidth}
    \begin{tcolorbox}[enhanced,colback=poppurple,colframe=popink,boxrule=2.2pt,arc=10pt,
        drop shadow=popink,left=11pt,right=11pt,top=10pt,bottom=10pt,height=3.1cm,valign=center]
      \tikz[baseline=-0.7ex]\node[circle,fill=white,draw=popink,line width=1.3pt,minimum size=1.6cm,inner sep=0]{\popdisplay\fontsize{24}{24}\selectfont\color{poppurple}+};%
      \hspace{8pt}\begin{minipage}[c]{0.6\linewidth}
        {\popdisplay\fontsize{11}{12.5}\selectfont\color{white}\MakeUppercase{Passe à OnBuch+}}\par\vspace{4pt}
        {\raggedright\popsemi\fontsize{8.4}{11}\selectfont\color{white}Tous les cours signés, Tuteur IA illimité, fiches PDF — onglet OnBuch+ de l'app\par}
      \end{minipage}
    \end{tcolorbox}
  \end{minipage}
  \vspace{8pt}
  \popcontenu
  \vfill
  \signatureonbuch
  \restoregeometry
  \newpage
}

% Rangée « Ce cours contient » (page de garde)
\newcommand{\poptile}[3]{% #1 couleur #2 gros texte #3 légende
  \begin{tcolorbox}[enhanced,colback=#1,colframe=popink,boxrule=2pt,arc=9pt,shadow={2.5pt}{-2.5pt}{0pt}{popink},
      left=6pt,right=6pt,top=9pt,bottom=9pt,halign=center,valign=center,height=1.75cm,width=\linewidth,nobeforeafter]
    {\popdisplay\fontsize{12.5}{13}\selectfont\color{white}\MakeUppercase{#2}}\par\vspace{3pt}
    {\centering\popsemi\fontsize{7.8}{9.5}\selectfont\color{white}#3\par}
  \end{tcolorbox}}
\newcommand{\popcontenu}{%
  \par\noindent{\popxboldls\fontsize{7.5}{7.5}\selectfont\color{popmuted}CE COURS SIGNÉ CONTIENT}\par\vspace{7pt}
  \noindent\begin{minipage}[t]{0.235\linewidth}\poptile{popblue}{Cours}{complet, illustré, pas à pas}\end{minipage}\hfill
  \begin{minipage}[t]{0.235\linewidth}\poptile{poporange}{Méthodes}{et exercices résolus}\end{minipage}\hfill
  \begin{minipage}[t]{0.235\linewidth}\poptile{poppink}{Exercices}{type examen + corrigés}\end{minipage}\hfill
  \begin{minipage}[t]{0.235\linewidth}\poptile{popgreen}{Fiche bilan}{l'essentiel à retenir}\end{minipage}\par}

% Sommaire
\newtcolorbox{sommaire}{enhanced,colback=white,colframe=popink,boxrule=2.2pt,arc=10pt,
  drop shadow=popink,left=18pt,right=18pt,top=13pt,bottom=13pt,before skip=6pt,after skip=20pt,
  before upper={\poppill{Sommaire}{poppurple}{white}\par\vspace{9pt}\popsemi\fontsize{10.6}{14.5}\selectfont\color{popink}}}

% Signature de fin de cours — poussée tout en bas de la dernière page
\newcommand{\findecours}{%
  \par\vfill\nopagebreak
  \begin{center}\popsquiggle{poporange}\end{center}\vspace{4pt}
  \signatureonbuch
}
\AtBeginDocument{\def\llbracket{\mathopen{[\mkern-3.5mu[}}\def\rrbracket{\mathclose{]\mkern-3.5mu]}}} % intervalles d'entiers (glyphe ⟦ absent de la police)
% Glyphes absents de Fira Math : redéfinis à partir de glyphes disponibles
\AtBeginDocument{%
  \def\setminus{\mathbin{\backslash}}\let\smallsetminus\setminus
  \def\vdots{\mathord{\vbox{\baselineskip3.2pt\lineskiplimit0pt\kern4pt\hbox{.}\hbox{.}\hbox{.}}}}%
  \def\ddots{\mathinner{\mkern1mu\raise6.5pt\hbox{.}\mkern2mu\raise3.6pt\hbox{.}\mkern2mu\raise0.7pt\hbox{.}\mkern1mu}}%
  \def\triangle{\mathord{\scalebox{0.9}{$\symup{\Delta}$}}}%
  \def\nabla{\mathord{\raisebox{\depth}{\scalebox{1}[-1]{$\symup{\Delta}$}}}}%
  \def\checkmark{\popcheck{popdark}}%
  \def\star{\mathbin{\ast}}\def\bigstar{\mathord{\ast}}%
  \def\diamond{\mathbin{\scalebox{0.6}{$\Diamond$}}}%
  \def\bigcup{\mathop{\mathchoice{\raisebox{-0.3em}{\scalebox{1.7}{$\cup$}}}{\scalebox{1.25}{$\cup$}}{\cup}{\cup}}\displaylimits}%
  \def\bigcap{\mathop{\mathchoice{\raisebox{-0.3em}{\scalebox{1.7}{$\cap$}}}{\scalebox{1.25}{$\cap$}}{\cap}{\cap}}\displaylimits}%
  \def\lesseqqgtr{\mathrel{\lessgtr}}\def\bigoplus{\mathop{\oplus}\displaylimits}%
}
\providecommand{\ssi}{si et seulement si} % abréviation fréquente
\AtBeginDocument{\providecommand{\wideparen}{\overparen}} % arc de cercle

%% FIGFIX-BEGIN v4 — correctifs globaux des figures (docs/onbuch-generation/tools/figfix)
\usepackage{adjustbox}\usepackage{etoolbox}
% toute figure TikZ/pgfplots plus large que la ligne est réduite à la largeur disponible
\BeforeBeginEnvironment{tikzpicture}{\begin{adjustbox}{max width=\linewidth}}
\AfterEndEnvironment{tikzpicture}{\end{adjustbox}}
% légendes pgfplots lisibles par-dessus les courbes
\pgfplotsset{every axis/.append style={legend style={fill=white,fill opacity=0.92,text opacity=1,draw=black!15,inner sep=3pt}}}
% étiquettes d'arêtes (syntaxe edge["..."]) avec fond blanc
\tikzset{every edge quotes/.append style={fill=white,inner sep=1.2pt}}
%% FIGFIX-END
\begin{document}

\pagegarde{Production d’écrits utilitaires : la requête}{Français}{Tle Tech}{Français – classe de Terminale technique}{N11}{15 séances (fiche de progression harmonisée)}{Cette leçon outille l'élève de Terminale technique pour produire des écrits utilitaires, en particulier la requête administrative et professionnelle. À travers l'analyse de textes fonctionnels et iconiques authentiques et la rédaction progressive (entête, introduction, corps, conclusion), l'élève maîtrise la structure externe et interne de la requête et mobilise les actes de langage, le lexique spécialisé et l'emprunt.
\vspace{4pt}
\begin{multicols}{2}\small
\begin{itemize}
\item À la fin de cette leçon, je sais identifier la nature, l'émetteur, le destinataire et la fonction d'un texte fonctionnel (requête, affiche, formulaire)
\item À la fin de cette leçon, je sais distinguer les actes locutoire, illocutoire et perlocutoire dans un énoncé de la vie courante
\item À la fin de cette leçon, je sais décrire la structure externe (entête, objet, formules, signature) et interne (introduction, corps, conclusion) d'une requête
\item À la fin de cette leçon, je sais rédiger l'entête, les formules d'introduction et la conclusion d'une requête conformément aux normes administratives camerounaises
\item À la fin de cette leçon, je sais employer un lexique spécialisé (administratif, juridique, technique) adapté au destinataire et à la situation
\item À la fin de cette leçon, je sais reconnaître et utiliser correctement les emprunts (anglicismes, mots régionaux) dans un écrit utilitaire
\end{itemize}\end{multicols}}

\begin{sommaire}
\begin{multicols}{2}
\begin{enumerate}
\item Les textes fonctionnels et les actes de langage
\item La structure externe et interne de la requête
\item Le lexique spécialisé et l'emprunt
\item Rédaction guidée : entête, introduction, corps et conclusion
\item Exploitation des textes iconiques pour préparer la requête
\item Analyse de situation, plan détaillé et rédaction complète
\item Confrontation, amélioration des productions, compte rendu et remédiation
\item Activité d'intégration
\item Méthodes et exercices résolus
\item Exercices
\item Corrigés des exercices
\item Fiche bilan
\end{enumerate}
\end{multicols}
\end{sommaire}

\begin{objectifs}
\begin{itemize}
\item À la fin de cette leçon, je sais identifier la nature, l'émetteur, le destinataire et la fonction d'un texte fonctionnel (requête, affiche, formulaire)
\item À la fin de cette leçon, je sais distinguer les actes locutoire, illocutoire et perlocutoire dans un énoncé de la vie courante
\item À la fin de cette leçon, je sais décrire la structure externe (entête, objet, formules, signature) et interne (introduction, corps, conclusion) d'une requête
\item À la fin de cette leçon, je sais rédiger l'entête, les formules d'introduction et la conclusion d'une requête conformément aux normes administratives camerounaises
\item À la fin de cette leçon, je sais employer un lexique spécialisé (administratif, juridique, technique) adapté au destinataire et à la situation
\item À la fin de cette leçon, je sais reconnaître et utiliser correctement les emprunts (anglicismes, mots régionaux) dans un écrit utilitaire
\item À la fin de cette leçon, je sais analyser une situation de communication, élaborer un plan détaillé et rédiger une requête complète
\item À la fin de cette leçon, je sais confronter ma production à une grille d'évaluation, l'améliorer et justifier mes choix rédactionnels
\item À la fin de cette leçon, je sais exploiter un texte iconique (organigramme, panneau, formulaire) pour préparer la rédaction d'un écrit utilitaire
\end{itemize}
\end{objectifs}

\begin{prerequis}
\begin{itemize}
\item La lettre administrative et la lettre de motivation (classes de Première et de 3e)
\item Les types et formes de phrases (déclarative, impérative, interrogative)
\item Les connecteurs logiques et la cohérence textuelle (Première)
\item Les registres de langue et le niveau de langue soutenu
\item La correspondance commerciale simple vue en début d'année
\end{itemize}
\end{prerequis}

\newpage

\coursec{Les textes fonctionnels et les actes de langage}

\courssub{Situation d'accroche et problématique}

Imaginez deux scènes de la vie quotidienne au Cameroun. À Yaoundé, au marché Mokolo, M.~Etoundi découvre avec stupeur que sa boutique a été attribuée à un autre commerçant par une erreur administrative. Il décide de saisir le Préfet du Mfoundi. Son fils, pressé, rédige une lettre brouillonne, truffée de fautes, sans formules de politesse administrative, sans référence au décret d'attribution initial, sans pièce jointe. Résultat : le courrier est classé sans suite, M.~Etoundi perd son local. À Douala-Bonabéri, Carine, couturière, reçoit une machine industrielle commandée chez un fournisseur agréé. La machine est défectueuse. Elle rédige une requête précise : en-tête complet, référence de la facture, description technique de la panne, demande de remboursement fondée sur la loi n°2011/012 sur la protection du consommateur, copies des pièces justificatives. Quinze jours plus tard, elle est remboursée intégralement.

\textbf{Problématique :} Qu'est-ce qui distingue une requête efficace d'un courrier inopérant ? Comment rédiger un écrit utilitaire qui obtient réellement gain de cause auprès d'une administration ou d'une entreprise camerounaise ? Cette section pose les fondements théoriques : la nature du texte fonctionnel et la puissance des actes de langage.

\courssub{Le texte fonctionnel : définition et caractéristiques}

\begin{definition}[Texte fonctionnel]
Un \cle{texte fonctionnel} (ou texte utilitaire) est un texte produit dans une situation de communication précise, en réponse à un besoin concret, dans le but d'obtenir un résultat tangible : une information, une autorisation, une réparation, un emploi, une décision administrative. Il se distingue du texte littéraire par sa \cle{visée pratique} et son \cle{efficacité attendue}.
\end{definition}

\begin{aretenir}[Trois piliers du texte fonctionnel]
\begin{enumerate}
    \item \textbf{Visée pratique :} Le texte est un \emph{outil d'action}. On n'écrit pas pour le plaisir d'écrire, mais pour \emph{faire faire} ou \emph{faire savoir}.
    \item \textbf{Destinataire précis (identifié) :} On s'adresse à une personne morale ou physique déterminée (le Préfet, le Directeur Général, le Maire, le Chef de service) qui a le \textbf{pouvoir de décider} ou d'agir.
    \item \textbf{Efficacité attendue :} La réussite se mesure au résultat obtenu (accord, réponse, réparation). L'efficacité impose des normes strictes : clarté, concision, correction linguistique, respect des conventions de forme (structure externe) et de fond (argumentation).
\end{enumerate}
\end{aretenir}

\begin{exemplebox}[Texte littéraire vs Texte fonctionnel]
\begin{center}
\begin{tabularx}{\linewidth}{|l|X|X|}\hline
\textbf{Critère} & \textbf{Texte littéraire (ex. roman)} & \textbf{Texte fonctionnel (ex. requête)} \\ \hline
\rowcolor{popblueL} \textbf{Finalité} & Esthétique, émotion, exploration & Pragmatique, action, résultat \\ \hline
\textbf{Destinataire} & Indéterminé (lecteur universel) & Déterminé (décideur identifié) \\ \hline
\textbf{Structure} & Libre, créative & Codifiée, normative (entête, corps, formules) \\ \hline
\textbf{Langage} & Connoté, figuré, subjectif & Dénoté, précis, objectif, standardisé \\ \hline
\textbf{Critère de réussite} & Plaisir du lecteur, valeur artistique & Obtention de la satisfaction (réponse favorable) \\ \hline
\end{tabularx}
\end{center}
\end{exemplebox}

\courssub{Lecture guidée : Texte fonctionnel N°1 -- Une requête authentique}

Analysons l'extrait suivant, inspiré d'une requête réelle adressée au Préfet du Mfoundi.

\begin{exemplebox}[Extrait de requête -- Demission de régularisation d'attribution de boutique]
\textbf{République du Cameroun}\\
\textbf{Paix -- Travail -- Patrie}\\[0.5em]
Yaoundé, le 15 mars 2025\\[1em]
\textbf{À}\\
\textbf{Monsieur le Préfet du Département du Mfoundi}\\
\textbf{Service des Affaires Foncières et Domaniales}\\
\textbf{Yaoundé}\\[1em]
\textbf{Objet :} Requête en régularisation d'attribution de boutique n°45, Marché Mokolo.\\[1em]
\textbf{Monsieur le Préfet,}\\[0.5em]
J'ai l'honneur de solliciter votre bienveillance pour la régularisation de ma situation locative concernant la boutique n°45 sise au Marché Mokolo, quartier Mokolo II, arrondissement de Yaoundé II.\\[0.5em]
Par arrêté municipal n°012/YII/2023 du 10 janvier 2023, j'ai été régulièrement attributaire de ladite boutique. Contre toute attente, un avis d'attribution en date du 05 mars 2025, portant le n°045/YII/2025, a été affiché au nom de M.~Kamga Jean, m'attribuant le même local.\\[0.5em]
Cette erreur matérielle porte préjudice à mon activité commerciale, unique source de revenus pour ma famille. Je joins en annexe : la copie de l'arrêté d'attribution initial, la copie de ma CNI, les quittances de loyer à jour.\\[0.5em]
En espérant une suite favorable à ma demande, je vous prie d'agréer, Monsieur le Préfet, l'expression de ma haute considération.\\[0.5em]
\textbf{Le requérant}
\end{exemplebox}


\coursec{La structure externe et interne de la requête}

\courssub{Transition : des actes de langage à l'écrit administratif codifié}

Dans la section précédente, nous avons analysé les \textbf{actes de langage} (locutoire, illocutoire, perlocutoire) qui fondent toute communication. Nous avons vu qu'un énoncé comme « Je vous demande de bien vouloir m'accorder... » n'est pas une simple description : c'est un \emph{acte illocutoire} de type \textbf{directif} (une requête, une prière, un ordre) visant à obtenir une action de la part du destinataire.

La \textbf{requête administrative} est l'archétype de cet acte directif mis par écrit dans un cadre institutionnel. Elle ne tolère ni l'improvisation ni la familiarité : elle obéit à un \textbf{rituel graphique strict} (structure externe) et à une \textbf{rhétorique codifiée} (structure interne) qui garantissent sa recevabilité et son efficacité. Maîtriser ce genre, c'est passer de la compétence langagière (savoir demander) à la compétence \emph{scripturale professionnelle} (savoir rédiger une demande recevable).

\begin{definition}[La requête : définition fonctionnelle]
La \cle{requête} est un \textbf{texte fonctionnel (écrit utilitaire)} par lequel un \textbf{requérant} (administré, employé, justiciable) expose \textbf{respectueusement} une \textbf{demande précise} (droit, service, réparation, faveur) à une \textbf{autorité compétente} (administrative, hiérarchique ou judiciaire) dont il relève ou qui a le pouvoir de statuer.
\end{definition}

\begin{savaistu}[Le timbre fiscal au Cameroun]
Au Cameroun, la recevabilité de nombreuses requêtes administratives (demande d'acte de naissance, de certificat de nationalité, de licence, recours gracieux auprès du Préfet...) est conditionnée par l'apposition d'un \textbf{timbre fiscal} (souvent 1000 FCFA ou plus selon la nature de l'acte). Une requête non timbrée peut être déclarée \emph{irrecevable} sans examen du fond. Vérifiez toujours la réglementation en vigueur auprès du secrétariat de l'administration concernée.
\end{savaistu}

\courssub{La structure externe : l'architecture visible du document}

La structure externe est la \textbf{mise en page normalisée}. Elle permet au secrétariat de trier, classer et router le courrier instantanément. L'omission d'un seul élément peut entraîner le rejet du dossier.

\begin{methode}[Vérifier la structure externe : les 8 éléments obligatoires]
Avant de signer, cochez cette liste de contrôle (ordre chronologique de lecture) :
\begin{enumerate}
    \item \textbf{Lieu et date} : en haut à droite (ex. : \emph{Yaoundé, le 15 octobre 2025}).
    \item \textbf{Coordonnées du requérant} : en haut à gauche (Nom, Prénoms, Adresse complète, Téléphone, Email, Matricule si agent).
    \item \textbf{Mention du destinataire} : en haut à gauche, sous le requérant. \textbf{Titre exact et complet} de l'autorité (ex. : \emph{Monsieur le Préfet du Mfoundi}, \emph{Monsieur le Directeur Général de la SNH}, \emph{Monsieur le Procureur de la République près le Tribunal de Grande Instance du Mfoundi}).
    \item \textbf{Objet} : centré, en gras, concis et précis. Commence par un nom : \emph{Demande de...}, \emph{Réclamation pour...}, \emph{Recours contre...}.
    \item \textbf{Références} (si applicable) : numéro de dossier, décision attaquée, lettre précédente (ex. : \emph{Réf. : Votre lettre n° 045/MINFOPRA/SG/DGRH du 10/09/2025}).
    \item \textbf{Corps de la lettre} : introduction, développement, conclusion (voir structure interne).
    \item \textbf{Formule de politesse finale + Signature} : formule adaptée au grade du destinataire, suivie de la signature manuscrite (précédée de \emph{Fait à..., le...} si lieu/date non mis en en-tête).
    \item \textbf{Pièces jointes} : liste numérotée en bas à gauche (ex. : \emph{Pièces jointes : 1. Copie CNI ; 2. Attestation de résidence ; 3. Quittance de timbre fiscal}).
\end{enumerate}
\end{methode}

\begin{attention}[Piège fréquent : l'objet « fourre-tout »]
\textbf{Erreur :} Écrire \emph{Objet : Requête} ou \emph{Objet : Demande}.
\textbf{Conséquence :} Le secrétariat ne sait pas de quoi il s'agit ; le dossier traîne.
\textbf{Correction :} L'objet doit être une \textbf{phrase nominale informative complète}.
\begin{itemize}
    \item \emph{Objet : Demande de réattribution de la boutique n°12 du Marché Central de Yaoundé.}
    \item \emph{Objet : Recours gracieux contre la décision de radiation des effectifs n° 012/MINESEC/SG/DRH du 05/09/2025.}
    \item \emph{Objet : Demande de stage de fin d'études (Licence en Génie Civil) du 02 janvier au 28 février 2026.}
\end{itemize}
\end{attention}

\begin{popfigure}
\begin{tikzpicture}[
    box/.style={rectangle, draw=popdark, thick, rounded corners=3pt, fill=popcream, inner sep=4pt, font=\small},
    arrow/.style={->, >=Stealth, thick, popblue},
    label/.style={font=\footnotesize\bfseries, text=poppurple, anchor=west},
    title/.style={font=\small\bfseries, text=popdark}
]
% Page frame
\draw[popline, thick] (0,0) rectangle (14,19);

% 1. Lieu/Date (Top Right)
\node[box, anchor=north east] (lieu) at (13.5, 18.5) {Yaoundé, le 15 octobre 2025};
\draw[arrow] (10, 18.8) -- (11.5, 18.8) node[label, midway, above] {1. Lieu \& Date};

% 2. Requérant (Top Left)
\node[box, anchor=north west, text width=5cm, align=left] (requerant) at (0.5, 18.5) {
    \textbf{NOM Prénom(s)}\\
    BP 12345 Yaoundé\\
    Tel: 6XX XX XX XX\\
    Email: nom@email.cm\\
    Matricule: 2025/ABC/123
};
\draw[arrow] (5.5, 18.8) -- (7, 18.8) node[label, midway, above] {2. Coordonnées Requérant};

% 3. Destinataire (Below Requérant)
\node[box, anchor=north west, text width=5cm, align=left] (dest) at (0.5, 15.8) {
    \textbf{À l'attention de}\\
    \textbf{Monsieur le Préfet}\\
    \textbf{du Mfoundi}\\
    \textbf{Yaoundé}
};
\draw[arrow] (5.5, 16.5) -- (7, 16.5) node[label, midway, above] {3. Destinataire (Titre exact)};

% 4. Objet (Centered)
\node[box, anchor=north, fill=popblueL, font=\bfseries\small] (objet) at (7, 14.2) {Objet : Demande de réattribution de la boutique n°12 du Marché Central};
\draw[arrow] (7, 13.5) -- (7, 12.5) node[label, below] {4. Objet (Précis \& Complet)};

% 5. Références (Optional, below Objet)
\node[box, anchor=north, font=\small] (ref) at (7, 12.0) {Réf. : Arrêté Préfectoral n° 045/MP/MFD du 10/09/2025};
\draw[arrow] (9.5, 11.8) -- (11, 11.8) node[label, midway, right] {5. Références (si existant)};

% 6. Corps de la lettre (Large box)
\node[box, anchor=north, text width=12.5cm, align=justify, minimum height=7cm] (corps) at (7, 10.5) {
    \textbf{Monsieur le Préfet,}\\[4pt]
    \textit{\textbf{Introduction (Identité + Motif) :}}\\
    J'ai l'honneur de solliciter votre bienveillance pour...\\[6pt]
    \textit{\textbf{Corps (Faits + Arguments + Justificatifs) :}}\\
    En effet, suite à l'arrêté cité en référence, ma boutique a été attribuée à un tiers alors que je suis à jour de mes redevances (pièce 3). Cette décision porte atteinte à mon droit de propriété commerciale...\\[6pt]
    \textit{\textbf{Conclusion (Demande claire + Remerciements) :}}\\
    En foi de quoi, je vous prie, Monsieur le Préfet, de bien vouloir annuler ladite attribution et me rétablir dans mes droits.\\
    Je vous prie d'agréer, Monsieur le Préfet, l'expression de ma haute considération.
};
\draw[arrow] (0.5, 7.5) -- (-1.5, 7.5) node[label, left, align=right] {6. Corps de la lettre\\(Intro, Développement,\\Conclusion)};

% 7. Formule + Signature (Bottom Right)
\node[box, anchor=south east, text width=6cm, align=left] (sig) at (13.5, 1.5) {
    Fait à Yaoundé, le 15 octobre 2025\\
    \vspace{1.5cm}
    \textbf{Signature}\\
    \textit{(NOM Prénom)}
};
\draw[arrow] (13.5, 4.5) -- (15, 4.5) node[label, right, align=left] {7. Formule finale\\(Politesse + Signature)};

% 8. Pièces jointes (Bottom Left)
\node[box, anchor=south west, text width=5cm, align=left] (pj) at (0.5, 1.5) {
    \textbf{Pièces jointes :}\\
    1. Copie CNI\\
    2. Copie Arrêté attaqué\\
    3. Quittances de loyer\\
    4. Timbre fiscal 1000 FCFA
};
\draw[arrow] (0.5, 4.5) -- (-1.5, 4.5) node[label, left, align=right] {8. Pièces jointes\\(Liste numérotée)};
\end{tikzpicture}
\legende{Schéma annoté d'une page de requête type : repérage visuel des 8 zones obligatoires de la structure externe.}
\end{popfigure}

\courssub{La structure interne : la rhétorique de la persuasion administrative}

Si l'externe assure la \emph{recevabilité formelle}, l'interne assure la \emph{recevabilité du fond}. Elle suit le schéma classique de l'argumentation \emph{exordium - narratio - confirmatio - conclusio}, adapté à l'administration.

\begin{definition}[Les trois séquences de la structure interne]
\begin{enumerate}
    \item \textbf{L'Introduction (Exordium) :} \emph{Identité du requérant} (qualité, matricule, domicile) + \emph{Motif immédiat} de l'écriture (référence à un fait, une décision, une situation). \textbf{Max 3 à 4 lignes.}
    \item \textbf{Le Corps / Développement (Narratio + Confirmatio) :}
        \begin{itemize}
            \item \textbf{Exposé des faits :} Chronologique, factuel, sans pathos.
            \item \textbf{Arguments juridiques / réglementaires :} Citation des textes (loi, décret, arrêté, règlement intérieur, contrat).
            \item \textbf{Justificatifs :} Renvoi explicite aux pièces jointes (\emph{« comme en attestent les pièces 2 et 3 ci-jointes »}).
        \end{itemize}
    \item \textbf{La Conclusion (Conclusio) :}
        \begin{itemize}
            \item \textbf{Formulation explicite de la demande (la « pétition ») :} Verbe à l'infinitif ou subjonctif (\emph{« annuler », « accorder », « réintégrer », « délivrer »}).
            \item \textbf{Formule de politesse finale adaptée au grade.}
        \end{itemize}
\end{enumerate}
\end{definition}

\begin{exemplebox}[Découpage d'une introduction et d'une conclusion modèles]
\textbf{Introduction (Modèle) :}
\begin{quote}
\textit{Monsieur le Préfet,}\par
\textit{Je soussigné(e) \textbf{MBARGA Jean}, commerçant, titulaire de la carte de contribuable n° 2025/045/MFD, demeurant au quartier Mvog-Ada, bp 876 Yaoundé, ai l'honneur de porter à votre connaissance la situation relative à la boutique n°12 du Marché Central, m'attribuée par l'arrêté n° 012/MP/MFD du 05/01/2024.}
\end{quote}
\textit{Analyse :} Identité complète (nom, qualité, identifiant fiscal, adresse) + Référence précise du droit acquis (arrêté d'attribution initial). \textbf{3 lignes.}

\textbf{Conclusion (Modèle) :}
\begin{quote}
\textit{En conséquence, je vous saurais gré de bien vouloir \textbf{annuler la décision d'attribution litigieuse} et de \textbf{me rétablir dans mes droits} de titulaire de la boutique n°12.}\par
\textit{En vous remerciant par avance de l'intérêt que vous voudrez bien réserver à ma doléance, je vous prie d'agréer, Monsieur le Préfet, l'assurance de ma très haute considération.}
\end{quote}
\textit{Analyse :} Deux verbes d'action précis (\emph{annuler, rétablir}) + Formule de politesse « haute considération » réservée aux Préfets, Ministres, Directeurs Généraux, Procureurs.
\end{exemplebox}

\begin{aretenir}[Formules de politesse finales : le grade dicte la formule]
\begin{center}
\begin{tabularx}{\linewidth}{|l|X|}\hline
\textbf{Destinataire} & \textbf{Formule de politesse (usage camerounais)} \\ \hline
Préfet, Ministre, DG, Procureur, Recteur & \emph{Je vous prie d'agréer, [Titre], l'assurance de ma très haute considération.} \\ \hline
Sous-Préfet, Chef de Service, Proviseur, Directeur d'École & \emph{Je vous prie d'agréer, [Titre], l'expression de ma haute considération.} \\ \hline
Chef de Bureau, Supérieur hiérarchique direct & \emph{Je vous prie d'agréer, [Titre], l'expression de ma considération distinguée.} \\ \hline
Collègue, Administration sans grade élevé (secrétariat) & \emph{Je vous prie de croire, [Titre], à l'assurance de mes salutations distinguées.} \\ \hline
\end{tabularx}
\end{center}
\textbf{Règle d'or :} On ne « salue » pas une autorité, on lui « prie d'agréer l'expression de sa considération ». Le tutoiement et les formules familières (\emph{« Cordialement », « Bien à vous »}) sont \textbf{interdits}.
\end{aretenir}

\courssub{Typologie des requêtes : ne pas confondre les genres}

Le programme exige de distinguer la requête administrative des écrits voisins. Le tableau ci-dessous clarifie les frontières pour l'examen (sujet de dissertation ou cas pratique).

\begin{center}
\begin{tabularx}{\linewidth}{|l|X|X|X|X|}\hline
\textbf{Critère} & \textbf{Requête Administrative} & \textbf{Demande d'Emploi (Candidature)} & \textbf{Réclamation (Réclamation commerciale/technique)} & \textbf{Requête Judiciaire (Requête introductive d'instance)} \\ \hline
\textbf{Destinataire} & Autorité administrative (Préfet, Ministre, Maire, DG) & Employeur (DRH, Patron, DG) & Fournisseur, Vendeur, Prestataire de service & Juge (Tribunal de Grande Instance, Conseil de Prud'hommes) \\ \hline
\textbf{Fondement} & Droit administratif, Règlement intérieur, Discrétionnaire & Contrat de travail (offre d'emploi), Besoin de l'entreprise & Contrat de vente/prestation, Garantie, Droit de la consommation & Code de procédure civile, Droit substantiel \\ \hline
\textbf{Objet} & Droit administratif, Service public, Avantage, Grâce, Rétablissement & Poste de travail, Stage, Formation & Remplacement, Remboursement, Réparation, Indemnisation & Saisine du juge pour faire valoir un droit (assignation) \\ \hline
\textbf{Forme} & Libre mais codifiée (usage) & Très codifiée (CV + Lettre de motivation) & Libre (LRAR recommandé) & \textbf{Très stricte} (avocat obligatoire souvent, mentions obligatoires art. 58 CPC) \\ \hline
\textbf{Issues} & Décision administrative (Arrêté, Lettre) & Embauche / Refus / Stage & Accord commercial / Médiation / Tribunal & Jugement / Ordonnance \\ \hline
\end{tabularx}
\end{center}

\begin{attention}[Distinction Bac : Requête Administrative vs Requête Judiciaire]
\begin{itemize}
    \item La \textbf{requête administrative} (ex. : recours gracieux au Préfet) est un \textbf{acte unilatéral} adressé à l'auteur de la décision ou à son supérieur. Pas d'avocat obligatoire. Délai : souvent 2 mois (recours gracieux) ou 60 jours (recours hiérarchique).
    \item La \textbf{requête judiciaire} (ex. : saisine du Tribunal Administratif ou TGI) est un \textbf{acte de procédure}. Elle \textbf{doit} contenir les mentions de l'article 58 du Code de Procédure Civile (nom du demandeur, domicile, objet, moyens, conclusions, signature d'avocat). Elle interrompt la prescription.
    \item \textbf{Piège :} Ne pas rédiger une requête administrative avec les mentions d'une assignation (ex. « Vu l'article 58 CPC... ») et inversement.
\end{itemize}
\end{attention}

\courssub{Analyse collective d'une requête modèle : repérage et étiquetage}

Appliquons la méthode à un cas concret complet.

\begin{exoresolu}[Analyse structurelle d'une requête de recours gracieux]
\textbf{Énoncé :} Voici le texte intégral d'une requête. Identifiez et étiquetez chaque composante de la structure externe et interne.

\begin{quote}
\textbf{Yaoundé, le 20 novembre 2025}\hfill \textit{(1)}

\textbf{MBALLA Pierre}\\
BP 54321 Yaoundé\\
Tel: 690 12 34 56\\
Matricule: 2018/0045/MINEDUB \hfill \textit{(2)}

\textbf{À}\\
\textbf{Monsieur le Ministre}\\
\textbf{des Enseignements Secondaires}\\
\textbf{Yaoundé} \hfill \textit{(3)}

\hfill \textbf{Objet : Recours gracieux contre la décision de mutation d'office n° 045/MINESEC/SG/DRH du 15/10/2025.} \hfill \textit{(4)}

\textbf{Réf. : Lettre de notification n° 112/DRH/SDP du 20/10/2025.} \hfill \textit{(5)}

\textbf{Monsieur le Ministre,}\par
\textit{Par la présente, je soussigné MBALLA Pierre, professeur de Mathématiques au Lycée de Bafia (Matricule 2018/0045/MINEDUB), viens respectueusement former un recours gracieux contre la décision de mutation d'office susvisée m'affectant au Lycée de Bertoua.} \hfill \textit{(6) Introduction}

\textit{En effet, cette décision, bien que relevant de votre pouvoir discrétionnaire, méconnaît les dispositions de l'article 12 du Statut Général de la Fonction Publique qui dispose que « tout fonctionnaire a droit à une affectation tenant compte de sa situation familiale ». Je suis marié, père de trois enfants scolarisés à Yaoundé, et mon épouse est fonctionnaire au Ministère des Finances à Yaoundé (Pièce 1 : Acte de mariage ; Pièce 2 : Attestation de scolarité ; Pièce 3 : Attestation de service de l'épouse).} \hfill \textit{(7) Corps : Faits + Argument juridique + Pièces}

\textit{Par ailleurs, mon état de santé nécessite un suivi médical régulier au Centre Hospitalier Universitaire de Yaoundé (Pièce 4 : Certificat médical). Une mutation à Bertoua (400 km) compromettrait gravement ma prise en charge thérapeutique et la stabilité de ma famille.} \hfill \textit{(8) Corps : Argument humanitaire/médical + Pièce}

\textit{Au regard de ce qui précède, je vous saurais gré de bien vouloir \textbf{annuler ladite décision de mutation} et de \textbf{me maintenir à mon poste actuel au Lycée de Bafia} ou, à défaut, de m'affecter dans un établissement de la ville de Yaoundé.} \hfill \textit{(9) Conclusion : Pétition claire (Annuler + Maintenir/Affecter)}

\textit{En espérant que ma demande recevra une suite favorable, je vous prie d'agréer, Monsieur le Ministre, l'assurance de ma très haute considération.} \hfill \textit{(10) Formule de politesse}

\textbf{Fait à Yaoundé, le 20 novembre 2025}\\
\textbf{Signature} \hfill \textit{(11) Lieu/Date + Signature}

\textbf{Pièces jointes :}\\
1. Copie certifiée conforme de l'acte de mariage\\
2. Attestations de scolarité des enfants (3)\\
3. Attestation de service de l'épouse\\
4. Certificat médical du CHU Yaoundé (daté de moins de 3 mois)\\
5. Copie de la décision de mutation attaquée\\
6. Quittance de timbre fiscal 1000 FCFA \hfill \textit{(12) Pièces jointes}
\end{quote}

\textbf{Solution détaillée (Étiquetage) :}

\begin{center}
\begin{tabularx}{\linewidth}{|c|l|X|}\hline
\textbf{Zone} & \textbf{Structure} & \textbf{Contenu repéré / Commentaire} \\ \hline
(1) & \textbf{Externe} & Lieu et Date (Haut droit) \\ \hline
(2) & \textbf{Externe} & Coordonnées Requérant (Haut gauche : Identité, Adresse, Tel, Matricule) \\ \hline
(3) & \textbf{Externe} & Destinataire (Titre exact : Ministre des Enseignements Secondaires) \\ \hline
(4) & \textbf{Externe} & \textbf{Objet Précis} (Nature : Recours gracieux + Référence décision attaquée + Date + Numéro) \\ \hline
(5) & \textbf{Externe} & Références (Lettre de notification reçue) \\ \hline
(6) & \textbf{Interne : Introduction} & \textbf{Identité complète} (Nom, Grade, Poste actuel, Matricule) + \textbf{Motif} (Recours contre décision précise). \emph{Respect de la limite ~3 lignes.} \\ \hline
(7) & \textbf{Interne : Corps (Arg. Juridique)} & \textbf{Fait :} Mutation Bertoua. \textbf{Droit :} Art 12 Statut Fonction Publique (Situation familiale). \textbf{Preuves :} Pièces 1, 2, 3. \\ \hline
(8) & \textbf{Interne : Corps (Arg. Médical)} & \textbf{Fait :} Suivi CHU Yaoundé. \textbf{Conséquence :} Rupture soins. \textbf{Preuve :} Pièce 4. \\ \hline
(9) & \textbf{Interne : Conclusion} & \textbf{Pétition double :} 1. Annulation. 2. Maintien ou Affectation Yaoundé. \emph{Verbes d'action précis.} \\ \hline
(10) & \textbf{Interne : Conclusion} & Formule de politesse \textbf{adaptée au Ministre} (\emph{très haute considération}). \\ \hline
(11) & \textbf{Externe} & Lieu/Date (bas droit) + Signature manuscrite. \\ \hline
(12) & \textbf{Externe} & Pièces jointes \textbf{numérotées}, cohérentes avec les renvois du corps (1 à 4) + pièces de procédure (5) + timbre (6). \\ \hline
\end{tabularx}
\end{center}
\end{exoresolu}

\courssub{Les qualités d'une requête recevable et efficace}

Au-delà de la structure, l'examinateur (ou l'administration) évalue la \textbf{qualité rédactionnelle}. Le programme en liste cinq : clarté, concision, exactitude, courtoisie, respect de la hiérarchie.

\begin{propriete}[Les 5 qualités canoniques de la requête (Grille d'évaluation Bac)]
\begin{enumerate}
    \item \textbf{Clarté :} Une phrase = une idée. Vocabulaire précis (pas de « truc », « chose », « problème »). Plan visible (paragraphes distincts Intro/Corps/Conclusion). Connecteurs logiques (\emph{donc, par conséquent, en outre, toutefois}).
    \item \textbf{Concision :} \textbf{15 à 25 lignes au total} (hors en-tête). Introduction $\le$ 3 lignes. Conclusion $\le$ 3 lignes. Pas de répétitions, pas de digressions narratives (« Il était une fois... »).
    \item \textbf{Exactitude :} Dates, numéros, noms propres, articles de loi, numéros de pièces \textbf{vérifiés}. Une erreur factuelle (mauvais n° d'arrêté) discrédite tout le dossier.
    \item \textbf{Courtoisie :} Registre soutenu. Verbes déontiques au conditionnel (\emph{je vous saurais gré, je vous prie de bien vouloir}). \textbf{Interdit :} « Je veux », « Il faut que », « Vous devez », menaces, chantage affectif excessif.
    \item \textbf{Respect de la hiérarchie :} Titre exact du destinataire. Formule de politesse calibrée. Pas de familiarité (pas de prénom du ministre). Reconnaissance implicite de l'autorité du décideur (\emph{« bien vouloir examiner », « daigner accorder »}).
\end{enumerate}
\end{propriete}

\begin{exemplebox}[Chiffres clés : la « métrique » de la requête type]
\begin{center}
\begin{tabular}{|l|c|c|}\hline
\textbf{Composante} & \textbf{Lignes moyennes} & \textbf{Pourcentage du corps} \\ \hline
En-tête externe (Lieu, Date, Coord, Dest, Objet, Ref) & 6 à 10 lignes & — \\ \hline
Introduction & 2 à 3 lignes & $\approx$ 10\% \\ \hline
Corps (Faits + 2 à 3 arguments + renvois pièces) & 10 à 15 lignes & $\approx$ 70\% \\ \hline
Conclusion (Pétition + Formule politesse) & 3 à 4 lignes & $\approx$ 20\% \\ \hline
Signature + Pièces jointes & 3 à 5 lignes & — \\ \hline
\textbf{TOTAL CORPS DE TEXTE} & \textbf{15 à 22 lignes} & \textbf{100\%} \\ \hline
\end{tabular}
\end{center}
\textit{Conseil :} Si votre brouillon dépasse 30 lignes, coupez dans le narratif (faits), gardez l'argumentaire juridique.
\end{exemplebox}

\begin{popfigure}
\begin{tikzpicture}[
    grow'=right, level distance=3.5cm,
    level 1/.style={sibling distance=3.5cm},
    level 2/.style={sibling distance=3cm},
    every node/.style={draw=popdark, thick, rounded corners=4pt, fill=popcream, font=\small, inner sep=4pt, align=center},
    edge from parent/.style={draw=popblue, thick, ->, >=Stealth},
    root/.style={fill=popblue, text=white, font=\small\bfseries, rounded corners=4pt},
    childnode/.style={fill=popblueL, font=\small\bfseries},
    leaf/.style={fill=popgreenL, font=\footnotesize, align=left, text width=3.5cm}
]
\node[root] {Écrits Utilitaires}
    child { node[childnode] {Requête Administrative}
        child { node[leaf] {Recours gracieux (à l'auteur de la décision)} }
        child { node[leaf] {Recours hiérarchique (au supérieur)} }
        child { node[leaf] {Demande de service / droit (bourse, logement, acte)} }
        child { node[leaf] {Demande de grâce / amnistie} }
    }
    child { node[childnode] {Requête Professionnelle (Interne)}
        child { node[leaf] {Demande de congé (annuel, maladie, maternité)} }
        child { node[leaf] {Demande de mutation / détachement} }
        child { node[leaf] {Demande de formation / stage / avancement} }
        child { node[leaf] {Réclamation interne (salaire, carrière)} }
    }
    child { node[childnode] {Requête Judiciaire (Contentieux)}
        child { node[leaf] {Requête introductive d'instance (TGI, TA)} }
        child { node[leaf] {Référé (urgence)} }
        child { node[leaf] {Appel / Pourvoi en cassation} }
        child { node[leaf] {Requête en référé-liberté / suspension} }
    };
\end{tikzpicture}
\legende{Arbre de classification des requêtes : trois grandes familles selon le destinataire et le fondement juridique. La leçon porte sur la branche « Administrative » (et professionnelle interne).}
\end{popfigure}

\courssub{Synthèse et passage à la rédaction guidée}

Nous avons maintenant une \textbf{check-list complète} pour construire une requête :
\begin{enumerate}
    \item \textbf{Externe :} 8 blocs visuels (Lieu/Date, Requérant, Destinataire, Objet, Réf, Corps, Formule/Sign, PJ).
    \item \textbf{Interne :} 3 séquences logiques (Intro : Qui/Quoi $\rightarrow$ Corps : Faits/Droits/Preuves $\rightarrow$ Conclusion : Demande/Politesse).
    \item \textbf{Typage :} Administrative (autorité publique) vs Professionnelle (employeur) vs Judiciaire (juge).
    \item \textbf{Qualité :} Clarté, Concision (15-25 lignes), Exactitude, Courtoisie, Hiérarchie.
\end{enumerate}

La section suivante (\textbf{Section 3}) mobilisera le \textbf{lexique spécialisé} (vocabulaire administratif : \emph{vu, attendu, considérant, arrête, décide}) et les mécanismes de l'\textbf{emprunt lexical} pour outiller la rédaction précise de chaque zone, avant de passer à la rédaction guidée pas à pas (Section 4).

\begin{exercice}[Entraînement : Diagnostic structurel]
Voici l'en-tête d'une requête rédigée par un élève. Relevez \textbf{4 erreurs majeures} de structure externe et proposez la correction.

\begin{quote}
\textbf{Douala, le 12/03/25}\\
\textbf{Etudiant en Terminale T}\\
\textbf{Lycée Technique de Bonabéri}\\
\textbf{Au Proviseur}\\
\textbf{Objet : Requête}\\
\textbf{Monsieur le Proviseur, je veux changer de série...}
\end{quote}
\end{exercice}

\begin{corrige}[Exercice n°1 : Diagnostic structurel]
\textbf{Erreurs repérées :}
\begin{enumerate}
    \item \textbf{Date incomplète :} Format « 12/03/25 » (chiffres) au lieu de « Douala, le 12 mars 2025 » (lettres pour le mois, 4 chiffres pour l'année).
    \item \textbf{Coordonnées requérant insuffisantes :} Pas de nom/prénom, pas d'adresse domicile, pas de téléphone, pas de matricule. « Étudiant en Terminale T » est une qualité, pas une identité.
    \item \textbf{Destinataire imprécis :} « Au Proviseur » $\rightarrow$ Doit être : \textbf{À l'attention de Monsieur le Proviseur du Lycée Technique de Bonabéri}.
    \item \textbf{Objet non informatif :} « Objet : Requête » $\rightarrow$ Doit être : \textbf{Objet : Demande de changement de série (Série T vers Série F4) pour l'année scolaire 2025-2026}.
    \item \textbf{(Bonus) Familiarité dans l'amorce :} « je veux changer » $\rightarrow$ Remplacer par : « J'ai l'honneur de solliciter votre bienveillance pour un changement de série... ».
\end{enumerate}
\end{corrige}

\coursec{Le lexique spécialisé et l'emprunt}

Dans la section précédente, nous avons appris à construire la \cle{structure externe} (en-tête, objet, formules, signature) et la \cle{structure interne} (introduction, corps, conclusion) de la requête. Mais une requête bien structurée peut être gâchée par des mots mal choisis : un vocabulaire trop familier, un terme emprunté à l'anglais sans adaptation, ou pire, une expression de camfranglais. Cette section nous apprend donc à \textbf{choisir les bons mots} : ceux du \cle{lexique spécialisé} du domaine concerné, et à gérer correctement les \cle{emprunts} qui envahissent notre français quotidien.

\courssub{Le lexique spécialisé : définition et rôle}

Commençons par une observation simple. Quand un mécanicien parle de « culasse », un comptable de « bilan » et un greffier d'« assignation », chacun utilise des mots que les autres comprennent mal. Ces mots ne sont pas du vocabulaire courant : ils appartiennent à un \textbf{domaine d'activité précis} et garantissent la précision du message.

\begin{definition}[Lexique spécialisé]
Le \cle{lexique spécialisé} est l'ensemble des termes propres à un domaine d'activité (administration, droit, commerce, industrie, informatique, médecine, etc.). Il constitue le vocabulaire technique de ce domaine et garantit la \textbf{précision} et la \textbf{crédibilité} des écrits utilitaires. Dans une requête, l'emploi juste du lexique spécialisé montre au destinataire que le requérant maîtrise son sujet.
\end{definition}

Pourquoi est-ce si important dans une requête ? Parce que le destinataire (un préfet, un directeur, un chef de service) juge la demande dès les premières lignes. Une requête qui dit « je demande un papier pour prouver que j'ai travaillé chez vous » fera moins d'effet qu'une requête qui demande « un \textbf{certificat de travail} » : le terme exact existe, il faut l'utiliser.

\begin{exemplebox}[Le même besoin, deux formulations]
\begin{itemize}
\item \emph{Formulation vague :} « Je vous écris pour que vous me donniez quelque chose qui montre que j'ai payé. »
\item \emph{Formulation spécialisée :} « J'ai l'honneur de solliciter la délivrance d'une \textbf{quittance de loyer} relative au mois de mars 2026. »
\end{itemize}
La seconde version emploie le terme juridique et commercial exact : le destinataire sait immédiatement quel document produire.
\end{exemplebox}

\courssub{Le champ lexical de la requête administrative}

La requête administrative possède son propre vocabulaire codifié. Ces expressions, parfois anciennes, sont \textbf{attendues} par l'administration : les omettre ou les déformer est une faute de registre.

\begin{center}
\begin{tabularx}{\linewidth}{|l|X|}
\hline
\rowcolor{popblueL} \textbf{Expression consacrée} & \textbf{Emploi et sens} \\
\hline
« Je soussigné(e)... » & Formule d'identification placée en tête du corps de la requête ; signifie « moi qui signe au bas de ce document ». \\
\hline
« domicilié(e) à... » & Indique l'adresse officielle du requérant (et non « habitant quelque part »). \\
\hline
« J'ai l'honneur de... » & Formule de politesse qui introduit la demande elle-même. \\
\hline
« pièces jointes » & Documents annexés à la requête (photocopies, justificatifs), énumérés en bas de page. \\
\hline
« suite favorable » & Réponse positive que le requérant espère : « dans l'espoir d'une suite favorable... ». \\
\hline
« à toutes fins utiles » & Formule finale signifiant « pour tous les usages où ce document pourra servir ». \\
\hline
\end{tabularx}
\end{center}

\begin{attention}[Piège : déformer les formules consacrées]
On lit souvent « Je sous-signé » (faute : c'est \textbf{soussigné}, en un seul mot), ou « requette » (faute d'orthographe : \textbf{requête}, sans double \emph{t}). Ces formules sont figées : elles s'écrivent toujours de la même manière. Apprends-les par cœur comme des formules mathématiques.
\end{attention}

\courssub{Le lexique spécialisé selon les filières techniques}

L'élève de Terminale technique rédigera demain des requêtes \textbf{professionnelles} : demande de stage, demande d'emploi, réclamation auprès d'un fournisseur, demande de congé. Chaque filière possède son vocabulaire, qu'il faut mobiliser avec exactitude.

\begin{center}
\begin{tabularx}{\linewidth}{|l|X|}
\hline
\rowcolor{popgreenL} \textbf{Filière} & \textbf{Termes spécialisés utiles dans une requête} \\
\hline
Comptabilité / Gestion & bilan, facture, avoir, échéance, grand livre, rapprochement bancaire, amortissement \\
\hline
Commerce / Vente & bon de commande, bon de livraison, devis, stock, marge, réclamation, avoir commercial \\
\hline
Industrie / Maintenance & cahier des charges, maintenance préventive, pièce de rechange, chaîne de production, contrôle qualité \\
\hline
Informatique & cahier des charges fonctionnel, maintenance, licence logicielle, sauvegarde, réseau \\
\hline
\end{tabularx}
\end{center}

\begin{exemplebox}[Requête professionnelle dans une filière technique]
Un élève de la filière Gestion demandant un stage écrira : « J'ai l'honneur de solliciter un \textbf{stage académique} au sein de votre \textbf{service comptable}, afin de mettre en pratique les notions de \textbf{tenue des livres} et de \textbf{rapprochement bancaire} acquises en formation. » Chaque terme en gras appartient au lexique spécialisé de sa filière : la demande gagne immédiatement en sérieux.
\end{exemplebox}

\courssub{L'emprunt : définition et procédés}

Notre français quotidien est plein de mots venus d'ailleurs : « planning », « weekend », « challenge »... D'où viennent-ils et comment sont-ils entrés dans la langue ?

\begin{definition}[L'emprunt]
Un \cle{emprunt} est un mot ou une expression pris à une langue étrangère et intégré, avec ou sans adaptation, dans la langue d'accueil. On distingue trois procédés :
\begin{itemize}
\item l'\textbf{emprunt direct} : le mot est repris tel quel (\emph{planning}, \emph{weekend}, \emph{briefing}) ;
\item l'\textbf{emprunt adapté} : le mot est francisé dans sa graphie ou sa prononciation (\emph{wagon}, venu de l'anglais \emph{waggon}, a perdu un \emph{g} ; \emph{match} est prononcé à la française) ;
\item le \textbf{calque} : on traduit mot à mot une expression étrangère (\emph{gratte-ciel} calque l'anglais \emph{skyscraper}).
\end{itemize}
\end{definition}

\begin{experience}[Observation d'un corpus : où sont les emprunts ?]
\textbf{Protocole :} une classe relève, dans un corpus de 40 mots empruntés à l'anglais trouvés dans des articles de la presse camerounaise, le domaine auquel chaque mot appartient. \textbf{Résultats :} 18 mots relèvent de l'informatique (\emph{software}, \emph{mail}, \emph{login}...), 10 du commerce (\emph{marketing}, \emph{deal}...), 8 du sport (\emph{coach}, \emph{fair-play}...) et 4 de l'administration (\emph{staff}, \emph{briefing}...). \textbf{Vérification :} $18 + 10 + 8 + 4 = 40$ mots, soit respectivement 45 \%, 25 \%, 20 \% et 10 \% du corpus. \textbf{Interprétation :} les emprunts pénètrent d'abord les domaines techniques et nouveaux, où le français n'a pas toujours créé d'équivalent ; l'administration, domaine codifié, résiste davantage. C'est précisément dans l'écrit administratif que l'emprunt est le moins toléré.
\end{experience}

\begin{popfigure}
\begin{center}
\begin{tikzpicture}
\begin{axis}[
    ybar, bar width=0.9cm,
    width=0.85\linewidth, height=6cm,
    symbolic x coords={Informatique, Commerce, Sport, Administration},
    xtick=data,
    ymin=0, ymax=20,
    ylabel={Nombre de mots relevés},
    xlabel={Domaine d'usage},
    nodes near coords,
    every node near coord/.append style={font=\small},
    ymajorgrids=true, grid style={popline!40},
]
\addplot[fill=popblue!70, draw=popdark] coordinates {(Informatique,18) (Commerce,10) (Sport,8) (Administration,4)};
\end{axis}
\end{tikzpicture}
\end{center}
\legende{Fréquence des emprunts à l'anglais par domaine, d'après un corpus de 40 mots relevés dans la presse camerounaise (enquête de classe). L'informatique concentre près de la moitié des emprunts (45 \%), l'administration seulement 10 \%.}
\end{popfigure}

\courssub{Emprunts et français camerounais : usages acceptables et limites}

Au Cameroun, le français vit au contact de l'anglais et des langues nationales (ewondo, douala, fulfuldé...). De ce contact naissent des emprunts légitimes dans la \textbf{conversation} : « il m'a \emph{briefé} », « on se voit au \emph{weekend} », et même des mots du \textbf{camfranglais} comme « le ndem », « être en galère avec les \emph{mbeng} ». Mais l'écrit utilitaire formel obéit à d'autres règles.

\begin{aretenir}[La règle d'or pour la requête]
Dans une requête administrative ou professionnelle :
\begin{itemize}
\item on \textbf{proscrit} totalement le camfranglais et les emprunts non adaptés (\emph{briefing}, \emph{deadline}, \emph{ok}) ;
\item on \textbf{tolère} les emprunts intégrés au dictionnaire français quand aucun équivalent courant n'existe (\emph{weekend} reste admis, mais « fin de semaine » est préférable) ;
\item on \textbf{préfère} toujours l'équivalent français recommandé lorsqu'il existe.
\end{itemize}
\end{aretenir}

\begin{popfigure}
\begin{center}
\begin{tikzpicture}[font=\small]
\draw[fill=popblueL, draw=popdark, rounded corners] (0,0) rectangle (7.2,3.4);
\draw[fill=popgreenL, draw=popdark, rounded corners] (7.8,0) rectangle (15,3.4);
\node[font=\bfseries] at (3.6,2.9) {Emprunt à éviter};
\node[font=\bfseries] at (11.4,2.9) {Équivalent recommandé};
\node[align=center] at (3.6,2.1) {deadline};
\node[align=center] at (11.4,2.1) {date limite / délai};
\node[align=center] at (3.6,1.4) {briefing};
\node[align=center] at (11.4,1.4) {réunion d'information};
\node[align=center] at (3.6,0.7) {planning};
\node[align=center] at (11.4,0.7) {emploi du temps / calendrier};
\draw[->, thick, poporange] (7.25,2.1) -- (7.75,2.1);
\draw[->, thick, poporange] (7.25,1.4) -- (7.75,1.4);
\draw[->, thick, poporange] (7.25,0.7) -- (7.75,0.7);
\end{tikzpicture}
\end{center}
\legende{Tableau de correspondance : chaque emprunt courant possède un équivalent français soutenu, exigible dans une requête.}
\end{popfigure}

\begin{methode}[Choisir le bon niveau de lexique]
\begin{enumerate}
\item \textbf{Identifier le destinataire} : un préfet, un directeur des ressources humaines ou un chef de service exigent le registre soutenu et le lexique administratif.
\item \textbf{Sélectionner les termes du domaine} : repérer dans le champ lexical de sa filière les termes exacts correspondant à sa demande (certificat de travail, bon de commande, stage académique...).
\item \textbf{Traquer les emprunts} : relire chaque phrase ; pour tout mot d'origine étrangère, vérifier dans un dictionnaire français s'il est intégré ; sinon, le remplacer par l'équivalent recommandé.
\item \textbf{Bannir le camfranglais} : toute expression parlée (« je suis dans le ndem », « le dossier a été \emph{chapé} ») doit être traduite en français standard (« je suis dans une situation difficile », « le dossier a disparu »).
\item \textbf{Se relire à voix haute} : si une phrase ressemble à une conversation de rue, elle n'a pas sa place dans la requête.
\end{enumerate}
\end{methode}

\begin{attention}[Erreur fréquente : le calque fautif]
Écrire « Je suis au \textbf{chômage technique} depuis \textbf{le mois passé} » est un double calque fautif de l'anglais : on dira correctement « Je suis \textbf{en chômage technique} depuis \textbf{le mois dernier} » (ou mieux, dans une requête : « mon contrat a été suspendu pour chômage technique depuis le mois d'octobre »). De même, « je vais \emph{checker} le dossier » doit devenir « je vais \textbf{vérifier} le dossier ». Règle pratique : vérifie chaque emprunt dans un dictionnaire avant de l'écrire.
\end{attention}

\begin{savaistu}[Le mot « ok »]
« OK » est probablement l'emprunt le plus répandu au monde : né aux États-Unis au XIX\textsuperscript{e} siècle (abréviation populaire de \emph{oll korrect}, orthographe plaisante de \emph{all correct}), il est compris sur tous les continents et s'est intégré à des centaines de langues, y compris au camfranglais. Pourtant, malgré cette carrière mondiale, il reste \textbf{proscrit dans une requête administrative} : on lui préférera « d'accord », « conforme » ou « approuvé » selon le contexte.
\end{savaistu}

\courssub{Application : réécriture d'une requête fautive}

\begin{exoresolu}[Assainir une requête pleine d'emprunts]
\textbf{Énoncé.} Voici l'extrait d'une requête rédigée par un élève : « Je soussigné Mbarga Jean, je vous \emph{briefe} que je suis dans le \emph{ndem} : mon \emph{planning} de stage a été changé sans me prévenir. La \emph{deadline} pour déposer mon rapport, c'est le \emph{weekend} prochain. Merci de me \emph{checker} ça vite, \emph{ok} ? » Réécris cet extrait en français soutenu, adapté à une requête administrative.
\tcblower
\textbf{Solution détaillée.} On repère d'abord les fautes de registre : \emph{briefe} (emprunt verbalisé), \emph{ndem} (camfranglais), \emph{planning} et \emph{deadline} (emprunts non adaptés), \emph{checker} et \emph{ok} (proscrits). On remplace chacun par un équivalent soutenu et on reconstruit les formules consacrées :

« Je soussigné Mbarga Jean, domicilié à Yaoundé, quartier Nkolbisson, ai l'honneur de porter à votre connaissance que je me trouve dans une situation difficile : mon emploi du temps de stage a été modifié sans information préalable. Or, la date limite de dépôt de mon rapport de stage est fixée à la fin de la semaine prochaine. C'est pourquoi je sollicite votre bienveillante intervention pour l'examen rapide de ma situation. Dans l'espoir d'une suite favorable, je vous prie d'agréer, Monsieur le Directeur, l'expression de ma haute considération. »

Chaque emprunt a été remplacé : \emph{briefe} $\rightarrow$ « porter à votre connaissance » ; \emph{ndem} $\rightarrow$ « situation difficile » ; \emph{planning} $\rightarrow$ « emploi du temps » ; \emph{deadline} $\rightarrow$ « date limite » ; \emph{weekend} $\rightarrow$ « fin de la semaine » ; \emph{checker ça vite} $\rightarrow$ « examen rapide de ma situation » ; \emph{ok} $\rightarrow$ supprimé et remplacé par la formule de politesse.
\end{exoresolu}

\begin{exercice}[À toi de jouer]
Réécris en français soutenu la phrase suivante, extraite d'une demande d'emploi : « Je vous envoie mon \emph{CV} pour le \emph{job} de comptable, parce que franchement je suis au \emph{chômage} depuis le mois passé et je gère bien le \emph{marketing} aussi. » (On corrigera le calque fautif et on remplacera tous les emprunts non adaptés.)
\end{exercice}

\begin{corrige}[Exercice : réécriture]
Proposition de correction : « J'ai l'honneur de vous adresser mon curriculum vitae en vue du poste de comptable vacant dans votre établissement. En recherche d'emploi depuis le mois dernier, je maîtrise en outre les techniques de mercatique (ou : les techniques commerciales). » Points attendus : \emph{job} $\rightarrow$ « poste » ; « depuis le mois passé » $\rightarrow$ « depuis le mois dernier » ; « je suis au chômage » $\rightarrow$ « en recherche d'emploi » ; \emph{marketing} $\rightarrow$ « mercatique » ou « techniques commerciales » ; \emph{CV} est accepté (abréviation intégrée) mais « curriculum vitae » est plus soutenu.
\end{corrige}

\begin{aretenir}[L'essentiel de la section]
Le \cle{lexique spécialisé} est le vocabulaire technique d'un domaine : il donne précision et crédibilité à la requête. La requête administrative possède des formules figées (« je soussigné », « pièces jointes », « suite favorable ») à reproduire exactement. L'\cle{emprunt} (direct, adapté ou calque) est toléré à l'oral mais doit être remplacé par son équivalent français dans tout écrit utilitaire formel ; le camfranglais y est totalement proscrit. Avant d'envoyer une requête : identifier le destinataire, mobiliser les termes du domaine, vérifier chaque emprunt au dictionnaire.
\end{aretenir}

\coursec{Rédaction guidée : entête, introduction, corps et conclusion}

Dans la section précédente, nous avons enrichi notre vocabulaire : nous savons désormais choisir un \cle{lexique spécialisé} adapté à l'administration et reconnaître les \cle{emprunts} qui peuplent notre langue administrative (dossier, formulaire, récépissé). Mais un beau vocabulaire ne suffit pas : encore faut-il savoir le disposer dans le texte. Une requête, comme une maison, se construit dans un ordre précis : on ne pose pas le toit avant les fondations. Cette section vous guide pas à pas dans la rédaction effective des quatre parties de la requête : l'\cle{entête}, l'\cle{introduction}, le \cle{corps} et la \cle{conclusion}.

\courssub{L'entête : la carte d'identité de la requête}

\textbf{Intuition.} Imaginez que vous receviez une lettre sans nom d'expéditeur, sans date, sans adresse. Que feriez-vous ? Vous la jetteriez, car vous ne pourriez ni vérifier son authenticité ni y répondre. L'entête joue exactement ce rôle : il rend la requête \emph{identifiable}, \emph{datable} et \emph{acheminable}.

\begin{definition}[L'entête]
L'\cle{entête} est la partie supérieure de la requête. Il comprend obligatoirement quatre éléments, dans cet ordre : le \cle{lieu et la date} de rédaction, l'\cle{identité complète du requérant} (nom, prénom, date et lieu de naissance, profession, adresse), la \cle{mention exacte du destinataire} (titre officiel complet) et l'\cle{objet} de la demande, résumé en une seule ligne. La mention des pièces jointes, bien que facultative à cet emplacement, est fortement recommandée pour la clarté du dossier.
\end{definition}

\textbf{Observation.} Prenez trois requêtes authentiques déposées dans une sous-préfecture : vous constaterez que toutes commencent par la même formule de datation (« Yaoundé, le ... »), que le requérant se présente toujours à la troisième personne (« Le soussigné... ») et que l'objet tient en une ligne. Cette régularité n'est pas un hasard : elle obéit à un code administratif que toute administration camerounaise attend.

\begin{exemplebox}[Modèle complet d'entête]
\begin{center}
\begin{tabularx}{\linewidth}{|l|X|}
\hline
\rowcolor{popblueL} \textbf{Élément} & \textbf{Rédaction} \\
\hline
Lieu et date & Douala, le 14 novembre 2026. \\
\hline
Identité du requérant & Le soussigné MBAPPE Arnaud, né le 12/03/1990 à Bafoussam, commerçant, BP 4021 Douala, tél. 699 00 00 00. \\
\hline
Destinataire & À Monsieur le Maire de la Communauté Urbaine de Douala. \\
\hline
Objet & Objet : demande de renouvellement de patente. \\
\hline
Pièces jointes & Pièces jointes : 3 (copie de l'ancienne patente, quitus fiscal, photocopie de la CNI). \\
\hline
\end{tabularx}
\end{center}
\end{exemplebox}

\begin{attention}[Pièges de l'entête]
Trois erreurs reviennent sans cesse dans les copies :
\begin{itemize}
\item écrire la date sans le lieu (« le 14 novembre 2026 » seul est insuffisant) ;
\item oublier la profession ou l'adresse du requérant, ce qui empêche toute réponse ;
\item rédiger un objet de trois lignes : l'objet doit tenir en \textbf{une seule ligne}, sans verbe conjugué (on écrit « demande de renouvellement », pas « je demande le renouvellement »).
\end{itemize}
\end{attention}

\courssub{Les formules d'introduction : ouvrir la porte avec respect}

\textbf{Intuition.} Dans nos traditions, on n'entre jamais chez un chef sans saluer d'abord. La formule d'introduction est la salutation protocolaire de la requête : elle annonce qui parle, à qui, et pour quoi, tout en marquant le respect dû au rang du destinataire.

\begin{definition}[La formule d'introduction]
La \cle{formule d'introduction} est la première phrase du texte proprement dit. Elle présente le requérant et annonce la demande, selon des tournures consacrées : « J'ai l'honneur de solliciter de votre haute bienveillance... » ou « Le soussigné, ..., a l'honneur de venir très respectueusement auprès de votre autorité... ».
\end{definition}

\begin{exemplebox}[Deux modèles d'introduction]
\begin{itemize}
\item \emph{Modèle à la première personne} : « J'ai l'honneur de solliciter de votre haute bienveillance le renouvellement de ma patente de commerçant pour l'année fiscale 2027. »
\item \emph{Modèle à la troisième personne} : « Le soussigné MBAPPE Arnaud, commerçant au marché Mboppi, a l'honneur de venir très respectueusement auprès de votre autorité solliciter le renouvellement de sa patente. »
\end{itemize}
Les deux modèles sont corrects ; il faut simplement en choisir un et s'y tenir jusqu'à la fin de la requête (cohérence de la personne).
\end{exemplebox}

\courssub{Le corps de la requête : raconter, prouver, demander}

\textbf{Intuition.} Le corps est le cœur de la requête. C'est là que vous convainquez. Or on ne convainc pas une administration avec des émotions, mais avec des \emph{faits datés}, des \emph{preuves jointes} et une \emph{demande claire}. Un bon corps de requête ressemble à un dossier d'enquête : chronologique, factuel, documenté.

\begin{methode}[Rédiger le corps de la requête en trois étapes]
\begin{enumerate}
\item \textbf{Rappeler les faits avec dates et lieux} : exposez la situation dans l'ordre chronologique, en phrases courtes, sans commentaire inutile. Chaque fait doit être localisé et daté (« Le 03 janvier 2021, j'ai obtenu ma première patente au service des impôts de Douala 1\textsuperscript{er} »).
\item \textbf{Joindre les preuves} : annoncez chaque pièce justificative au moment où vous citez le fait qu'elle prouve (« ... comme l'atteste la copie de l'ancienne patente ci-jointe »).
\item \textbf{Formuler la demande en une phrase claire} : terminez le corps par une phrase unique qui dit exactement ce que vous attendez (« C'est pourquoi je sollicite le renouvellement de ma patente pour l'année 2027 »).
\end{enumerate}
Le ton doit rester \textbf{respectueux et ferme} : respectueux dans les formules, ferme dans l'exposé des droits et des faits.
\end{methode}

\begin{attention}[Ton et registre]
Erreur fréquente : mélanger les formules (« je vous prie d'agréer l'expression de ma \emph{haute considération distinguée} », qui n'existe pas) ou tutoyer implicitement le destinataire par un ton trop familier (« vous savez bien que je paie toujours mes impôts »). On vouvoie toujours, on ne plaisante jamais, on n'accuse personne sans preuve. La fermeté s'exprime par les faits, jamais par l'agressivité.
\end{attention}

\courssub{La conclusion : refermer la porte avec protocole}

\textbf{Intuition.} Après avoir plaidé votre cause, il faut prendre congé. La formule de conclusion n'est pas une simple politesse : c'est un code qui indique à l'administration que vous maîtrisez les usages, donc que votre demande est sérieuse.

\begin{definition}[La formule de conclusion]
La \cle{formule de conclusion} exprime l'espoir d'une suite favorable et la considération due au destinataire. Modèle type : « Dans l'espoir d'une suite favorable, je vous prie d'agréer, Monsieur le Préfet, l'expression de ma haute considération. » Elle est suivie de la \cle{signature} du requérant.
\end{definition}

\begin{propriete}[Concordance des titres]
La formule de politesse finale reprend \textbf{toujours} le titre exact du destinataire mentionné dans l'entête. Si l'entête dit « À Monsieur le Préfet du Wouri », la conclusion dira « je vous prie d'agréer, \emph{Monsieur le Préfet}, ... ». Toute discordance (« Monsieur le Maire » dans l'entête, « Monsieur le Préfet » dans la conclusion) est une faute de cohérence sanctionnée à l'examen.
\end{propriete}

\textbf{Adapter la formule au rang du destinataire.} La hiérarchie administrative impose une gradation des formules : plus le destinataire est élevé dans le protocole, plus la formule est solennelle.

\begin{popfigure}
\begin{center}
\begin{tikzpicture}[scale=0.85, font=\small]
\draw[fill=popblueL, rounded corners=2pt] (0,3.6) rectangle (14,4.6);
\node[font=\bfseries] at (7,4.1) {Formules de politesse selon le destinataire};
\draw[fill=popgoldL] (0,2.7) rectangle (4.5,3.5);
\node[font=\bfseries, align=center] at (2.25,3.1) {Destinataire (Rang protocolaire)};
\draw[fill=popgoldL] (4.5,2.7) rectangle (14,3.5);
\node[font=\bfseries, align=center] at (9.25,3.1) {Formule de conclusion adaptée};
\draw[fill=popcream] (0,1.8) rectangle (4.5,2.7);
\node[align=center] at (2.25,2.25) {Préfet, Gouverneur,\\ Ministre, Secrétaire d'État};
\draw[fill=white] (4.5,1.8) rectangle (14,2.7);
\node[align=center] at (9.25,2.25) {« Dans l'espoir d'une suite favorable, je vous prie d'agréer, [Titre], l'expression de ma \textbf{haute considération}. »};
\draw[fill=popcream] (0,0.9) rectangle (4.5,1.8);
\node[align=center] at (2.25,1.35) {Directeur, Maire,\\ Chef d'établissement, Chef de service};
\draw[fill=white] (4.5,0.9) rectangle (14,1.8);
\node[align=center] at (9.25,1.35) {« Dans l'espoir d'une suite favorable, je vous prie d'agréer, [Titre], l'expression de ma \textbf{considération distinguée}. »};
\draw[fill=popcream] (0,0) rectangle (4.5,0.9);
\node[align=center] at (2.25,0.45) {Chef d'entreprise,\\ Responsable privé, Particulier};
\draw[fill=white] (4.5,0) rectangle (14,0.9);
\node[align=center] at (9.25,0.45) {« Dans l'espoir d'une suite favorable, je vous prie d'agréer, [Titre], mes \textbf{salutations distinguées}. »};
\end{tikzpicture}
\end{center}
\legende{Gradation des formules de politesse finales selon le rang protocolaire du destinataire : 1. haute considération (autorités administratives supérieures) ; 2. considération distinguée (cadres dirigeants, élus locaux, chefs de service) ; 3. salutations distinguées (secteur privé, relations courantes).}
\end{popfigure}

\begin{savaistu}[Une formalité qui coûte cher]
Au Cameroun, une requête sans date ni signature est \textbf{irrecevable} devant la plupart des administrations : elle n'est même pas enregistrée au bureau d'ordre. De nombreux usagers ont perdu des mois de démarches pour avoir oublié de dater ou de signer leur demande. La date fixe aussi le point de départ des délais légaux de réponse de l'administration : sans elle, impossible de faire valoir vos droits.
\end{savaistu}

\courssub{Vue d'ensemble : l'architecture complète de la requête}

Le schéma en escalier ci-dessous synthétise le parcours de rédaction, avec le nombre de lignes recommandé pour chaque étape dans une requête tenant sur une page.

\begin{popfigure}
\begin{center}
\begin{tikzpicture}[font=\small, scale=0.9]
\draw[fill=popblueL, rounded corners=3pt] (0,3.6) rectangle (6.5,4.6);
\node[align=center] at (3.25,4.1) {\textbf{1. Entête}\\ 5 à 7 lignes};
\draw[-{Stealth}, very thick, popdark] (6.5,4.1) -- (7.5,4.1);
\draw[fill=popgreenL, rounded corners=3pt] (7.5,2.7) rectangle (14,3.7);
\node[align=center] at (10.75,3.2) {\textbf{2. Introduction}\\ 1 à 2 lignes};
\draw[-{Stealth}, very thick, popdark] (10.75,2.7) -- (10.75,2.1);
\draw[fill=poporangeL, rounded corners=3pt] (7.5,1.1) rectangle (14,2.1);
\node[align=center] at (10.75,1.6) {\textbf{3. Corps}\\ 8 à 12 lignes};
\draw[-{Stealth}, very thick, popdark] (7.5,1.6) -- (6.5,1.6);
\draw[fill=poppinkL, rounded corners=3pt] (0,0.5) rectangle (6.5,1.6);
\node[align=center] at (3.25,1.05) {\textbf{4. Conclusion}\\ 2 lignes + signature};
\end{tikzpicture}
\end{center}
\legende{Schéma en escalier de la requête : entête (5 à 7 lignes) $\rightarrow$ introduction (1 à 2 lignes) $\rightarrow$ corps (8 à 12 lignes) $\rightarrow$ conclusion (2 lignes + signature). Le corps est la partie la plus développée.}
\end{popfigure}

\begin{exoresolu}[Compléter une requête à trous]
Énoncé. Voici la fin d'une requête adressée au Préfet du Mfoundi. Complétez-la correctement : « Dans l'espoir d'une suite favorable, je vous prie d'agréer, .................., l'expression de .................. »
\tcblower
Solution détaillée. Deux vérifications s'imposent. Premièrement, la concordance des titres : l'entête mentionne « À Monsieur le Préfet du Mfoundi », donc la conclusion reprend « Monsieur le Préfet ». Deuxièmement, la gradation protocolaire : un Préfet est une autorité administrative supérieure, donc la formule exigée est « ma haute considération ». Réponse : « Dans l'espoir d'une suite favorable, je vous prie d'agréer, \textbf{Monsieur le Préfet}, l'expression de \textbf{ma haute considération}. »
\end{exoresolu}

\courssub{Atelier d'écriture en sous-groupes : l'assemblage collectif}

\begin{experience}[Atelier de rédaction collaborative]
\textbf{Situation donnée} : « Votre coopérative de mototaxis de Bafoussam demande au Maire l'autorisation d'occuper un emplacement près du marché B. »
\begin{enumerate}
\item Répartissez la classe en quatre sous-groupes. Le groupe A rédige l'entête complet (lieu, date, identité du président de la coopérative, destinataire, objet, pièces jointes). Le groupe B rédige la formule d'introduction. Le groupe C rédige le corps en trois étapes (faits datés, pièces jointes, demande en une phrase). Le groupe D rédige la conclusion et prépare la signature.
\item Chaque groupe présente sa partie au tableau (5 minutes par groupe).
\item Assemblage collectif : la classe lit la requête complète à voix haute et vérifie la cohérence (concordance des titres, gradation de la formule finale, chronologie des faits, respect de la personne choisie).
\item Amélioration : chaque groupe propose une correction sur la partie d'un autre groupe, puis la version définitive est recopiée au propre.
\end{enumerate}
\textbf{Observation attendue} : les élèves constatent qu'une requête réussie est un texte où chaque partie dialogue avec les autres — l'objet annonce le corps, le corps justifie la demande, la conclusion fait écho à l'entête.
\end{experience}

\begin{aretenir}[L'essentiel de la section]
\begin{itemize}
\item L'entête comprend quatre éléments obligatoires : lieu et date, identité complète du requérant, destinataire exact, objet en une ligne (pièces jointes recommandées).
\item L'introduction annonce la demande par une formule consacrée, à la première ou à la troisième personne (cohérence maintenue).
\item Le corps suit trois étapes : faits datés et localisés, preuves jointes, demande en une phrase claire ; ton respectueux et ferme.
\item La conclusion reprend le titre exact du destinataire et adapte la formule à son rang : haute considération, considération distinguée ou salutations distinguées.
\item Sans date ni signature, la requête est irrecevable.
\end{itemize}
\end{aretenir}

\begin{exercice}[À vous de rédiger]
À partir de la situation suivante : « Vous êtes NGONO Aline, coiffeuse à Yaoundé, et vous sollicitez du Directeur des Impôts de votre centre un délai de paiement de votre patente », rédigez : 1) l'entête complet ; 2) la formule d'introduction ; 3) la formule de conclusion adaptée au rang du destinataire. Vous justifierez le choix de la formule finale.
\end{exercice}

\begin{corrige}[Exercice : à vous de rédiger]
1) Entête : « Yaoundé, le 20 novembre 2026. La soussignée NGONO Aline, née le 05/07/1995 à Ebolowa, coiffeuse, BP 5120 Yaoundé. À Monsieur le Directeur des Impôts du Centre des impôts de Yaoundé Centre. Objet : demande de délai de paiement de patente. Pièces jointes : 2 (échéancier proposé, justificatif d'activité). » 2) Introduction : « J'ai l'honneur de solliciter de votre haute bienveillance un délai de paiement de ma patente pour l'année fiscale 2026. » 3) Conclusion : « Dans l'espoir d'une suite favorable, je vous prie d'agréer, Monsieur le Directeur, l'expression de ma considération distinguée. » Justification : le Directeur des Impôts est un cadre dirigeant (chef de service déconcentré), non une autorité préfectorale ; la formule « considération distinguée » convient donc, et le titre « Monsieur le Directeur » reprend exactement celui de l'entête (concordance).
\end{corrige}

\coursec{Exploitation des textes iconiques pour préparer la requête}

Dans la section précédente, nous avons appris à rédiger, pas à pas, l'entête, les formules d'introduction, le corps et la conclusion d'une requête. Mais une question demeure : comment savoir \emph{à qui} adresser la requête, \emph{quel objet} lui donner et \emph{quelles références} mentionner ? Ces informations ne se devinent pas : elles se lisent, le plus souvent, sur des documents affichés dans les administrations — organigrammes, panneaux d'affichage, formulaires. Ces documents, où l'image et la disposition visuelle portent l'essentiel du message, sont appelés des \cle{textes iconiques}. Cette section vous apprend à les lire méthodiquement pour préparer une requête exacte et efficace.

\courssub{Qu'est-ce qu'un texte iconique ?}

\begin{definition}[Texte iconique]
Un \cle{texte iconique} est un document dont le message repose principalement sur des éléments \emph{visuels} : image, schéma, tableau, organigramme, panneau, signalétique, formulaire. Le lecteur n'y lit pas un discours suivi, mais il \emph{parcourt} le document en suivant des repères graphiques : flèches, cadres, couleurs, titres, rubriques.
\end{definition}

L'intuition est simple : lorsque vous entrez dans une préfecture ou une mairie, personne ne vous récite un discours pour vous indiquer où aller. Vous levez les yeux vers un panneau, vous lisez un organigramme accroché au mur, vous remplissez un formulaire. L'information est là, mais elle est \emph{organisée visuellement}. Savoir lire ces documents est donc une compétence de la vie courante autant que scolaire : au Baccalauréat technique, on peut vous demander d'exploiter un tel document pour rédiger une requête.

\begin{aretenir}[Les grands types de textes iconiques utiles à la requête]
\begin{itemize}
\item L'\cle{organigramme} : schéma qui représente la hiérarchie d'une administration ou d'une entreprise (qui commande qui, quel service dépend de quel chef).
\item Le \cle{panneau d'affichage} : support qui présente des avis, annonces, décisions, horaires, avis de recrutement.
\item Le \cle{formulaire} : document à rubriques (nom, date, objet, signature) qu'il faut remplir avec exactitude.
\item Le \cle{tableau} : présentation de données en lignes et colonnes (tarifs, dates limites, pièces à fournir).
\end{itemize}
\end{aretenir}

\courssub{Lecture du Texte iconique n°1 : l'organigramme d'une préfecture}

\textbf{Observation préalable.} Imaginez que vous vouliez obtenir la régularisation d'un acte de naissance. À qui écrire ? Au Préfet ? Au Sous-préfet ? Au chef de service des affaires administratives ? L'organigramme répond à cette question : il montre la \cle{voie hiérarchique}, c'est-à-dire le chemin que doit suivre votre demande pour aboutir au bon destinataire.

\begin{popfigure}
\begin{tikzpicture}[
  every node/.style={font=\small},
  box/.style={draw=popblue, fill=popblueL, rounded corners=2pt, align=center, minimum width=4.2cm, minimum height=0.9cm},
  sbox/.style={draw=poppurple, fill=poppurpleL, rounded corners=2pt, align=center, minimum width=3.6cm, minimum height=0.9cm},
  ubox/.style={draw=popgreen, fill=popgreenL, rounded corners=2pt, align=center, minimum width=3.2cm, minimum height=0.8cm},
  fleche/.style={-{Stealth[length=2.5mm]}, thick, popdark}
]
\node[box, align=center] (prefet) at (0,0){\textbf{Préfet}\\ (autorité administrative du département)};
\node[box, fill=poporangeL, draw=poporange, align=center] (sp) at (0,-1.8){\textbf{Sous-préfet}\\ (arrondissement)};
\node[sbox, align=center] (cs1) at (-4.6,-3.6){Chef de service\\ des affaires administratives};
\node[sbox, align=center] (cs2) at (0,-3.6){Chef de service\\ de l'état civil};
\node[sbox, align=center] (cs3) at (4.6,-3.6){Chef de service\\ des affaires économiques};
\node[ubox] (us) at (0,-5.2) {\textbf{Usagers}};
\draw[fleche] (prefet) -- (sp);
\draw[fleche] (sp) -- (cs1);
\draw[fleche] (sp) -- (cs2);
\draw[fleche] (sp) -- (cs3);
\draw[fleche] (cs1) -- (us);
\draw[fleche] (cs2) -- (us);
\draw[fleche] (cs3) -- (us);
\draw[fleche, poporange, dashed] (us) -- node[right, font=\footnotesize\itshape, text=poporange, xshift=0.3cm]{voie de la requête} (sp);
\end{tikzpicture}
\legende{Organigramme simplifié d'une préfecture camerounaise. Les flèches pleines indiquent la subordination (du sommet vers la base) ; la flèche orange en pointillés indique le chemin que suit la requête de l'usager vers l'autorité compétente.}
\end{popfigure}

\begin{methode}[Lire un organigramme en quatre étapes]
\begin{enumerate}
\item \textbf{Partir du sommet hiérarchique} : repérer l'autorité la plus élevée (ici, le Préfet). C'est souvent le destinataire formel de la requête, même si un service subordonné traite le dossier.
\item \textbf{Suivre les flèches de subordination} : descendre de cadre en cadre pour comprendre qui dépend de qui.
\item \textbf{Repérer le service compétent} : identifier le service dont l'intitulé correspond à votre problème (état civil pour un acte de naissance, affaires économiques pour un commerce).
\item \textbf{Formuler le destinataire exact} : écrire dans l'entête l'autorité et, si utile, la mention « à l'attention du chef de service de... ».
\end{enumerate}
\end{methode}

\begin{attention}[Erreur fréquente : le mauvais destinataire]
L'erreur la plus coûteuse est d'adresser la requête au mauvais service faute d'avoir lu l'organigramme : elle est alors réorientée, retardée, parfois classée sans suite. Avant toute rédaction, vérifiez toujours la \emph{compétence} du destinataire. Une requête concernant un litige de quartier relève du Sous-préfet ou du Préfet, pas du maire ; une requête concernant un marché communal relève du maire, pas du Préfet.
\end{attention}

\begin{exemplebox}[Une donnée qui fait réfléchir]
Lors d'une simulation menée en classe de Terminale, on a suivi le parcours de 30 requêtes adressées à un service préfectoral. Résultat : près d'un tiers, soit environ 10 requêtes, a dû être réorienté pour erreur de destinataire. Conséquences : perte de temps pour l'usager (parfois plusieurs semaines), surcharge de travail pour les services, et dossier qui risque de passer après les délais. Lire l'organigramme avant d'écrire, c'est gagner du temps pour soi et pour l'administration.
\end{exemplebox}

\begin{savaistu}[La hiérarchie administrative au Cameroun]
Au Cameroun, l'organisation territoriale et administrative de l'État est fixée par la Constitution et les lois sur la décentralisation (lois n°2004/017, 2004/018, 2004/019 modifiées). Ces textes définissent la chaîne hiérarchique que doit respecter toute requête administrative : l'usager s'adresse à l'autorité de proximité (Sous-préfet pour l'État, Maire pour la Commune), et la demande peut remonter, par voie hiérarchique, jusqu'au Préfet (chef de département), au Gouverneur (chef de région), voire au ministre compétent. Sauter un échelon sans motif expose la requête à être renvoyée à l'échelon inférieur.
\end{savaistu}

\courssub{Lecture du Texte iconique n°2 : le panneau d'affichage administratif}

Le second type de texte iconique à maîtriser est le panneau d'affichage : c'est là que l'administration publie les avis de recrutement, les décisions, les appels à candidatures, les dates limites. Sa lecture exige de la méthode, car l'information est dense et hiérarchisée visuellement.

\begin{popfigure}
\begin{tikzpicture}[font=\small]
\draw[fill=popcream, draw=popdark, thick] (0,0) rectangle (11,7.4);
\draw[fill=popblue, draw=popblue] (0,6.6) rectangle (11,7.4);
\node[white, font=\bfseries] at (5.5,7.0) {COMMUNE DE BAFOUSSAM -- PANNEAU D'AFFICHAGE OFFICIEL};
\draw[fill=white, draw=popline] (0.4,3.9) rectangle (6.6,6.3);
\node[align=left, anchor=north west, font=\footnotesize] at (0.55,6.2) {\textbf{AVIS DE RECRUTEMENT N°04/2026}\\ La Mairie de Bafoussam recrute :\\ -- 2 techniciens de génie civil ;\\ -- 1 comptable.\\ Dossier : CV, acte de naissance, diplôme.\\ \textbf{Date limite : 30 avril 2026.}};
\draw[fill=white, draw=popline] (7.0,3.9) rectangle (10.6,6.3);
\node[align=left, anchor=north west, font=\footnotesize] at (7.15,6.2) {\textbf{HORAIRES D'OUVERTURE}\\ Lundi--vendredi : 7 h 30 -- 15 h 30\\ Service des affaires administratives :\\ guichet n°2, bâtiment B};
\draw[fill=white, draw=popline] (0.4,0.3) rectangle (10.6,3.5);
\node[align=left, anchor=north west, font=\footnotesize] at (0.55,3.4) {\textbf{PIÈCES À FOURNIR POUR TOUTE REQUÊTE}\\ 1. Requête manuscrite ou dactylographiée, timbrée ; 2. Photocopie de la CNI ;\\ 3. Justificatif du motif (facture, attestation, contrat). Dépôt au secrétariat central, bâtiment A.};
\draw[-{Stealth[length=2.5mm]}, thick, poporange] (6.8,5.2) -- (7.6,2.2) node[midway, right, font=\footnotesize\itshape, text=poporange, align=left]{références utiles :\\ lieu, guichet, horaires};
\draw[-{Stealth[length=2.5mm]}, thick, poppink] (3.4,3.8) -- (3.4,2.6) node[midway, right, font=\footnotesize\itshape, text=poppink]{pièces à joindre};
\draw[-{Stealth[length=2.5mm]}, thick, popgreen] (2.0,6.5) -- (2.0,5.9) node[midway, left, font=\footnotesize\itshape, text=popgreen]{objet et date limite};
\end{tikzpicture}
\legende{Reproduction schématique d'un panneau d'affichage administratif. Flèche verte : l'objet et la date limite ; flèche rose : les pièces à fournir ; flèche orange : les références pratiques (lieu, guichet, horaires).}
\end{popfigure}

\begin{methode}[Extraire les données utiles d'un panneau ou d'un formulaire]
\begin{enumerate}
\item \textbf{Lire d'abord le bandeau-titre} : il indique l'autorité émettrice (la commune, la préfecture), donc le destinataire probable de la requête.
\item \textbf{Relever les références exactes} : numéro de l'avis, date de publication, intitulé du service. Ces références figureront dans l'entête ou dans l'objet de la requête.
\item \textbf{Noter les dates limites} : une requête hors délai est irrecevable ; la date limite conditionne le calendrier de rédaction et de dépôt.
\item \textbf{Lister les pièces exigées} : elles seront mentionnées dans le corps de la requête (« pièces jointes : ... »).
\item \textbf{Repérer le lieu et le guichet de dépôt} : utile pour la formule de conclusion et pour le suivi du dossier.
\end{enumerate}
\end{methode}

\courssub{Croisement image/texte : de l'organigramme à l'entête de la requête}

Tout l'art de la préparation consiste à \emph{transformer} une information visuelle en élément rédactionnel. Le tableau ci-dessous montre ce travail de conversion.

\begin{center}
\begin{tabularx}{\linewidth}{|l|X|X|}\hline
\rowcolor{popblueL} \textbf{Information lue sur le document iconique} & \textbf{Emplacement dans la requête} & \textbf{Formulation exacte} \\ \hline
Autorité au sommet de l'organigramme & Destinataire (entête, en haut à droite) & « À Monsieur le Préfet du département du... » \\ \hline
Service compétent repéré sur l'organigramme & Mention d'attention & « À l'attention du chef de service de l'état civil » \\ \hline
Numéro et date de l'avis affiché & Objet ou références & « Objet : demande d'acte de naissance (avis n°04/2026 du 12 mars 2026) » \\ \hline
Pièces listées sur le panneau & Corps de la requête & « Pièces jointes : photocopie de la CNI, attestation... » \\ \hline
Date limite affichée & Corps ou conclusion & « Dans l'attente d'une suite avant le 30 avril 2026... » \\ \hline
\end{tabularx}
\end{center}

\begin{exoresolu}[Exploiter un avis affiché pour rédiger un entête]
\textbf{Énoncé.} Sur le panneau de la mairie de Bafoussam figure l'avis suivant : « Avis n°12/2026 -- Les usagers souhaitant régulariser leur acte de naissance sont invités à adresser une requête au service de l'état civil avant le 15 mai 2026. Pièces : CNI, attestation de naissance. » L'organigramme de la commune indique : Maire $\rightarrow$ Secrétaire général $\rightarrow$ Chef de service de l'état civil. Rédigez l'entête complet (expéditeur, destinataire, lieu et date, objet) de la requête de Mme Kamga Aline, demeurant à Bafoussam, quartier Tougang.
\tcblower
\textbf{Solution détaillée.}
\begin{enumerate}
\item \textbf{Lecture de l'organigramme} : le service compétent est l'état civil, qui dépend du Maire. Le destinataire formel est donc le Maire, avec mention du service.
\item \textbf{Relevé des références} : avis n°12/2026, date limite du 15 mai 2026, pièces exigées (CNI, attestation).
\item \textbf{Rédaction de l'entête} :
\end{enumerate}
\begin{quote}
Kamga Aline\\
Quartier Tougang, Bafoussam\\
Tél. : 6XX XX XX XX\\[0.4em]
\emph{À Monsieur le Maire de la Commune de Bafoussam}\\
\emph{À l'attention du chef de service de l'état civil}\\[0.4em]
Bafoussam, le 20 avril 2026\\[0.4em]
\textbf{Objet :} demande de régularisation d'acte de naissance (avis n°12/2026).\\
\textbf{Pièces jointes :} photocopie de la CNI, attestation de naissance.
\end{quote}
L'entête respecte la structure externe de la requête et exploite fidèlement toutes les données des deux textes iconiques : destinataire exact, service compétent, références de l'avis, pièces annoncées. La date de rédaction (20 avril) précède la date limite (15 mai) : la requête est recevable.
\end{exoresolu}

\courssub{TP guidé : de l'avis de recrutement à l'esquisse de la requête}

\begin{experience}[Travaux pratiques : lire pour écrire]
\textbf{Situation.} Un avis de recrutement est affiché à la mairie de Bafoussam : « Avis de recrutement n°04/2026. La Commune de Bafoussam recrute deux techniciens de génie civil et un comptable. Les candidats adresseront une requête à Monsieur le Maire avant le 30 avril 2026. Dossier : CV, acte de naissance, photocopie du diplôme. Renseignements au guichet n°2, bâtiment B. »

\textbf{Matériel.} Reproduction de l'avis (ou affichage réel), organigramme de la commune, brouillon, grille de relevé.

\textbf{Protocole.}
\begin{enumerate}
\item Relevez sur une fiche : l'autorité émettrice, le numéro de l'avis, les postes offerts, la date limite, les pièces exigées, le lieu de renseignement.
\item À l'aide de l'organigramme, vérifiez que le Maire est bien le destinataire compétent.
\item Esquissez l'entête complet de la requête (expéditeur, destinataire, lieu et date, objet).
\item Formulez l'objet en une phrase précise mentionnant le poste visé et la référence de l'avis.
\end{enumerate}

\textbf{Observations attendues.} Chaque élève dispose d'un entête dont tous les éléments proviennent des documents iconiques, et non de l'imagination.

\textbf{Interprétation.} La requête administrative n'est pas un texte inventé : c'est un texte \emph{alimenté} par des documents officiels. La qualité de la lecture conditionne la qualité de l'écriture.
\end{experience}

\begin{exercice}[À vous de jouer]
À partir du même avis de recrutement n°04/2026, rédigez l'entête et l'objet de la requête de M. Fotso Ulrich, technicien en génie civil, demeurant à Bafoussam, quartier Banengo. Vous indiquerez également, en deux lignes, pourquoi il serait incorrect d'adresser cette requête au Gouverneur de la région de l'Ouest.
\end{exercice}

\begin{corrige}[Exercice]
\textbf{Entête attendu :} Fotso Ulrich, quartier Banengo, Bafoussam ; \emph{À Monsieur le Maire de la Commune de Bafoussam} ; Bafoussam, le [date antérieure au 30 avril 2026] ; \textbf{Objet :} candidature au poste de technicien de génie civil (avis de recrutement n°04/2026) ; pièces jointes : CV, acte de naissance, photocopie du diplôme.

\textbf{Justification :} le Gouverneur n'est pas l'autorité émettrice de l'avis ; l'organigramme communal montre que le recrutement relève du Maire. Adresser la requête au Gouverneur, c'est sauter la voie hiérarchique : la demande serait réorientée ou classée sans suite, avec risque de dépasser la date limite.
\end{corrige}

\begin{aretenir}[L'essentiel de la section]
\begin{itemize}
\item Un \cle{texte iconique} (organigramme, panneau, formulaire, tableau) porte son information par l'image et la disposition visuelle.
\item L'\cle{organigramme} se lit du sommet vers la base : il révèle le destinataire compétent et la voie hiérarchique.
\item Le \cle{panneau d'affichage} fournit les références exactes : numéro d'avis, dates limites, pièces à fournir.
\item Chaque information visuelle se convertit en un élément précis de la requête : destinataire, objet, références, pièces jointes.
\item Une requête adressée au mauvais service est une requête perdue : lire avant d'écrire.
\end{itemize}
\end{aretenir}

Ces compétences de lecture étant acquises, la section suivante vous conduira à l'étape décisive : analyser une situation complète, élaborer le plan détaillé et rédiger intégralement votre requête.

\coursec{Analyse de situation, plan détaillé et rédaction complète}

\courssub{Transition et rappel des acquis}

Après avoir exploité les textes iconiques (schémas d'organigramme, photos de matériel défectueux, captures d'écran de factures erronées) lors de la séance précédente, nous disposons maintenant de \textbf{matériaux concrets} pour étayer nos arguments. L'image n'est plus une simple illustration : elle devient une \emph{preuve visuelle} que l'on référence dans le corps de la requête (ex. : « \emph{Comme le montre la photo ci-jointe, référence Fig. 1, la courroie du tour CNC présente une usure anormale au point de repère 3} »).

Nous franchissons ici l'étape décisive : passer de la \emph{préparation} à la \emph{production} effective. Tout comme un technicien ne démarre pas l'usinage sans avoir validé son programme ISO et vérifié son brut, le rédacteur ne doit pas écrire la première phrase sans avoir \textbf{analysé la situation} et \textbf{figé son plan détaillé}. Cette section vous guide pas à pas, de la grille QQOQCP à la rédaction finale, en intégrant l'auto-évaluation critériée conforme aux attendus du Baccalauréat Technique.

\begin{methode}[Analyser une situation de communication avant de rédiger : la démarche QQOQCP]
\textbf{Objectif :} Transformer une consigne floue en cahier des charges précis pour la rédaction.
\begin{enumerate}
    \item \textbf{Qui ?} Identifier l'expéditeur (rôle, service, matricule) et le destinataire (nom, fonction, pouvoir de décision). \emph{Ex. : « Moi, NGOUMOU Jean, technicien maintenance atelier B, matricule 452, à M. Le Directeur de l'Usine. »}
    \item \textbf{Quoi ?} Nommer l'objet précis : réclamation, demande de stage, signalement de panne, demande de congé. Un seul objet par requête.
    \item \textbf{Où ?} Situer le lieu : atelier, bureau, chantier, établissement d'accueil (pour stage).
    \item \textbf{Quand ?} Dater les faits (panne du 12/03), la rédaction (fait à Douala, le 15/03/2026), les délais souhaités (réponse sous 48h, stage du 01/07 au 31/08).
    \item \textbf{Comment ?} Décrire le mode opératoire ou les circonstances : procédure respectée, dysfonctionnement constaté, pièces jointes (photos, devis, CV).
    \item \textbf{Pourquoi ?} Expliciter la finalité : obtenir réparation, correction de facture, convention de stage, remplacement pièce. C'est la \cle{cause profonde} qui justifie la \cle{demande finale}.
\end{enumerate}
\textbf{Astuce :} Remplissez un tableau QQOQCP \emph{au brouillon} avant toute phrase. Chaque case remplie = un paragraphe ou un argument du corps de la requête.
\end{methode}

\courssub{Analyse de la situation de communication : application de la grille QQOQCP}

Prenons une situation technique réaliste : \textbf{une panne critique sur une machine-outil à commande numérique (CN) bloquant la production d'une série de pièces pour un client majeur}.

\begin{popfigure}
\begin{tikzpicture}[
    node distance=1.2cm and 1.5cm,
    main/.style={draw=popblue, thick, fill=popblueL, rounded corners, minimum width=3.2cm, minimum height=1cm, align=center, font=\small},
    sub/.style={draw=poporange, fill=poporangeL, rounded corners, minimum width=2.8cm, minimum height=0.9cm, align=center, font=\footnotesize},
    arrow/.style={-Stealth, thick, popdark},
    qnode/.style={draw=poppurple, thick, fill=poppurpleL, circle, minimum size=1.2cm, font=\bfseries\small}
]
% Central node
\node[qnode] (center) {QQOQCP};

% Questions nodes
\node[main, above left=of center] (qui) {QUI ?};
\node[main, above=of center] (quoi) {QUOI ?};
\node[main, above right=of center] (ou) {OÙ ?};
\node[main, right=of center] (quand) {QUAND ?};
\node[main, below right=of center] (comment) {COMMENT ?};
\node[main, below=of center] (pourquoi) {POURQUOI ?};

% Answers - Qui
\node[sub, left=of qui] (qui_r) {Expéditeur : \textbf{Ngoumou Jean}\newline Technicien Maintenance Niv. 3\newline Matricule : 452 - Atelier B\newline Destinataire : \textbf{M. Le Directeur Général}\newline Usine MECACAM Douala};

% Answers - Quoi
\node[sub, above=of quoi] (quoi_r) {OBJET : \textbf{Requête en réparation urgente}\newline Tour CNC Fanuc 0i-TF (N° série 7845)\newline Panne broche : vibration + code alarme 417};

% Answers - Où
\node[sub, right=of ou] (ou_r) {Lieu : \textbf{Atelier B - Poste 4}\newline Usine MECACAM, Zone Industrielle Bassa\newline Pièce concernée : Broche principale};

% Answers - Quand
\node[sub, below=of quand] (quand_r) {Faits : \textbf{12/03/2026 - 06h15} (démarrage poste)\newline Rédaction : \textbf{13/03/2026}\newline Délai souhaité : \textbf{Intervention < 24h}\newline Livraison client : \textbf{20/03/2026}};

% Answers - Comment
\node[sub, right=of comment, align=center] (comment_r){Circonstances :\newline 1. Démarrage routine OK\\ 2. Usinage pièce 12/50 : bruit anormal\\ 3. Arrêt urgence + Alarme 417 (Surcharge)\\ 4. Vérif. niveau huile + filtres : OK\\ 5. Photos capteur vibration + écran alarme jointes};

% Answers - Pourquoi
\node[sub, left=of pourquoi] (pourquoi_r) {Finalité :\newline - Arrêt ligne = 500 000 FCFA/jour perte\newline - Pénalité retard client : 10\%/jour\newline - Risque casse outil (coût 150 000 FCFA)\newline DEMANDE : Intervention tech. Fanuc + pièce rechange};

% Arrows
\draw[arrow] (center) -- (qui); \draw[arrow] (center) -- (quoi); \draw[arrow] (center) -- (ou);
\draw[arrow] (center) -- (quand); \draw[arrow] (center) -- (comment); \draw[arrow] (center) -- (pourquoi);
\draw[arrow, popgreen] (qui) -- (qui_r); \draw[arrow, popgreen] (quoi) -- (quoi_r);
\draw[arrow, popgreen] (ou) -- (ou_r); \draw[arrow, popgreen] (quand) -- (quand_r);
\draw[arrow, popgreen] (comment) -- (comment_r); \draw[arrow, popgreen] (pourquoi) -- (pourquoi_r);

% Legend box
\node[draw=popline, fill=popcream, rounded corners, below left=-0.5cm and -2cm of center, anchor=north west, align=left, font=\footnotesize, text width=12cm] (leg) {
    \textbf{Légende :} Le cercle central (violet) = Outil d'analyse. Bleu = Questions standards. Orange = Réponses contextualisées (Données techniques réelles). Flèches vertes = Transfert vers le plan de rédaction.
};
\end{tikzpicture}
\legende{Schéma de la grille QQOQCP appliquée à une panne machine en entreprise (MECACAM Douala). Chaque réponse alimente directement une zone du plan détaillé.}
\end{popfigure}

\begin{aretenir}[Résultat de l'analyse QQOQCP]
L'analyse révèle trois \textbf{enjeux majeurs} qui structureront le corps de la requête :
\begin{itemize}
    \item \textbf{Urgence économique :} Perte journalière + pénalités contractuelles (chiffres clés).
    \item \textbf{Traçabilité technique :} Alarme 417, code erreur précis, procédure de premier niveau effectuée (huile, filtres).
    \item \textbf{Preuves matérielles :} Photos capteur, écran alarme, carnet de maintenance à jour.
\end{itemize}
Ces éléments passent du stade « informations brutes » à « arguments structurés » pour le corps de la lettre.
\end{aretenir}

\courssub{Élaboration du plan détaillé : de l'analyse à la structure interne}

Le plan détaillé est le \emph{squelette} de la requête. Il respecte la structure interne canonique : \textbf{Introduction -- Corps -- Conclusion}. Chaque section du plan correspond à une fonction illocutoire précise (informer, argumenter, demander).

\begin{definition}[Plan détaillé type pour une requête technique (Réclamation/Panne)]
\begin{enumerate}
    \item \textbf{ENTÊTE} (Structure externe) : Expéditeur / Destinataire / Lieu, Date / Objet (Précis + Référence machine/dossier).
    \item \textbf{INTRODUCTION} (1 paragraphe) :
        \begin{itemize}
            \item Formule d'appel (Monsieur le Directeur,).
            \item \textbf{Identité + Qualité} : « Je soussigné..., technicien maintenance... »
            \item \textbf{Motif immédiat} : « J'ai l'honneur de vous signaler... » / « Par la présente, je porte à votre connaissance... »
        \end{itemize}
    \item \textbf{CORPS} (2 à 3 paragraphes argumentés) -- \emph{Le cœur technique} :
        \begin{itemize}
            \item \textbf{§1 : Exposé des faits datés et localisés} (Le « Quoi/Quand/Où/Comment »). Chronologie précise. Vocabulaire technique exact (alarme 417, broche, roulements, vibromètre).
            \item \textbf{§2 : Analyse technique et diagnostic préliminaire} (Le « Comment » approfondi). Preuves des vérifications faites (niveau huile, filtres). Hypothèse : roulements broche / variateur. \textbf{Preuves visuelles jointes} (Photos Fig. 1, 2).
            \item \textbf{§3 : Conséquences chiffrées et Risques} (Le « Pourquoi » économique). Coût arrêt/jour, pénalités client, sécurité personnel. Justifie l'urgence.
        \end{itemize}
    \item \textbf{CONCLUSION} (1 paragraphe) :
        \begin{itemize}
            \item \textbf{Demande explicite, chiffrée, datée} : « Je sollicite l'intervention du technicien Fanuc avant le 14/03 14h + commande roulements ref. 7012-C-TR. »
            \item \textbf{Disponibilité} : « Je reste à disposition pour tout complément. »
            \item Formule de politesse de clôture (adaptée au destinataire).
            \item Signature + Pièces jointes listées.
        \end{itemize}
\end{enumerate}
\end{definition}

\begin{attention}[Piège fatal : rédiger sans plan détaillé]
\textbf{Erreur fréquente :} Se lancer dans l'écriture « au fil de l'eau » pour gagner du temps.
\textbf{Conséquences au Bac :}
\begin{itemize}
    \item Oubli de la référence machine (zéro point structure externe).
    \item Arguments mélangés : faits, diagnostic et conséquences dans le même paragraphe (incohérence structure interne - perte de 3 à 4 pts).
    \item Demande floue : « Je vous prie de faire le nécessaire » (refus implicite, 0 pt conclusion).
    \item Lexique courant au lieu du lexique spécialisé (perte 2-3 pts lexique).
\end{itemize}
\textbf{Règle d'or :} \textbf{10 minutes d'analyse + plan au brouillon} = 40 minutes de rédaction fluide et notée. Sans plan, c'est 50 minutes de réécriture et une note < 8/20.
\end{attention}

\courssub{Rédaction guidée complète : cas de la panne machine (Exercice résolu)}

Voici la rédaction finale issue du plan ci-dessus. Observez l'alignement strict Plan $\rightarrow$ Texte.

\begin{exoresolu}[Rédaction complète : Requête de réparation urgente Tour CNC]
\textbf{Consigne :} À partir de l'analyse QQOQCP (figure ci-dessus), rédigez la requête complète. Respectez la mise en page administrative.

\tcblower

\begin{center}
\textbf{REPUBLIQUE DU CAMEROUN}\\
\textbf{Paix -- Travail -- Patrie}\\
\vspace{0.3cm}
\textbf{USINE MECACAM DOUALA}\\
\textbf{Direction Générale}\\
\vspace{0.5cm}
\rule{\linewidth}{0.5pt}
\vspace{0.3cm}
\textbf{REQUÊTE EN RÉPARATION URGENTE}\\
\textbf{TOUR CNC FANUC 0i-TF (N° SÉRIE 7845 -- POSTE 4 ATELIER B)}
\end{center}

\vspace{0.5cm}
\noindent
\textbf{Expéditeur :} NGOUMOU Jean\\
Technicien Maintenance Niveau 3 -- Matricule 452\\
Atelier B -- Poste 4\\
Tél. interne : 2245

\vspace{0.3cm}
\noindent
\textbf{Destinataire :} Monsieur le Directeur Général\\
Usine MECACAM -- Zone Industrielle Bassa\\
Douala

\vspace{0.3cm}
\noindent
Fait à Douala, le 13 mars 2026

\vspace{0.5cm}
\noindent
\textbf{Objet :} Signalement panne critique Tour CNC Fanuc 0i-TF (N° 7845) -- Demande intervention urgente.

\vspace{0.5cm}
\noindent
Monsieur le Directeur Général,

\vspace{0.3cm}
\noindent
Je soussigné, \textbf{NGOUMOU Jean}, technicien de maintenance affecté à l'Atelier B (matricule 452), ai l'honneur de porter à votre connaissance une panne critique survenue sur le \textbf{Tour à Commande Numérique Fanuc 0i-TF, numéro de série 7845, installé au Poste 4 de l'Atelier B}, qui immobilise totalement la production de la série de référence \textbf{PIECE-A-2026} depuis le \textbf{12 mars 2026, 06h15}.

\vspace{0.3cm}
\noindent
\textbf{1. Exposé des faits et constat technique.}
Le \textbf{12 mars 2026 à 06h15}, lors du démarrage de la série (pièce n°12 sur 50), la machine a émis des vibrations anormales au niveau de la broche principale, suivies de l'affichage de l'\textbf{alarme 417 (Surcharge servo-moteur broche)} sur l'écran opérateur. Conformément à la procédure de premier niveau (Fiche MNT-PROC-04), j'ai procédé aux vérifications suivantes : niveau d'huile hydraulique (conforme), état des filtres à air et huile (propres), absence de corps étranger dans la zone de travail. Malgré un reset système, l'alarme 417 réapparaît instantanément à la mise en rotation (M03 S1500). Les relevés du vibromètre portable (Modèle VM-6360) indiquent un pic à \textbf{12,5 mm/s} (seuil alarme ISO 10816 : 7,1 mm/s pour cette classe de machine).

\vspace{0.3cm}
\noindent
\textbf{2. Diagnostic préliminaire et pièces jointes.}
Au vu du code alarme 417 couplé aux mesures vibratoires excessives, l'hypothèse privilégiée est une \textbf{défaillance des roulements à billes de contact oblique de la broche (Réf. constructeur : 7012-C-TR-P4)} ou, plus rarement, un défaut du variateur de fréquence (Drive SPMC-15i). \textbf{Deux photographies} (Pièces jointes n°1 et 2) documentent l'écran d'alarme et la position du capteur de vibration lors de la mesure. Le carnet de maintenance de la machine (cote MNT-7845) est à jour ; la dernière révision préventive (niveau 2) date du 15 janvier 2026.

\vspace{0.3cm}
\noindent
\textbf{3. Impact économique et risques opérationnels.}
L'arrêt de ce poste unique pour cette référence engendre une \textbf{perte de production estimée à 500 000 FCFA/jour} (coût machine + main-d'œuvre + overhead). Le contrat client CLT-2026-045 impose une livraison le \textbf{20 mars 2026} sous peine de \textbf{pénalités de retard de 10\% par jour de calendrier} (soit 1 200 000 FCFA/jour). Par ailleurs, la poursuite des essais en mode dégradé risquerait d'endommager irréversiblement le porte-outil et la pièce (coût outillage : 150 000 FCFA), sans compter le risque pour la sécurité de l'opérateur.

\vspace{0.5cm}
\noindent
Au vu de ce qui précède, \textbf{je sollicite respectueusement} :
\begin{enumerate}
    \item La programmation d'une \textbf{intervention du technicien agréé FANUC (Hotline contrat n° FNC-2024-012) avant le 14 mars 2026, 14h00}.
    \item Le lancement immédiat de la commande des \textbf{roulements de broche (Réf. 7012-C-TR-P4, quantité 2)} auprès du fournisseur agréé (SOTRAMAC Douala), délai annoncé 48h.
    \item L'autorisation de basculer la série sur le Tour Conventionnel MAZAK (Poste 2, Atelier A) après validation du Chef d'Atelier A, pour limiter le retard.
\end{enumerate}

\vspace{0.3cm}
\noindent
Je reste à votre entière disposition pour toute information complémentaire ou pour accompagner l'expert Fanuc lors de son intervention.

\vspace{0.5cm}
\noindent
Je vous prie d'agréer, Monsieur le Directeur Général, l'expression de ma considération distinguée.

\vspace{1cm}
\noindent
\begin{tabular}{ll}
\textbf{Pièces jointes :} & 1. Photo écran Alarme 417 (12/03/2026 06h18)\\
& 2. Photo mesure vibromètre (12,5 mm/s)\\
& 3. Extrait carnet maintenance MNT-7845\\
& 4. Bon de commande prévisionnel roulements (brouillon)
\end{tabular}

\vspace{1.5cm}
\noindent
\begin{minipage}{0.45\linewidth}
\hrulefill\\
\textbf{NGOUMOU Jean}\\
Technicien Maintenance\\
Signature :
\end{minipage}
\hfill
\begin{minipage}{0.45\linewidth}
\hrulefill\\
\textbf{Visa Chef Atelier B}\\
Nom : \dots\hfill Date : \dots\\
Signature :
\end{minipage}
\end{exoresolu}

\begin{aretenir}[Analyse de la rédaction modèle (Pourquoi cette note ?)]
\begin{itemize}
    \item \textbf{Structure externe (4/4) :} En-tête complet (République, Entité, Qualité), Objet précis (Machine + N° série + Localisation + Nature), Lieu/Date, Formules d'appel/clôture adaptées, Signature + Visa hiérarchique, Pièces jointes numérotées.
    \item \textbf{Structure interne (6/6) :} Intro (Identité+Motif) distincte. Corps en 3 § logiques : Faits techniques datés $\rightarrow$ Diagnostic argumenté + Preuves $\rightarrow$ Enjeux économiques chiffrés. Conclusion : Demande \textbf{explicite, chiffrée, datée, alternative} (basculement Mazak).
    \item \textbf{Lexique spécialisé (4/4) :} Alarme 417, Broche, Vibromètre, ISO 10816, Roulements 7012-C-TR-P4, Variateur SPMC-15i, Reset, M03 S1500, Hotline, Niveau 2, Overhead. Zéro langage familier.
    \item \textbf{Correction (4/4) :} Syntaxe administrative (subjonctif « je sollicite que... », passif « l'alarme est affichée »), accords parfaits, ponctuation rigoureuse.
    \item \textbf{Présentation (2/2) :} Aération, alignements, gras sur mots-clés techniques, numérotation demandes.
\end{itemize}
\textbf{Total : 20/20.}
\end{aretenir}

\courssub{Auto-évaluation critériée : la grille sur 20 points (Standard Bac Tech)}

L'auto-évaluation n'est pas une option, c'est une \textbf{compétence évaluée}. Vous devez être capable de noter votre propre production comme un correcteur. Voici la grille officielle type, à mémoriser et à appliquer systématiquement.

\begin{popfigure}
\begin{center}
\begin{tabularx}{\linewidth}{|>{\bfseries}l|X|c|}
\hline
\rowcolor{popblue} \textbf{\color{white}Critère} & \textbf{\color{white}Descriptif / Indicateurs} & \textbf{\color{white}Pts} \\
\hline
Structure externe (4 pts) & En-tête complet (Rép., Entité, Expéd./Destin., Lieu/Date), Objet précis (Réf. technique), Formules appel/clôture adaptées, Signature + Visa + PJ listées. & 4 \\
\hline
Structure interne -- Introduction (2 pts) & Identité + Qualité + Motif clair. Formule d'appel correcte. Transition vers le corps. & 2 \\
\hline
Structure interne -- Corps (6 pts) & 3 § distincts : 1.Faits datés/localisés 2.Diagnostic/Preuves 3.Contrôles/Enjeux. Chronologie, vocabulaire technique, arguments hiérarchisés, références PJ. & 6 \\
\hline
Structure interne -- Conclusion (4 pts) & Demande EXPLICITE, CHIFFREE, DATEE, ALTERNATIVE. Disponibilité. Formule clôture. Signature. & 4 \\
\hline
Lexique spécialisé \& Emprunts (4 pts) & Termes techniques exacts (normes, codes, réf.), registres adaptés (admin + tech), emprunts intégrés (ex: reset, hotline, deadline), précision lexicale. & 4 \\
\hline
Correction langagière (4 pts) & Syntaxe administrative (subjonctif, passif, impersonnel), accords, temps, ponctuation, orthographe lexicale/grammaticale. Zéro solécisme. & 4 \\
\hline
Présentation \& Mise en page (2 pts) & Aération paragraphes, alignements, lisibilité, gras/italique sur mots-clés, numérotation demandes, propreté. & 2 \\
\hline
\rowcolor{popgold} \multicolumn{2}{|c|}{\textbf{TOTAL}} & \textbf{20} \\
\hline
\end{tabularx}
\end{center}
\legende{Grille d'auto-évaluation critériée type Baccalauréat Technique (20 points). L'élève coche les descripteurs atteints.}
\end{popfigure}

\begin{exemplebox}[Application de la grille : Auto-correction guidée]
\textbf{Exercice :} Vous avez rédigé une requête pour une \emph{erreur de facturation fournisseur (double facturation matière première)}. Vous vous auto-évaluez.
\begin{itemize}
    \item \textbf{Structure externe :} Vous avez mis « Objet : Facture wrong ». \rightarrow \textbf{1/4} (Objet imprécis : pas de n° facture, date, montant).
    \item \textbf{Corps :} Vous avez mélangé la description de l'erreur et la demande de remboursement dans un seul paragraphe long. \rightarrow \textbf{3/6} (Manque de structure §, arguments noyés).
    \item \textbf{Lexique :} Vous avez écrit « l'argent » au lieu de « le montant indû », « le fournisseur » au lieu de « le cocontractant / le créancier ». \rightarrow \textbf{2/4}.
    \item \textbf{Correction :} 2 fautes d'accords, 1 solécisme (« je souhaite que vous me remboursez »). \rightarrow \textbf{2/4}.
    \item \textbf{Présentation :} Correcte. \rightarrow \textbf{2/2}.
    \item \textbf{Total : 10/20.} \textbf{Action corrective :} Revoir l'objet, découper le corps en 3 §, remplacer le lexique courant par le lexique comptable/juridique, corriger le subjonctif.
\end{itemize}
\end{exemplebox}

\courssub{Justification écrite de la démarche : expliciter ses choix}

Au-delà de la note, le Bac Technique (épreuve écrite ou orale de contrôle) peut exiger une \textbf{justification métacognitive} : « Expliquez pourquoi vous avez choisi cette formule d'appel, ce terme technique, cette organisation en trois paragraphes. » C'est la preuve de votre \cle{conscience professionnelle}.

\begin{methode}[Rédiger la justification de sa démarche (Modèle type)]
Structurez votre justification en 3 temps :
\begin{enumerate}
    \item \textbf{Choix de la situation d'énonciation (Formules) :} « J'ai choisi "Monsieur le Directeur Général" car le pouvoir de décision budgétaire (commande pièces, intervention externe) relève de sa signature. La formule de clôture "considération distinguée" est la norme administrative hiérarchique descendante. »
    \item \textbf{Choix du lexique spécialisé (Précision technique) :} « J'ai utilisé "Alarme 417" et non "code erreur" car c'est la terminologie constructeur Fanuc. "Roulements 7012-C-TR-P4" permet au magasinier de commander sans ambiguïté. L'emprunt "Hotline" (anglais technique) est consacré dans les contrats de maintenance. »
    \item \textbf{Choix de l'organisation (Plan) :} « Le §1 isole les faits bruts (traçabilité légale). Le §2 montre ma compétence (diagnostic + procédure respectée) et engage la responsabilité de l'expert. Le §3 traduit le technique en langage de gestion (FCFA, pénalités) pour décider le DG. La conclusion propose une solution de contournement (Mazak) prouvant mon sens de l'initiative. »
\end{enumerate}
\end{methode}

\begin{exoresolu}[Exemple de justification écrite (Extrait)]
\textbf{Consigne :} Justifiez en 150 mots l'organisation du corps de votre requête (panne Tour CNC).

\tcblower
\textbf{Justification de l'organisation du corps de la requête}

L'organisation en trois paragraphes distincts répond à une \textbf{logique argumentative progressive} adaptée au destinataire décideur.

Le \textbf{premier paragraphe} établit la \textbf{traçabilité factuelle incontestable} : date/heure précises (12/03 06h15), référence machine normalisée (Fanuc 0i-TF N°7845), code alarme constructeur (417), mesures physiques objectives (12,5 mm/s vs norme ISO 10816). Ce choix lexical technique (alarme, vibromètre, reset, M03) vise à \textbf{crédibiliser l'expéditeur} comme expert technique et à figer l'événement dans le temps pour la maintenance et l'assurance.

Le \textbf{deuxième paragraphe} opère le \textbf{passage du constat au diagnostic} (hypothèse roulements 7012-C-TR-P4). L'intégration des pièces jointes (photos, carnet maintenance) transforme l'affirmation en \textbf{dossier probant}. La mention de la révision préventive récente (15/01) écarte la négligence entretien.

Le \textbf{troisième paragraphe} \textbf{traduit la technique en enjeux économiques} (500k FCFA/j, pénalités 1,2M FCFA/j, risque sécurité). Ce registre gestionnaire est le seul audible par le Directeur Général pour déclencher le budget urgence. La conclusion propose une \textbf{alternative opérationnelle} (basculement Mazak), démontrant que le technicien ne subit pas la panne mais gère la continuité.
\end{exoresolu}

\begin{savaistu}[Le Bac Technique : Conformité au modèle = Langue]
Au Baccalauréat Technique, l'épreuve de production d'écrits utilitaires (Requête, Rapport, Note de service) est notée sur \textbf{la conformité au modèle administratif attendu} autant que sur la qualité de la langue.
\begin{itemize}
    \item Un texte en parfait français mais sans objet précis, sans pièces jointes listées, sans visa hiérarchique, sans demande datée $\rightarrow$ \textbf{Note maximale : 8/20} (Échec structure externe/interne).
    \item Un texte avec 3 fautes d'orthographe mais structure externe parfaite, plan interne respecté (3 §), lexique technique précis, demande explicite chiffrée $\rightarrow$ \textbf{Note : 14--15/20} (Réussite honorable).
\end{itemize}
\textbf{Enseignement :} La \textbf{forme administrative (le contenant)} est une condition \textbf{sine qua non} pour que le \textbf{fond technique (le contenu)} soit lu et noté. Maîtrisez le « squelette » avant d'habiller la phrase.
\end{savaistu}

\courssub{Deux autres situations types : Entraînement ciblé}

Pour être prêt le jour J, vous devez savoir décliner le même schéma (QQOQCP $\rightarrow$ Plan $\rightarrow$ Rédaction $\rightarrow$ Auto-éval) sur les deux autres grands types de requêtes du programme.

\begin{exercice}[Situation 1 : Erreur de facturation fournisseur (Gestion/Comptabilité)]
\textbf{Contexte :} Vous êtes Gestionnaire des Achats (Matricule 112) à la SOCATRAL. La facture n° F-2026-0445 du fournisseur \emph{ACIERS DU CAMEROUN} (livraison 10/03) facture 50 tonnes de tôle 3mm à 850 000 FCFA/T au lieu du prix contractuel 780 000 FCFA/T (Contrat CAD-2025-012). Écart : 3 500 000 FCFA.
\begin{enumerate}
    \item Remplissez la grille QQOQCP (5 min).
    \item Élaborez le plan détaillé (3 § corps : Constat écart / Référence contractuelle / Demande avoir).
    \item Rédigez la requête complète (20 min).
    \item Auto-évaluez sur 20 avec la grille.
\end{enumerate}
\end{exercice}

\begin{exercice}[Situation 2 : Demande de stage de fin d'études (Insertion professionnelle)]
\textbf{Contexte :} Élève Terminale TEE (Électrotechnique), Matricule 2026-TEE-045. Vous ciblez l'entreprise \emph{ENEO Douala-Bonabéri} pour un stage de 8 semaines (02/06 -- 25/07/2026) sur la maintenance des postes HTA/B. Vous avez le PV de soutenance de votre projet de fin d'études (PFE) sur « Protection différentielle réseaux HTA ».
\begin{enumerate}
    \item QQOQCP : Qui (Élève + École) $\rightarrow$ Qui (DRH ENEO). Quoi (Demande stage). Où (Poste HTA Bonabéri). Quand (Dates + Délai réponse). Comment (CV + Lettre motivation + PFE + Convention école). Pourquoi (Spécialisation PFE + Embauche potentielle).
    \item Plan corps : 1. Présentation cursus + PFE pertinent. 2. Adéquation stage / besoins ENEO (compétences relais). 3. Disponibilité + Cadre légal (Convention, Assurance).
    \item Rédigez. Auto-évaluez. \textbf{Point de vigilance :} La formule d'appel est « Madame la Directrice des Ressources Humaines, » ; la conclusion demande un \textbf{entretien} + envoi \textbf{convention tripartite}.
\end{enumerate}
\end{exercice}

\begin{corrige}[Exercice 1 -- Erreur facturation (Éléments attendus)]
\textbf{Objet :} Réclamation facture n° F-2026-0445 -- Écart prix contrat CAD-2025-012 -- Demande avoir 3 500 000 FCFA.
\textbf{Corps §1 :} Constat : Facture 50T $\times$ 850k = 42 500 000 FCFA. Bon de réception BL-2026-0110 confirme 50T livrées.
\textbf{Corps §2 :} Référence : Contrat CAD-2025-012 Art. 4.2 fixe prix ferme 780 000 FCFA/T HT. Calcul juste : 39 000 000 FCFA. Écart = 3 500 000 FCFA.
\textbf{Corps §3 :} Demande : Émission avoir n° [à définir] de 3 500 000 FCFA sous 8 jours. Règlement solde 39 000 000 FCFA à échéance 30j. Maintien relation contractuelle.
\textbf{Lexique :} Avoir, Écart, Prix ferme, Échéance, Bon de réception, Article contractuel, Solde.
\end{corrige}

\begin{corrige}[Exercice 2 -- Demande stage (Éléments attendus)]
\textbf{Objet :} Demande de stage de fin d'études (02/06 -- 25/07/2026) -- Maintenance Postes HTA/B -- Spécialisation Protection Différentielle.
\textbf{Corps §1 :} Élève TEE Lycée Technique X. PFE : « Protection différentielle réseaux HTA » (Note 16/20). Maîtrise relais Sepam, protection 51N/50N.
\textbf{Corps §2 :} Intérêt ENEO : Maintenance préventive relais numériques, analyse déclenchements, plans de consignation. Compétences immédiatement opérationnelles.
\textbf{Corps §3 :} Cadre : Convention tripartite (École/Entreprise/Étudiant) prête. Assurance responsabilité civile étudiante valide. Disponible pour entretien semaine 15--19/04.
\textbf{Lexique :} Convention tripartite, HTA/B, Relais de protection (Sepam), Consignation, Préventive/Curative, PFE, Tripartite.
\end{corrige}

\begin{attention}[Derniers pièges à éviter avant l'examen]
\begin{itemize}
    \item \textbf{Confondre Requête et Lettre de motivation :} La requête est un acte \textbf{administratif descendant/horizontal} (droit, procédure, preuve). La lettre de motivation est un acte \textbf{de séduction} (projet, personnalité). Ne mélangez pas les codes.
    \item \textbf{Oublier le « Visa » hiérarchique :} Dans l'entreprise, une requête technique monte toujours par la hiérarchie (Chef d'équipe $\rightarrow$ Chef d'atelier $\rightarrow$ Directeur). L'absence du visa du supérieur immédiat = procédure non respectée = requête irrecevable (0 pt structure externe).
    \item \textbf{Pièces jointes « fantômes » :} Ne listez que ce qui existe vraiment. « CV joint » sans CV = perte de crédibilité (perception correcteur).
    \item \textbf{Gestion du temps :} 10 min Analyse/Plan -- 35 min Rédaction -- 5 min Relecture/Grille. Tenez le chrono.
\end{itemize}
\end{attention}

\begin{aretenir}[Check-list finale « Prêt pour le Bac »]
Avant de rendre votre copie, cochez mentalement :
\fbox{\textbf{Structure Externe}} $\square$ République/Entité $\square$ Expéditeur/Destinataire complets $\square$ Lieu/Date $\square$ Objet : [Type] + [Réf. technique] + [Action demandée] $\square$ Formules appel/clôture $\square$ Signature + Visa + PJ numérotées.
\fbox{\textbf{Structure Interne}} $\square$ Intro : Identité + Motif $\square$ Corps : 3 § (Faits $\rightarrow$ Analyse/Preuves $\rightarrow$ Enjeux) $\square$ Conclusion : Demande Explicite/Chiffrée/Datée/Alternative $\square$ Transition fluide §.
\fbox{\textbf{Lexique}} $\square$ 5--7 termes techniques exacts $\square$ Registre admin + technique $\square$ Emprunts intégrés corrects $\square$ Zéro vague.
\fbox{\textbf{Langue}} $\square$ Subjonctif (demande) $\square$ Passif (constat) $\square$ Accords $\square$ Ponctuation $\square$ Orthographe lexicale.
\fbox{\textbf{Présentation}} $\square$ Aération $\square$ Alignements $\square$ Gras mots-clés $\square$ Propreté.
\end{aretenir}

\coursec{Confrontation, amélioration des productions, compte rendu et remédiation}

Dans la section précédente, vous avez analysé une situation, élaboré un plan détaillé et rédigé une requête complète. Mais un texte administratif ne vaut que par sa qualité finale : une requête mal présentée ou truffée de fautes a peu de chances d'aboutir, même si la demande est juste. Cette dernière section vous apprend donc à regarder votre propre production avec les yeux d'un correcteur : échanger, annoter, corriger, réécrire. C'est exactement la démarche que le correcteur du Probatoire et du Baccalauréat technique attend de vous, et c'est aussi la démarche professionnelle du bureau d'études ou du service administratif.

\courssub{La confrontation des productions}

\begin{definition}[Confrontation des productions]
La \cle{confrontation} est une phase d'échange croisé : chaque binôme remet sa requête à un autre binôme, qui l'annote à l'aide d'une \cle{grille critériée} (liste de critères de réussite). On ne juge pas l'auteur, on évalue le texte : chaque remarque doit être justifiée par un critère précis.
\end{definition}

Pourquoi confronter ? Parce que l'auteur d'un texte est presque toujours aveugle à ses propres erreurs : il relit ce qu'il a voulu écrire, pas ce qu'il a écrit. Le lecteur extérieur, lui, voit immédiatement une formule de politesse manquante ou un objet ambigu. La confrontation reproduit en classe la situation réelle du destinataire qui découvre la requête pour la première fois.

\begin{methode}[Conduire une confrontation croisée en binômes]
\begin{enumerate}
\item \textbf{Échanger} les requêtes entre binômes, sans commentaire oral préalable : le texte doit se défendre seul.
\item \textbf{Lire une première fois} sans rien écrire, pour saisir la demande globale.
\item \textbf{Annoter} avec la grille critériée : cocher chaque critère, noter en marge les passages fautifs avec un code (S = structure, L = langue, X = lexique, P = présentation).
\item \textbf{Rendre} la copie annotée et formuler oralement deux points forts et deux points à améliorer.
\item \textbf{Noter} les remarques reçues sans discuter : c'est l'auteur qui décidera des corrections à la phase suivante.
\end{enumerate}
\end{methode}

\begin{exemplebox}[Grille critériée de confrontation]
\begin{center}
\begin{tabularx}{\linewidth}{|l|X|c|}
\hline
\rowcolor{popblueL} \textbf{Domaine} & \textbf{Critère} & \textbf{Oui/Non} \\
\hline
Structure externe & Lieu et date, objet, formule d'appel (« Monsieur le Maire »), formule de politesse finale, signature & \\
\hline
Structure interne & Introduction (identité du requérant), corps (exposé des faits, demande claire), conclusion (remerciements) & \\
\hline
Langue & Orthographe, grammaire, ponctuation, phrases complètes & \\
\hline
Lexique & Registre soutenu, lexique administratif approprié, aucune familiarité & \\
\hline
Présentation & Aération, alinéas, lisibilité, paragraphes équilibrés & \\
\hline
\end{tabularx}
\end{center}
\end{exemplebox}

\courssub{L'amélioration guidée et la réécriture}

Une fois la copie annotée récupérée, commence le vrai travail : l'amélioration. L'erreur serait de corriger au hasard, faute par faute. Il faut procéder par relectures successives, chacune consacrée à un seul aspect.

\begin{methode}[Améliorer une production : la méthode des trois relectures]
\begin{enumerate}
\item \textbf{Première relecture — la structure} : vérifier la présence et la place de chaque élément de la structure externe (en-tête, objet, appel, politesse, signature) et interne (introduction, corps, conclusion). Ajouter ou déplacer ce qui manque.
\item \textbf{Deuxième relecture — la langue et le lexique} : corriger l'orthographe, les accords, la ponctuation ; remplacer les mots familiers par le lexique administratif (« je veux » devient « je sollicite », « donner » devient « accorder »).
\item \textbf{Troisième relecture — la présentation} : aérer, créer des alinéas, équilibrer les paragraphes, vérifier la lisibilité de la signature.
\item \textbf{Réécrire intégralement} la version finale au propre : on ne se contente jamais de ratures corrigées.
\end{enumerate}
\end{methode}

\begin{exoresolu}[Corriger un extrait de requête fautive]
Énoncé. Voici un extrait de requête rédigé par un élève. Relevez les 5 erreurs ciblées et proposez une version corrigée.

\emph{« Douala, le 12 mars 2026. Objet : demande de stage. Salut Monsieur le Directeur, je me nomme MBARGA Paul, élève en Terminale F4. Je veux faire un stage dans ton entreprise parce que j'ai besoin d'experience. Merci d'avance. »}
\tcblower
Solution détaillée. Les 5 erreurs :
\begin{enumerate}
\item \textbf{Formule d'appel familière} : « Salut Monsieur le Directeur » est inacceptable ; on écrit simplement « Monsieur le Directeur, ».
\item \textbf{Registre familier} : « je veux » doit devenir « je sollicite » ou « j'ai l'honneur de solliciter ».
\item \textbf{Tutoiement interdit} : « ton entreprise » devient « votre entreprise ».
\item \textbf{Orthographe} : « experience » prend un accent aigu : « expérience ».
\item \textbf{Conclusion absente} : « Merci d'avance » est trop sec ; il faut une formule de politesse développée.
\end{enumerate}
Version corrigée :

\emph{« Douala, le 12 mars 2026. Objet : demande de stage. Monsieur le Directeur, je me nomme MBARGA Paul, élève en classe de Terminale F4. J'ai l'honneur de solliciter un stage au sein de votre entreprise afin d'acquérir une expérience professionnelle. Dans l'attente d'une suite favorable, je vous prie d'agréer, Monsieur le Directeur, l'expression de ma haute considération. »}
\end{exoresolu}

\begin{attention}[Corriger l'orthographe ne suffit pas]
L'erreur la plus fréquente consiste à corriger uniquement l'orthographe et la grammaire, en oubliant de vérifier la \textbf{conformité de la structure externe} (en-tête, objet, appel, formule de politesse, signature). Or, au Probatoire et au Baccalauréat technique, la structure est fortement pondérée dans le barème : une requête sans faute mais sans objet ni formule de politesse perd la moitié des points. Toujours commencer la relecture par la structure.
\end{attention}

\begin{popfigure}
\begin{tikzpicture}[scale=1, every node/.style={font=\small}]
\node[draw=popblue, fill=popblueL, rounded corners, minimum width=3.2cm, minimum height=1cm] (a) at (0,0) {\textbf{1. Production}};
\node[draw=poporange, fill=poporangeL, rounded corners, minimum width=3.2cm, minimum height=1cm] (b) at (4.5,0) {\textbf{2. Confrontation}};
\node[draw=poppurple, fill=poppurpleL, rounded corners, minimum width=3.2cm, minimum height=1cm] (c) at (4.5,-2.5) {\textbf{3. Correction}};
\node[draw=popgreen, fill=popgreenL, rounded corners, minimum width=3.2cm, minimum height=1cm] (d) at (0,-2.5) {\textbf{4. Version finale}};
\draw[-{Stealth}, thick, popdark] (a) -- (b) node[midway, above]{échange croisé};
\draw[-{Stealth}, thick, popdark] (b) -- (c) node[midway, right]{annotations};
\draw[-{Stealth}, thick, popdark] (c) -- (d) node[midway, below]{3 relectures};
\draw[-{Stealth}, thick, popdark, dashed] (d) -- (a) node[midway, left]{nouvelle production};
\end{tikzpicture}
\legende{Le cycle d'amélioration d'une production écrite : chaque nouvelle production bénéficie des corrections de la précédente.}
\end{popfigure}

\courssub{Extension au rapport}

La même démarche s'applique au \cle{rapport}, autre texte utilitaire du programme. Les binômes échangent leurs rapports, les annotent avec une grille adaptée (page de garde, introduction, développement ordonné, conclusion, langue), puis réécrivent la version finale. Les critères changent — un rapport expose des faits constatés, il ne demande rien — mais la méthode reste identique : confrontation, trois relectures, réécriture intégrale. Cette permanence de la démarche montre qu'elle constitue une compétence transférable à tout texte professionnel.

\courssub{Le compte rendu de la leçon}

Le \cle{compte rendu} est une synthèse collective, construite au tableau avec l'enseignant, qui fixe les acquis de toute la leçon. Chaque élève le recopie et le mémorise : c'est l'outil de révision avant l'évaluation.

\begin{aretenir}[Synthèse des acquis de la leçon]
\begin{itemize}
\item \textbf{La requête} est un texte fonctionnel par lequel un citoyen sollicite une autorité ; sa structure externe comprend lieu et date, objet, formule d'appel, formule de politesse et signature ; sa structure interne comprend introduction, corps (faits + demande) et conclusion.
\item \textbf{Les actes de langage} : tout énoncé comporte un acte locutoire (ce qu'on dit), illocutoire (ce qu'on fait en le disant : demander, remercier) et perlocutoire (l'effet obtenu sur le destinataire).
\item \textbf{Le lexique spécialisé} administratif (« solliciter », « agréer », « suite favorable ») garantit le registre soutenu exigé.
\item \textbf{L'emprunt} est un mot venant d'une langue étrangère (ex. « stage », « planning ») ; dans une requête, on préfère l'équivalent français quand il existe.
\item \textbf{L'amélioration} passe par la confrontation croisée, trois relectures ciblées et une réécriture intégrale.
\end{itemize}
\end{aretenir}

\courssub{La remédiation ciblée}

La remédiation consiste à traiter immédiatement, par des exercices différenciés, les difficultés relevées pendant la confrontation. Le diagramme ci-dessous présente la répartition typique des erreurs dans une classe de Terminale technique.

\begin{popfigure}
\begin{tikzpicture}[scale=1]
% Langue 40% = 144 deg ; Structure 30% = 108 deg ; Lexique 20% = 72 deg ; Présentation 10% = 36 deg
\fill[popblue] (0,0) -- (90:2.4) arc (90:234:2.4) -- cycle;
\fill[poporange] (0,0) -- (234:2.4) arc (234:342:2.4) -- cycle;
\fill[poppurple] (0,0) -- (342:2.4) arc (342:54:2.4) -- cycle;
\fill[popgreen] (0,0) -- (54:2.4) arc (54:90:2.4) -- cycle;
\draw[popdark, thick] (0,0) circle (2.4);
\node[font=\small\bfseries, white] at (162:1.5) {40 \%};
\node[font=\small\bfseries, white] at (288:1.5) {30 \%};
\node[font=\small\bfseries, white] at (18:1.5) {20 \%};
\node[font=\small\bfseries, white] at (72:1.5) {10 \%};
% Légende
\node[draw=popblue, fill=popblue, minimum width=0.35cm, minimum height=0.35cm] at (3.4,1.5) {};
\node[right, font=\small] at (3.7,1.5) {Langue (orthographe, grammaire) : 40 \%};
\node[draw=poporange, fill=poporange, minimum width=0.35cm, minimum height=0.35cm] at (3.4,0.7) {};
\node[right, font=\small] at (3.7,0.7) {Structure externe et interne : 30 \%};
\node[draw=poppurple, fill=poppurple, minimum width=0.35cm, minimum height=0.35cm] at (3.4,-0.1) {};
\node[right, font=\small] at (3.7,-0.1) {Lexique administratif : 20 \%};
\node[draw=popgreen, fill=popgreen, minimum width=0.35cm, minimum height=0.35cm] at (3.4,-0.9) {};
\node[right, font=\small] at (3.7,-0.9) {Présentation : 10 \%};
\end{tikzpicture}
\legende{Répartition des erreurs relevées dans les productions de la classe : langue 40 \%, structure 30 \%, lexique 20 \%, présentation 10 \% (total : 100 \%). La remédiation porte d'abord sur les domaines les plus fautifs.}
\end{popfigure}

\begin{exemplebox}[Exercices différenciés de remédiation]
\begin{itemize}
\item \textbf{Difficulté : formules de politesse.} Compléter : « Je vous prie d'\ldots{}, Monsieur le Préfet, l'expression de ma \ldots{} considération. » Puis réécrire trois formules adaptées à trois destinataires différents (maire, directeur, ministre).
\item \textbf{Difficulté : orthographe administrative.} Dictée de dix mots du lexique administratif : « solliciter », « agréer », « honorabilité », « pièces jointes », « expérience », « autorisation ».
\item \textbf{Difficulté : plan.} À partir d'une situation donnée (demande d'acte de naissance), remettre dans l'ordre les cinq éléments de la structure externe, puis rédiger uniquement l'introduction et la conclusion.
\end{itemize}
\end{exemplebox}

\begin{savaistu}[L'efficacité de la remédiation immédiate]
Les recherches en pédagogie par objectifs montrent qu'une remédiation menée \emph{immédiatement} après la confrontation des productions améliore nettement les performances à la production suivante. Corriger à chaud, quand l'erreur est encore présente à l'esprit, vaut mieux qu'une correction différée de plusieurs semaines : c'est pourquoi cette séance précède directement l'évaluation sommative.
\end{savaistu}

\courssub{Évaluation sommative}

La leçon se clôt par une évaluation individuelle, notée sur 20, en situation nouvelle : aucune aide, aucun échange, comme à l'examen.

\begin{exercice}[Évaluation sommative : rédaction d'une requête complète (20 points)]
Situation : vous êtes élève en Terminale technique au Lycée technique de Bafoussam. Vous avez perdu votre relevé de notes du Probatoire, indispensable à votre inscription à un concours professionnel. Rédigez une requête complète adressée à Monsieur le Proviseur pour solliciter un duplicata.

Barème indicatif : structure externe (6 pts), structure interne et clarté de la demande (5 pts), langue et ponctuation (4 pts), lexique administratif et registre soutenu (3 pts), présentation (2 pts). Durée : 45 minutes.
\end{exercice}

\begin{corrige}[Évaluation sommative]
Attendus essentiels : lieu et date en haut à droite ; objet précis (« Demande de duplicata du relevé de notes du Probatoire ») ; formule d'appel « Monsieur le Proviseur, » ; introduction présentant l'identité du requérant (nom, classe) ; corps exposant les faits (perte du document, circonstances, urgence) et formulant clairement la demande ; conclusion avec remerciements ; formule de politesse développée (« Je vous prie d'agréer, Monsieur le Proviseur, l'expression de ma respectueuse considération. ») ; signature. Toute familiarité, tout tutoiement ou toute absence d'élément structurel entraîne une perte de points selon le barème.
\end{corrige}

Vous disposez désormais de la démarche complète de production d'un texte utilitaire : analyser la situation, planifier, rédiger, confronter, corriger et réécrire. Cette compétence, évaluée au Baccalauréat technique, vous servira dans toute votre vie professionnelle et administrative.

\newpage

\coursec{Activité d'intégration}

\courssub{Objectifs de l'activité}

À l'issue de cette activité d'intégration, l'élève doit être capable de :

\begin{itemize}
\item analyser une situation de communication professionnelle à l'aide de la grille \cle{QQOQCP} (Qui ? Quoi ? Où ? Quand ? Comment ? Pourquoi ?) ;
\item exploiter des textes fonctionnels (bon de livraison, facture) et un texte iconique (organigramme) pour identifier le destinataire compétent ;
\item reconnaître et produire les \cle{actes de langage} adaptés : actes locutoire, illocutoire (informer, réclamer, solliciter) et perlocutoire attendu ;
\item mobiliser un \cle{lexique spécialisé} (commercial et administratif) correct et éviter les \cle{emprunts} fautifs ;
\item rédiger une \cle{requête} complète respectant la structure externe (entête, références, objet, formules de politesse) et la structure interne (introduction, corps, conclusion) ;
\item évaluer et améliorer une production à l'aide d'une grille critériée, puis justifier oralement ses choix.
\end{itemize}

\courssub{Situation}

L'atelier de maintenance informatique \og Nkoulou Tech \fg{}, situé au quartier Mbanga à Ngaoundéré (région de l'Adamaoua), a livré, le \textbf{15 janvier 2025}, dix ordinateurs de bureau au lycée technique de Banyo, dans le cadre d'un marché passé avec la Délégation régionale des Enseignements secondaires de l'Adamaoua. Le montant total de la facture n° 2025-014 s'élève à \textbf{4 500 000 FCFA}, payable sous trente jours.

Six mois plus tard, le \textbf{18 juillet 2025}, la facture demeure impayée, malgré deux relances orales effectuées auprès du comptable de la Délégation les 20 mars et 12 mai 2025. Le gérant de Nkoulou Tech, M. Nkoulou Ibrahim, décide d'adresser une requête écrite à l'autorité compétente.

\textbf{Dossier documentaire remis aux élèves :}

\textbf{Document 1 — Extrait du bon de livraison n° BL-2025-014 (texte fonctionnel) :} \og Reçu ce jour 15/01/2025, dix (10) ordinateurs de marque HP, complets, en bon état de fonctionnement. Signature du receveur : le Proviseur du lycée technique de Banyo. \fg{}

\textbf{Document 2 — Organigramme simplifié de la Délégation régionale (texte iconique) :}

\begin{popfigure}
\begin{center}
\begin{tikzpicture}[node distance=1.1cm,
box/.style={draw=popblue, fill=popblueL, rounded corners, text width=4.6cm, align=center, minimum height=0.9cm, font=\small},
fleche/.style={-{Stealth}, thick, popdark}]
\node[box, fill=popgoldL, draw=popgold] (del) {\textbf{Délégué régional} des Enseignements secondaires};
\node[box, below left=1.1cm and 0.2cm of del] (sg) {Secrétariat général};
\node[box, below right=1.1cm and 0.2cm of del] (afc) {Agence comptable (comptable des matières)};
\node[box, below=2.6cm of del] (prov) {Proviseurs des lycées};
\draw[fleche] (del) -- (sg);
\draw[fleche] (del) -- (afc);
\draw[fleche] (del) -- (prov);
\end{tikzpicture}
\end{center}
\legende{Organigramme simplifié : le Délégué régional est l'autorité de tutelle ; l'agence comptable exécute les paiements.}
\end{popfigure}

\textbf{Document 3 — Extrait de la facture (données chiffrées) :}

\begin{center}
\begin{tabularx}{\linewidth}{|l|X|r|}\hline
\rowcolor{popblueL} \textbf{Désignation} & \textbf{Détail} & \textbf{Montant (FCFA)} \\ \hline
Ordinateurs HP & 10 unités à 400 000 & 4 000 000 \\ \hline
Installation et mise en service & Forfait atelier & 500 000 \\ \hline
\textbf{Total TTC} & Facture n° 2025-014 du 15/01/2025 & \textbf{4 500 000} \\ \hline
\end{tabularx}
\end{center}

\textbf{Document 4 — Extraits de courriers mal rédigés (à corriger) :}

\og Cher monsieur, ça fait longtemps que vous ne m'avez pas payé mon cash. Je suis vraiment disappointed. Faites vite sinon je vais tout déballer sur les réseaux sociaux. \fg{} (Courrier A)

\og Monsieur, je vous demande de bien vouloir me payer. Merci d'avance. \fg{} (Courrier B, sans entête, sans objet, sans références, sans date.)

\courssub{Consignes}

Par binômes, et à l'aide du dossier documentaire, vous allez :

\begin{enumerate}
\item \textbf{Analyser la situation} en remplissant la grille QQOQCP : qui écrit ? à qui ? de quoi s'agit-il ? où et quand les faits se sont-ils déroulés ? comment (circonstances, relances) ? pourquoi écrire (but poursuivi) ?
\item \textbf{Identifier le destinataire compétent} sur l'organigramme (document 2) et justifier ce choix : pourquoi ne pas écrire directement au comptable ou au proviseur ?
\item \textbf{Relever les actes illocutoires} nécessaires à cette requête (informer des faits, réclamer le paiement, solliciter une suite favorable) et préciser l'effet perlocutoire recherché sur le destinataire.
\item \textbf{Corriger les courriers A et B} : relevez au moins cinq fautes (registre de langue, emprunts fautifs, menace inappropriée, structure externe absente) et proposez des reformulations correctes.
\item \textbf{Élaborer le plan détaillé} de la requête : entête, références, objet, introduction (rappel des faits), corps (exposé du préjudice et demande argumentée), conclusion (formule de politesse administrative).
\item \textbf{Rédiger la requête complète} (une page maximum) en employant au moins huit termes du lexique commercial et administratif (facture, bon de livraison, règlement, créance, délai, relance, ordonnancement, suite favorable...) et sans aucun emprunt fautif.
\item \textbf{Échanger votre production} avec un autre binôme et l'évaluer à l'aide de la grille critériée sur 20 points (structure externe : 4 pts ; structure interne : 4 pts ; exactitude des faits et des chiffres : 4 pts ; lexique spécialisé et correction de la langue : 4 pts ; formules et registre soutenu : 4 pts). Améliorez ensuite votre texte.
\item \textbf{Justifier oralement}, devant la classe, deux de vos choix : la formule d'appel, l'organisation du corps et la formule de politesse finale.
\end{enumerate}

\courssub{Ressources à mobiliser}

\begin{aretenir}[Outils de la leçon à réactiver]
\begin{itemize}
\item \textbf{Actes de langage} : l'acte \cle{locutoire} (ce qu'on dit), l'acte \cle{illocutoire} (la force de ce qu'on dit : informer, réclamer, solliciter), l'acte \cle{perlocutoire} (l'effet produit : obtenir le paiement).
\item \textbf{Structure externe de la requête} : coordonnées de l'expéditeur, lieu et date, coordonnées du destinataire, références, objet, formule d'appel (\og Monsieur le Délégué régional \fg{}), signature.
\item \textbf{Structure interne} : introduction (rappel précis et daté des faits), corps (exposé du préjudice, demande claire et argumentée), conclusion (formule de politesse : \og Je vous prie d'agréer, Monsieur le Délégué régional, l'expression de ma haute considération \fg{}).
\item \textbf{Lexique spécialisé} : employer les termes exacts du domaine commercial et administratif ; \textbf{emprunts} : proscrire les mots anglais non intégrés (\emph{cash}, \emph{disappointed}) et leur préférer \og argent liquide / paiement \fg{}, \og déçu \fg{}.
\item \textbf{Grille QQOQCP} pour l'analyse de toute situation de communication.
\end{itemize}
\end{aretenir}

\begin{attention}[Pièges à éviter]
Ne confondez pas requête et réclamation agressive : la requête reste courtoise et factuelle. Ne menacez jamais le destinataire. Vérifiez que tous les éléments de l'entête sont présents et que les chiffres cités (4 500 000 FCFA, dates, numéros de pièces) correspondent exactement aux documents.
\end{attention}

\courssub{Démarche et corrigé commenté}

\begin{exoresolu}[La requête de Nkoulou Tech — résolution guidée]
Mobiliser l'ensemble des savoir-faire de la leçon pour produire la requête complète.

\tcblower

\textbf{Étape 1 — Grille QQOQCP.} \emph{Qui ?} M. Nkoulou Ibrahim, gérant de l'atelier Nkoulou Tech, créancier. \emph{À qui ?} Le Délégué régional des Enseignements secondaires de l'Adamaoua, autorité compétente. \emph{Quoi ?} Une facture de 4 500 000 FCFA impayée. \emph{Où ?} Ngaoundéré / lycée technique de Banyo. \emph{Quand ?} Livraison le 15/01/2025 ; relances les 20/03 et 12/05/2025 ; requête le 18/07/2025. \emph{Comment ?} Livraison conforme attestée par le bon BL-2025-014 signé ; deux relances orales restées sans suite. \emph{Pourquoi ?} Obtenir le règlement de la créance.

\textbf{Étape 2 — Destinataire.} Sur l'organigramme, le Délégué régional est l'autorité de tutelle : il ordonnance les dépenses. Le comptable ne fait qu'exécuter le paiement ; le proviseur n'est que receveur du matériel. On écrit donc au Délégué régional.

\textbf{Étape 3 — Actes de langage.} Acte locutoire : le texte de la lettre. Actes illocutoires : informer (rappel des faits), réclamer (le paiement dû), solliciter (une suite rapide). Acte perlocutoire visé : le déclenchement du règlement.

\textbf{Étape 4 — Correction des courriers.} Courrier A : familiarité et registre relâché (\og Cher monsieur, ça fait \fg{}), emprunts fautifs (\emph{cash}, \emph{disappointed}), menace inacceptable (\og déballer sur les réseaux sociaux \fg{}), absence de faits précis. Reformulation : \og Monsieur le Délégué régional, la facture n° 2025-014 demeure impayée six mois après la livraison ; je sollicite son règlement dans les meilleurs délais. \fg{} Courrier B : il faut ajouter entête, date, références, objet, faits, montant et formule de politesse.

\textbf{Étape 5 — Plan détaillé.} 1. Entête complet. 2. Objet : demande de règlement de la facture n° 2025-014. 3. Introduction : rappel du marché et de la livraison conforme. 4. Corps : constat de l'impayé, relances infructueuses, préjudice pour l'entreprise, demande argumentée. 5. Conclusion : remerciement et formule de politesse.

\textbf{Étape 6 — Requête rédigée (modèle).}

\emph{Nkoulou Tech, Atelier de maintenance informatique, quartier Mbanga, BP 312 Ngaoundéré. — Ngaoundéré, le 18 juillet 2025. — À Monsieur le Délégué régional des Enseignements secondaires de l'Adamaoua, Ngaoundéré. — Réf. : facture n° 2025-014 ; bon de livraison BL-2025-014. — Objet : demande de règlement de la facture n° 2025-014 d'un montant de 4 500 000 FCFA.}

\og Monsieur le Délégué régional,

J'ai l'honneur de vous rappeler que mon entreprise a livré, le 15 janvier 2025, dix ordinateurs au lycée technique de Banyo, conformément au marché conclu avec votre délégation. La réception du matériel a été attestée par le bon de livraison BL-2025-014, dûment signé.

Cependant, la facture n° 2025-014, d'un montant de quatre millions cinq cent mille francs CFA, demeure impayée à ce jour, soit six mois après la livraison, et ce malgré deux relances effectuées auprès de votre agence comptable les 20 mars et 12 mai 2025. Cette créance non réglée fragilise la trésorerie d'une petite entreprise qui emploie six techniciens.

C'est pourquoi je sollicite respectueusement votre bienveillante intervention afin que soit ordonnancé le règlement de cette facture dans les meilleurs délais.

Dans l'espoir d'une suite favorable, je vous prie d'agréer, Monsieur le Délégué régional, l'expression de ma haute considération. \fg{}

\emph{Signature : Nkoulou Ibrahim, gérant. Pièces jointes : copie de la facture, copie du bon de livraison.}

\textbf{Étapes 7 et 8 — Évaluation croisée et justification orale.} La grille sur 20 points permet de vérifier chaque critère ; à l'oral, l'élève justifie par exemple le choix de la formule d'appel (adaptée au rang du destinataire) et l'ordre chronologique du corps (faits d'abord, demande ensuite), ce qui garantit clarté et courtoisie.
\end{exoresolu}

\courssub{Bilan de l'activité}

Cette activité a permis de mettre en évidence que la requête est un texte utilitaire complet, qui exige bien plus qu'une simple demande : une \textbf{analyse préalable de la situation} (QQOQCP), une \textbf{lecture croisée de documents fonctionnels et iconiques} pour viser le bon destinataire, la maîtrise des \textbf{actes illocutoires} pour doser information, réclamation et sollicitation, et un \textbf{lexique spécialisé} exempt d'emprunts fautifs. L'évaluation croisée et la justification orale ont enfin montré qu'un écrit professionnel efficace est un écrit structuré, exact dans ses chiffres, courtois dans son ton — trois qualités décisives dans la vie professionnelle d'un technicien camerounais.

\newpage

\coursec{Méthodes et exercices résolus}

\courssub{Les textes fonctionnels et les actes de langage}

\begin{methode}[Identifier les actes de langage dans un texte fonctionnel]
\textbf{Objectif :} Reconnaître la dimension performative du langage dans une requête (demander, supplier, exiger, remercier d'avance).
\begin{enumerate}
    \item \textbf{Repérer le verbe performatif} (souvent à la 1\textsuperscript{re} personne du présent de l'indicatif : \emph{je demande, je sollicite, je prie}) ou la structure modalisée (\emph{je voudrais que, il serait souhaitable que}).
    \item \textbf{Distinguier les trois niveaux} (Austin/Searle) :
    \begin{itemize}
        \item \textbf{Locutoire :} Ce qui est dit (phrases, sens littéral). Ex: \emph{« Je vous écris pour demander un stage. »}
        \item \textbf{Illocutoire :} La force de l'énoncé (l'acte posé). Ex: \emph{Demander / Solliciter} (directive).
        \item \textbf{Perlocutoire :} L'effet visé sur le destinataire. Ex: \emph{Obtenir une réponse favorable / un rendez-vous}.
    \end{itemize}
    \item \textbf{Qualifier l'acte} selon la classification de Searle : \textbf{Directives} (demander, ordonner), \textbf{Expressifs} (remercier, s'excuser), \textbf{Engagements} (promettre), \textbf{Assertifs} (affirmer), \textbf{Déclaratifs} (nommer).
    \item \textbf{Justifier} par des indices linguistiques : modes (impératif, subjonctif), modaux (devoir, pouvoir), politesse (conditionnel, atténuateurs).
\end{enumerate}
\cle{Acte illocutoire} \cle{Force performative} \cle{Verbe performatif}
\end{methode}

\begin{exoresolu}[Analyse d'extraits de requête : actes de langage]
\textbf{Énoncé :} Lisez les extraits suivants d'une lettre de demande de bourse adressée au Ministre de l'Enseignement Supérieur. Pour chaque extrait, identifiez l'acte locutoire, l'acte illocutoire (type Searle) et l'effet perlocutoire visé.

\textbf{Extrait 1 :} \emph{« Monsieur le Ministre, j'ai l'honneur de solliciter votre bienveillance pour l'obtention d'une bourse d'études pour l'année académique 2025-2026. »}

\textbf{Extrait 2 :} \emph{« Je joins à ce courrier l'attestation de réussite au Baccalauréat, le relevé de notes et un certificat de visite médicale. »}

\textbf{Extrait 3 :} \emph{« En espérant une suite favorable à ma doléance, je vous prie d'agréer, Monsieur le Ministre, l'expression de ma très haute considération. »}

\tcblower
\textbf{Solution détaillée :}

\textbf{Extrait 1 :}
\begin{itemize}
    \item \textbf{Locutoire :} Énoncé déclaratif informant le destinataire de l'intention de l'expéditeur (solliciter une bourse).
    \item \textbf{Illocutoire :} \textbf{Directive} (acte de \textbf{solliciter / demander}). Le verbe \emph{solliciter} est un verbe performatif explicite (1\textsuperscript{re} pers. présent). La formule \emph{j'ai l'honneur de} renforce la dimension rituelle et respectueuse.
    \item \textbf{Perlocutoire :} Susciter la bienveillance du Ministre ; obtenir l'octroi de la bourse.
\end{itemize}

\textbf{Extrait 2 :}
\begin{itemize}
    \item \textbf{Locutoire :} Énoncé assertif (constatatif) énumérant des pièces jointes.
    \item \textbf{Illocutoire :} \textbf{Assertif} (acte d'\textbf{informer / attester}). Le verbe \emph{joins} constate une action matérielle qui prouve la recevabilité du dossier.
    \item \textbf{Perlocutoire :} Rassurer le destinataire sur la complétude du dossier ; prouver la sérieux du candidat (effet de crédibilité).
\end{itemize}

\textbf{Extrait 3 :}
\begin{itemize}
    \item \textbf{Locutoire :} Phrase complexe : subordonnée (\emph{En espérant...}) + principale impérative polie (\emph{je vous prie d'agréer...}).
    \item \textbf{Illocutoire :} \textbf{Expressif} (acte de \textbf{remercier / saluer} anticipé) + \textbf{Directive} (acte de \textbf{prier} pour l'agrément des salutations). \emph{En espérant} exprime un souhait (expressif) qui fonctionne comme une relance directive implicite.
    \item \textbf{Perlocutoire :} Laisser une impression finale positive (politesse de clôture) ; inciter à traiter le dossier rapidement (pression douce via \emph{en espérant une suite favorable}).
\end{itemize}

\textbf{Conclusion :} Une requête efficace enchaîne des actes \textbf{directifs} (le cœur de la demande), étayés par des \textbf{assertifs} (les preuves), et encadrés d'\textbf{expressifs} (politesse, espoir) pour maximiser l'effet perlocutoire (obtention).
\end{exoresolu}

\courssub{La structure externe et interne de la requête}

\begin{methode}[Analyser et respecter la structure canonique de la requête administrative]
\textbf{Objectif :} Produire un document conforme aux normes administratives camerounaises (format A4, marges, plis).
\begin{enumerate}
    \item \textbf{Structure externe (Mise en page) :}
    \begin{itemize}
        \item \textbf{En-tête expéditeur} (haut gauche) : Nom, Prénom, Adresse, Téléphone, Email, Matricule/Spécialité.
        \item \textbf{En-tête destinataire} (haut droite, sous expéditeur) : Qualité, Nom/Service, Adresse.
        \item \textbf{Lieu et Date} (ligne séparée, alignée à droite généralement).
        \item \textbf{Objet} (gras, centré ou gauche, concis : \emph{Objet : Demande de stage / Bourse / Document}).
        \item \textbf{Références} (si réponse à un courrier : \emph{Réf : Votre lettre n°... du...}).
        \item \textbf{Corps de la lettre} : Formules d'appel $\rightarrow$ Introduction $\rightarrow$ Exposition des motifs (Corps) $\rightarrow$ Conclusion (La demande formelle) $\rightarrow$ Formules de politesse (salutation).
        \item \textbf{Bas de page} : Signature (manuscrite) + Nom prénom tapé $\rightarrow$ Mentions \emph{Pièces jointes :} (liste numérotée) $\rightarrow$ \emph{Copie à :} (si nécessaire).
    \end{itemize}
    \item \textbf{Structure interne (Logique argumentative) :}
    \begin{itemize}
        \item \textbf{Introduction (Exorde) :} Qui suis-je ? (Identité, qualité d'étudiant). Pourquoi j'écris ? (Objet rappelé). Sur quel fondement ? (Texte de loi, arrêté, coutume).
        \item \textbf{Corps (Exposition / Argumentation) :} Les motifs (situation familiale, mérite scolaire, nécessité professionnelle). Arguments ordonnés (chronologique, importance croissante). Preuves (pièces jointes citées).
        \item \textbf{Conclusion (Pétition) :} La demande explicite, claire, unique (\emph{« Je vous prie de bien vouloir m'accorder... »}). Perspective future (disponibilité pour entretien).
    \end{itemize}
\end{enumerate}
\cle{Structure externe} \cle{Structure interne} \cle{Pétition} \cle{Exorde}
\end{methode}

\begin{exoresolu}[Restructuration d'une requête désordonnée]
\textbf{Énoncé :} Voici les éléments d'une demande de stage de fin d'études rédigés dans le désordre. Réorganisez-les selon la structure externe et interne canonique. Indiquez pour chaque élément sa zone (En-tête, Objet, Introduction, Corps, Conclusion, Salutation, Bas de page).

\begin{enumerate}
    \item \emph{« Je suis élève en Terminale TEE au Lycée Technique de Douala-Bassa. »}
    \item \emph{« Pièces jointes : 1. CV ; 2. Lettre de motivation ; 3. Attestation de scolarité. »}
    \item \emph{« Douala, le 15 mars 2025 »}
    \item \emph{« À l'attention du Directeur Général de l'ENEO, B.P. 1837 Douala. »}
    \item \emph{« Je vous prie d'agréer, Monsieur le Directeur Général, l'assurance de ma considération distinguée. »}
    \item \emph{« Objet : Demande de stage de fin d'études (Juillet-Août 2025). »}
    \item \emph{« Nguimegni Armand, B.P. 456 Douala, Tel : 6XX XX XX XX. »}
    \item \emph{« Ce stage me permettrait de mettre en pratique mes connaissances en électrotechnique et de découvrir le milieu professionnel. »}
    \item \emph{« En vous remerciant par avance de l'attention portée à ma candidature, je reste à votre disposition pour tout entretien. »}
    \item \emph{« Je sollicite votre bienveillance pour effectuer un stage de deux mois au sein de votre direction régionale. »}
\end{enumerate}

\tcblower
\textbf{Solution détaillée :}

\textbf{1. EN-TÊTE (Externe)}
\begin{itemize}
    \item \textbf{Expéditeur (Haut gauche) :} \textbf{(7)} \emph{Nguimegni Armand, B.P. 456 Douala, Tel : 6XX XX XX XX.}
    \item \textbf{Destinataire (Haut droite) :} \textbf{(4)} \emph{À l'attention du Directeur Général de l'ENEO, B.P. 1837 Douala.}
    \item \textbf{Lieu et Date :} \textbf{(3)} \emph{Douala, le 15 mars 2025}
\end{itemize}

\textbf{2. OBJET}
\begin{itemize}
    \item \textbf{(6)} \emph{Objet : Demande de stage de fin d'études (Juillet-Août 2025).}
\end{itemize}

\textbf{3. CORPS DE LA LETTRE (Interne)}
\begin{itemize}
    \item \textbf{Formule d'appel (implicite standard) :} \emph{Monsieur le Directeur Général,}
    \item \textbf{INTRODUCTION (Exorde) :} \textbf{(1)} \emph{Je suis élève en Terminale TEE au Lycée Technique de Douala-Bassa.} (Identité + Qualité).
    \item \textbf{CORPS (Exposition des motifs) :}
    \begin{enumerate}
        \item \textbf{(10)} \emph{Je sollicite votre bienveillance pour effectuer un stage de deux mois au sein de votre direction régionale.} (La demande brute / Pétition anticipée souvent placée ici ou en conclusion).
        \item \textbf{(8)} \emph{Ce stage me permettrait de mettre en pratique mes connaissances en électrotechnique et de découvrir le milieu professionnel.} (Argumentation : utilité pédagogique + professionnelle).
    \end{enumerate}
    \item \textbf{CONCLUSION (Pétition formelle + Disponibilité) :} \textbf{(9)} \emph{En vous remerciant par avance de l'attention portée à ma candidature, je reste à votre disposition pour tout entretien.}
\end{itemize}

\textbf{4. FORMULE DE POLITESSE (Salutation)}
\begin{itemize}
    \item \textbf{(5)} \emph{Je vous prie d'agréer, Monsieur le Directeur Général, l'assurance de ma considération distinguée.}
\end{itemize}

\textbf{5. BAS DE PAGE}
\begin{itemize}
    \item \textbf{Signature (espace pour manuscrite) + Nom tapé :} \emph{Nguimegni Armand}
    \item \textbf{Pièces jointes :} \textbf{(2)} \emph{Pièces jointes : 1. CV ; 2. Lettre de motivation ; 3. Attestation de scolarité.}
\end{itemize}

\textbf{Conclusion :} La structure externe assure la \textbf{traçabilité administrative} (qui, où, quand, quoi). La structure interne assure la \textbf{persuasion} (identité $\rightarrow$ motifs $\rightarrow$ demande claire).
\end{exoresolu}

\courssub{Le lexique spécialisé et l'emprunt}

\begin{methode}[Mobiliser le lexique spécialisé et gérer l'emprunt dans la requête]
\textbf{Objectif :} Adapter le registre de langue (formel, administratif, technique) et intégrer les termes étrangers (emprunts) sans faute.
\begin{enumerate}
    \item \textbf{Lexique spécialisé (Administratif / Technique) :}
    \begin{itemize}
        \item \textbf{Repérer le champ lexical} de la situation (éducation, entreprise, santé, justice).
        \item \textbf{Substituer le terme courant par le terme précis :}
        \begin{center}
            \begin{tabular}{|l|l|}
            \hline
            \textbf{Courant} & \textbf{Spécialisé (Administratif/Technique)} \\
            \hline
            Demander & \textbf{Solliciter, requérir, postuler} \\
            Donner & \textbf{Octroyer, accorder, délivrer} \\
            Regarder / Étudier & \textbf{Examiner, instruire, apprécier} \\
            Réponse & \textbf{Suite, issue, décision} \\
            Papiers & \textbf{Pièces justificatives, dossier, attestations} \\
            Stage & \textbf{Stage d'immersion, stage de perfectionnement, PFE} \\
            \hline
            \end{tabular}
        \end{center}
    \end{itemize}
    \item \textbf{L'emprunt (Anglicismes / Francophonismes) :}
    \begin{itemize}
        \item \textbf{Identifier l'emprunt :} Mot venu d'une autre langue (souvent l'anglais en milieu technique).
        \item \textbf{Règle d'intégration :} Orthographe francisée (accent, pluriel en \emph{-s}), genre (souvent masculin), prononciation adaptée.
        \item \textbf{Attention aux « faux emprunts » ou anglicismes syntaxiques :}
        \begin{center}
            \begin{tabular}{|l|l|l|}
            \hline
            \textbf{Anglicisme / Erreur} & \textbf{Correction (Français standard)} & \textbf{Type} \\
            \hline
            \emph{Faire un stage \textbf{sur} l'électrotechnique} & \emph{Faire un stage \textbf{dans / en} électrotechnique} & Préposition \\
            \emph{\textbf{Opportunité} de stage} & \emph{\textbf{Possibilité / Occasion} de stage} & Faux ami (Opportunité = à propos) \\
            \emph{\textbf{Valider} mon stage} & \emph{\textbf{Valider / Réussir / Obtenir la validation de} mon stage} & Verbe transitif direct \\
            \emph{Un \textbf{stage} de 2 mois} & \emph{Une \textbf{période de stage} de 2 mois} & Précision \\
            \hline
            \end{tabular}
        \end{center}
        \item \textbf{Emprunts techniques acceptés (ne pas « sur-traduire ») :} \emph{Software, hardware, email, stage, planning, process, management, marketing, leader, briefing, reporting}. Les écrire en \emph{italique} si très récents, sinon roman.
    \end{itemize}
\end{enumerate}
\cle{Lexique spécialisé} \cle{Emprunt intégré} \cle{Anglicisme syntaxique} \cle{Faux ami}
\end{methode}

\begin{exoresolu}[Correction lexicale et stylistique d'un corps de requête]
\textbf{Énoncé :} Le paragraphe suivant est extrait du corps d'une requête adressée au Proviseur pour une autorisation d'absence. Il contient des fautes de registre, des anglicismes syntaxiques et des imprécisions. Réécrivez-le en français administratif standard.

\emph{« Monsieur le Proviseur, je veux vous demander la permission de ne pas venir au lycée la semaine prochaine parce que j'ai un \textbf{stage} chez \textbf{ENEO}. C'est une grosse \textbf{opportunité} pour mon \textbf{curriculum vitae}. Je vais \textbf{apprendre} plein de trucs sur le \textbf{software} des compteurs intelligents. Je \textbf{valide} mon BTS l'année prochaine, donc c'est important. Je vous joins les \textbf{papiers} demandés. Merci de \textbf{valider} ma demande. »}

\tcblower
\textbf{Solution détaillée :}

\textbf{Analyse des erreurs :}
\begin{enumerate}
    \item \emph{« je veux vous demander la permission de ne pas venir »} $\rightarrow$ Registre familier / oral. \textbf{Correction :} \emph{Je sollicite l'autorisation de m'absenter.}
    \item \emph{« la semaine prochaine »} $\rightarrow$ Imprécis (date manquante). \textbf{Correction :} \emph{du [date] au [date].}
    \item \emph{« chez ENEO »} $\rightarrow$ Familier. \textbf{Correction :} \emph{au sein de l'entreprise ENEO / à la direction régionale de l'ENEO.}
    \item \emph{« grosse opportunité »} $\rightarrow$ Anglicisme sémantique (\emph{opportunity} = occasion favorable). \textbf{Correction :} \emph{excellente occasion / opportunité pertinente (si sens « à propos » gardé, mais « occasion » plus sûr).}
    \item \emph{« curriculum vitae »} $\rightarrow$ Emprunt latin intégré, correct. Garder \emph{CV} (sigle) ou \emph{curriculum vitæ}.
    \item \emph{« apprendre plein de trucs »} $\rightarrow$ Très familier. \textbf{Correction :} \emph{acquérir des compétences / approfondir mes connaissances.}
    \item \emph{« software »} $\rightarrow$ Emprunt technique anglais. \textbf{Correction :} \emph{logiciel / progiciel / firmware (si embarqué).}
    \item \emph{« Je valide mon BTS »} $\rightarrow$ Anglicisme syntaxique (\emph{to validate} $\neq$ valider un diplôme en français standard, on \emph{obtient} ou \emph{réussit}). \textbf{Correction :} \emph{Je prépare / Je suis en vue de l'obtention du BTS.}
    \item \emph{« papiers »} $\rightarrow$ Familier. \textbf{Correction :} \emph{pièces justificatives / documents requis.}
    \item \emph{« valider ma demande »} $\rightarrow$ Anglicisme (\emph{approve}). \textbf{Correction :} \emph{accorder une suite favorable / donner votre accord.}
\end{enumerate}

\textbf{Version corrigée (Registre administratif soutenu) :}
\begin{quote}
\emph{« Monsieur le Proviseur,}
\emph{J'ai l'honneur de solliciter votre bienveillance afin d'obtenir une autorisation d'absence pour la période du [Date de début] au [Date de fin].}
\emph{Cette absence est motivée par la réalisation d'un stage d'immersion professionnelle au sein de la direction régionale de l'ENEO. Cette expérience constitue une occasion précieuse d'enrichir mon curriculum vitæ et d'acquérir des compétences pratiques en électrotechnique, notamment sur les logiciels de gestion des compteurs intelligents.}
\emph{Actuellement élève en Terminale TEE, en préparation du Brevet de Technicien Supérieur (BTS), ce stage représente une étape déterminante de ma formation.}
\emph{Vous trouverez ci-joint les pièces justificatives constituant mon dossier de stage (convention, attestation d'assurance, programme prévisionnel).}
\emph{En vous remerciant par avance de l'attention que vous voudrez bien réserver à ma demande, je vous prie d'agréer, Monsieur le Proviseur, l'expression de mon respectueux dévouement. »}
\end{quote}
\end{exoresolu}

\courssub{Rédaction guidée : entête, introduction, corps et conclusion}

\begin{methode}[Rédiger les quatre segments obligatoires de la requête]
\textbf{Objectif :} Maîtriser les formules figées (rituels) et la progression argumentative.
\begin{enumerate}
    \item \textbf{L'En-tête (Identification) :}
    \begin{itemize}
        \item \textbf{Expéditeur :} Prénom NOM (majuscules), Adresse complète, Contacts. \textbf{Spécialité / Classe} (ex: \emph{Terminale TEE / Matricule 12345}).
        \item \textbf{Destinataire :} \emph{À l'attention de [Qualité] [Nom si connu] / [Service] / [Adresse]}. Qualité exacte : \emph{Monsieur le Directeur Général, Madame la Proviseure, Monsieur le Ministre}.
        \item \textbf{Lieu, Date} : \emph{Fait à [Ville], le [Jour] [Mois] [Année]}.
        \item \textbf{Objet} : \textbf{Gras, court, nominal}. \emph{Objet : Demande de [nature] / Autorisation de [action] / Renouvellement de [document]}.
    \end{itemize}
    \item \textbf{L'Introduction (Exorde - « Capter l'attention ») :}
    \begin{itemize}
        \item \textbf{Formule d'appel :} \emph{Monsieur le Directeur Général,} (sans « Cher », sans nom propre).
        \item \textbf{Identité + Qualité :} \emph{Je soussigné(e) [Prénom NOM], élève en [Classe/Spécialité] au [Établissement], matricule n°[...],} 
        \item \textbf{Motif d'écriture (Rappel Objet) :} \emph{ai l'honneur de solliciter / de vous adresser la présente pour...}
        \item \textbf{Fondement (Optionnel mais fort) :} \emph{Conformément à l'arrêté n°... / Dans le cadre de ma formation...}
    \end{itemize}
    \item \textbf{Le Corps (Exposition des motifs - « Convaincre ») :}
    \begin{itemize}
        \item \textbf{Paragraphes distincts} (1 idée = 1 paragraphe).
        \item \textbf{Connecteurs logiques :} \emph{En premier lieu, Par ailleurs, De surcroît, En effet, C'est pourquoi.}
        \item \textbf{Arguments types :} Situation personnelle (éloignement, santé), Mérite (notes, rang), Nécessité curriculaire (stage obligatoire), Contraintes extérieures.
        \item \textbf{Preuves :} \emph{Comme en attestent les pièces jointes n°1 et 2...}
    \end{itemize}
    \item \textbf{La Conclusion (Pétition + Salutation) :}
    \begin{itemize}
        \item \textbf{La Demande (Pétition) :} \emph{En conséquence, je vous prie de bien vouloir [verbe infinitif : m'accorder, m'autoriser, délivrer, signer]...} (Clair, direct, unique).
        \item \textbf{Disponibilité :} \emph{Je me tiens à votre entière disposition pour tout complément d'information / entretien.}
        \item \textbf{Formule de politesse finale (Codifiée) :} Adapter au rang du destinataire.
        \begin{center}
            \begin{tabular}{|l|l|}
            \hline
            \textbf{Destinataire} & \textbf{Formule type (Bac)} \\
            \hline
            Ministre / Directeur Général / Proviseur & \emph{Je vous prie d'agréer, [Qualité], l'assurance de ma considération distinguée.} \\
            Chef de service / Enseignant & \emph{Je vous prie de croire, [Qualité], à l'assurance de mes salutations respectueuses.} \\
            Administration générale & \emph{En vous remerciant par avance, je vous prie d'agréer, [Qualité], mes salutations distinguées.} \\
            \hline
            \end{tabular}
        \end{center}
    \end{itemize}
\end{enumerate}
\cle{Exorde} \cle{Pétition} \cle{Formules codifiées} \cle{Connecteurs logiques}
\end{methode}

\begin{exoresolu}[Rédaction segmentée d'une demande de duplicata de diplôme]
\textbf{Énoncé :} Vous êtes \textbf{Tchounga Grâce}, titulaire du Baccalauréat Technique Série F4 (Électrotechnique) session 2023, numéro de décision 2023/0045/MINESEC. Vous avez égaré votre diplôme original. Rédigez \textbf{uniquement} l'En-tête, l'Introduction, le Corps (2 arguments) et la Conclusion (Pétition + Formule de politesse) de la requête adressée au \textbf{Ministre des Enseignements Secondaires (MINESEC)}. Inventez les coordonnées plausibles.

\tcblower
\textbf{Solution détaillée :}

\textbf{EN-TÊTE}
\begin{flushleft}
Tchounga Grâce \\
B.P. 1245 Yaoundé - Quartier Mvog-Ada \\
Tel : 6 9X XX XX XX / Email : grâce.tchounga@email.cm \\
\textbf{Terminale F4 Électrotechnique (Ancien élève) - Décision n° 2023/0045/MINESEC}
\vspace{0.5cm}

À l'attention de Monsieur le Ministre des Enseignements Secondaires \\
B.P. 4510 Yaoundé \\
\vspace{0.5cm}

\textbf{Fait à Yaoundé, le 20 mai 2025}

\textbf{Objet : Demande de délivrance d'un duplicata de diplôme de Baccalauréat Technique.}
\end{flushleft}

\textbf{INTRODUCTION}
\begin{quote}
Monsieur le Ministre,

Je soussignée \textbf{Tchounga Grâce}, ancienne élève de la Terminale F4 Électrotechnique au Lycée Technique de Yaoundé, titulaire du Baccalauréat Technique session 2023 sous le numéro de décision \textbf{2023/0045/MINESEC}, ai l'honneur de solliciter votre haute bienveillance pour l'obtention d'un duplicata de mon diplôme original.
\end{quote}

\textbf{CORPS (Exposition des motifs)}
\begin{quote}
\textbf{1. Circonstances de la perte.} Par une maladresse regrettable lors de mon déménagement récent vers le quartier Mvog-Ada, j'ai égaré le dossier contenant mon diplôme original ainsi que mon relevé de notes. Malgré des recherches approfondies, ce document est demeuré introuvable. Une déclaration de perte a été déposée au Commissariat du 2\textsuperscript{e} Arrondissement de Yaoundé le 15 mai 2025 (récépissé n° 4587/YPJ joint en pièce n°1).

\textbf{2. Nécessité urgente.} Ce duplicata m'est indispensable pour finaliser mon inscription en première année de Licence Professionnelle en Maintenance des Systèmes Électriques à l'Institut Universitaire de la Côte (IUC) pour l'année académique 2025-2026. La date limite de dépôt des dossiers étant fixée au 31 mai 2025, l'urgence de la présente démarche est évidente.
\end{quote}

\textbf{CONCLUSION}
\begin{quote}
En conséquence, \textbf{je vous prie de bien vouloir m'accorder la délivrance d'un duplicata de mon diplôme de Baccalauréat Technique Série F4, session 2023.}

Je me tiens à votre entière disposition pour tout complément d'information ou vérification auprès de mes anciens établissements.

En vous remerciant par avance de l'attention que vous voudrez bien réserver à ma demande, \textbf{je vous prie d'agréer, Monsieur le Ministre, l'assurance de ma très haute considération.}

\textbf{Tchounga Grâce} \\
\textit{(Signature)}
\end{quote}

\textbf{BAS DE PAGE (Indicatif)}
\begin{flushleft}
\textbf{Pièces jointes :} \\
1. Déclaration de perte (Commissariat) ; 2. Photocopie du diplôme (si disponible) / Attestation de réussite ; 3. Photocopie CNI ; 4. Timbres fiscaux (selon tarif en vigueur) ; 5. Enveloppe timbrée à mon adresse.
\end{flushleft}
\end{exoresolu}

\courssub{Exploitation des textes iconiques pour préparer la requête}

\begin{methode}[Analyser un document iconique (tableau, graphique, organigramme) pour étayer une requête]
\textbf{Objectif :} Transformer une information visuelle en argument écrit (chiffres $\rightarrow$ mots, structure $\rightarrow$ justification).
\begin{enumerate}
    \item \textbf{Identifier la nature et la source du document :} Type (tableau stats, organigramme service, plan local, courbe évolution), Titre, Source, Date, Unités.
    \item \textbf{Extraire les données pertinentes (Sélection) :} Ne pas tout décrire. Choisir les chiffres/relations qui \textbf{justifient la demande}.
    \begin{itemize}
        \item \emph{Ex: Graphique « Taux d'insertion pro par filière »} $\rightarrow$ Chiffre de sa propre filière (argument force) ou filière cible (argument besoin).
        \item \emph{Ex: Organigramme entreprise} $\rightarrow$ Identifier le bon service/destinataire + montrer qu'on connaît la structure.
    \end{itemize}
    \item \textbf{Verbaliser (Traduire en phrases) :} Passer du visuel au textuel.
    \begin{center}
        \begin{tabular}{|l|l|}
        \hline
        \textbf{Visuel} & \textbf{Verbalisation type} \\
        \hline
        Barre à 85\% & \emph{« Le taux d'insertion atteint 85\%, ce qui témoigne de... »} \\
        Flèche montante & \emph{« On observe une progression constante de +12\% sur 3 ans. »} \\
        Case « Chef Service Maintenance » & \emph{« Votre service Maintenance, dirigé par M. X, compte 15 techniciens... »} \\
        \hline
        \end{tabular}
    \end{center}
    \item \textbf{Intégrer dans le Corps de la requête :} Insérer la verbalisation comme \textbf{preuve (argument d'autorité / fait)}.
    \begin{itemize}
        \item \emph{« Comme le montre le tableau ci-joint (Doc 1), le taux d'encadrement... »}
        \item \emph{« L'organigramme de votre direction (cf. pièce jointe) révèle que... »}
    \end{itemize}
    \item \textbf{Citer la source :} \emph{(Source : Rapport d'activité 2024 / Site web officiel / Enquête MINESEC)}.
\end{enumerate}
\cle{Verbalisation} \cle{Transcodage} \cle{Argument d'autorité} \cle{Source}
\end{methode}

\begin{exoresolu}[Exploitation d'un tableau statistique pour une demande de bourse]
\textbf{Énoncé :} Vous postulez à une bourse d'excellence « Femmes \& Sciences ». Le dossier demande de justifier le choix de votre filière (F4 Électrotechnique) par des données chiffrées. On vous fournit le tableau ci-dessous (extrait de l'\emph{Annuaire Statistique MINESEC 2024}).

\begin{center}
\begin{tabular}{|l|c|c|c|}\hline
\textbf{Filière Technique} & \textbf{Effectif Filles (2024)} & \textbf{Effectif Total} & \textbf{Taux de féminisation} \\ \hline
F1 Mécanique & 45 & 1 200 & 3,75\% \\ \hline
\textbf{F4 Électrotechnique} & \textbf{112} & \textbf{980} & \textbf{11,4\%} \\ \hline
F7 Génie Civil & 78 & 1 500 & 5,2\% \\ \hline
F10 Informatique & 320 & 2 100 & 15,2\% \\ \hline
\textbf{Moyenne Générale} & \textbf{---} & \textbf{---} & \textbf{8,5\%} \\ \hline
\end{tabular}
\end{center}

\textbf{Consigne :} Rédigez le paragraphe argumentatif du corps de la requête qui exploite ce tableau pour prouver : (1) votre mérite (réussite dans une filière peu féminisée) et (2) la pertinence de la bourse pour réduire l'écart.

\tcblower
\textbf{Solution détaillée :}

\textbf{Étape 1 : Analyse du tableau (Sélection)}
\begin{itemize}
    \item Ma filière : F4 Électrotechnique.
    \item Taux de féminisation F4 : \textbf{11,4\%} (supérieur à la moyenne générale 8,5\%, mais encore faible).
    \item Effectif filles : 112 sur 980 (minorité significative).
    \item Comparaison : F1 (3,75\%) et F7 (5,2\%) sont plus critiques ; F10 (15,2\%) est plus féminisée.
\end{itemize}

\textbf{Étape 2 : Verbalisation et Argumentation (Rédaction)}

\begin{quote}
\textbf{Paragraphe argumentatif (Corps de la requête) :}

« Monsieur le Ministre, mon engagement dans la filière F4 Électrotechnique s'inscrit dans une démarche volontariste de dépassement des stéréotypes de genre dans les métiers techniques. Les données officielles de l'Annuaire Statistique MINESEC 2024 révèlent en effet que, si le taux moyen de féminisation dans l'enseignement technique industriel s'établit à \textbf{8,5\%}, la filière Électrotechnique (F4) n'atteint que \textbf{11,4\%} avec seulement \textbf{112 filles} sur un effectif total de 980 élèves. 

Bien que ce taux soit légèrement supérieur à la moyenne, il témoigne d'une sous-représentation persistante des jeunes filles dans les métiers de l'énergie et de l'électricité. Mon admission en Terminale F4, couronnée par un rang de \textbf{3\textsuperscript{e} / 42} à l'examen blanc régional (Pièce n°3), démontre ma capacité à exceller dans cet environnement exigeant. 

L'obtention de la bourse « Femmes \& Sciences » constituerait, au-delà d'un soutien financier indispensable à la poursuite de mes études en BTS Électrotechnique, un signal fort encourageant la vocations scientifiques féminines dans les filières où l'écart de genre demeure structurel, conformément aux objectifs de la Stratégie Nationale de Développement 2020-2030 (Axe 3 : Capital humain). »
\end{quote}

\textbf{Commentaire de la solution :}
\begin{itemize}
    \item \textbf{Transcodage réussi :} Chiffres (11,4\%, 112/980, 8,5\%) intégrés dans des phrases syntaxiques correctes.
    \item \textbf{Double argumentation :} Mérite personnel (rang 3/42) + Cause collective (réduction écart genre).
    \item \textbf{Source citée :} « Annuaire Statistique MINESEC 2024 » donne de la crédibilité (argument d'autorité).
    \item \textbf{Lien avec l'appel à projet :} Référence à la stratégie nationale (contexte institutionnel).
\end{itemize}
\end{exoresolu}

\courssub{Analyse de situation, plan détaillé et rédaction complète}

\begin{methode}[De l'analyse du sujet à la rédaction finale : la méthode en 4 temps]
\textbf{Objectif :} Gérer le temps d'examen (2h-3h) et produire un texte cohérent, complet, sans hors-sujet.
\begin{enumerate}
    \item \textbf{Temps 1 : Analyse de la situation de communication (15 min) - QQOQCP}
    \begin{center}
        \begin{tabular}{|l|l|l|}
        \hline
        \textbf{Question} & \textbf{Réponse à extraire du sujet} & \textbf{Impact sur l'écriture} \\
        \hline
        \textbf{Qui ?} (Expéditeur) & Identité, Qualité, Spécialité, Matricule & En-tête, Introduction (Exorde) \\
        \textbf{À Qui ?} (Destinataire) & Qualité exacte, Service, Adresse & En-tête, Formule d'appel, Politesse finale \\
        \textbf{Quoi ?} (Objet précis) & Nature exacte de la demande & Objet, Pétition (Conclusion) \\
        \textbf{Pourquoi ?} (Motifs) & Causes, Justificatifs, Enjeux & Corps (Argumentation) \\
        \textbf{Quand ?} (Délais, Dates) & Urgence, Dates de stage, Échéances & Corps (Argument urgence), Date lettre \\
        \textbf{Comment ?} (Pièces, Procédure) & Dossier requis, Convention, Timbres & Bas de page (P.J.), Corps (réf. procédure) \\
        \hline
    \end{tabular}
    \end{center}
    \item \textbf{Temps 2 : Élaboration du Plan Détaillé (15 min)}
    \begin{itemize}
        \item \textbf{Introduction :} [Formule appel] + [Identité/Qualité] + [Objet] + [Fondement].
        \item \textbf{Corps - §1 (Contexte/Fait déclencheur) :} [Situation actuelle] + [Preuve/Pièce 1].
        \item \textbf{Corps - §2 (Argument principal/Mérite/Besoin) :} [Argument fort] + [Chiffre/Exemple/Pièce 2].
        \item \textbf{Corps - §3 (Argument secondaire/Engagement) :} [Disponibilité/Projet futur] + [Pièce 3].
        \item \textbf{Conclusion :} [Pétition formelle] + [Disponibilité entretien] + [Formule politesse codifiée].
        \item \textbf{Bas de page :} [Signature] + [Liste P.J. numérotée].
    \end{itemize}
    \item \textbf{Temps 3 : Rédaction soignée (45-60 min)}
    \begin{itemize}
        \item Respecter la \textbf{mise en page} (marges, sauts de ligne).
        \item Utiliser le \textbf{lexique spécialisé} (solliciter, octroyer, pièces justificatives).
        \item Soigner la \textbf{syntaxe} : phrases complexes (subordination), connecteurs, concordance des temps.
        \item \textbf{Éviter :} Familier, abréviations non standards (sauf sigles connus : MINESEC, CNI, BTS), répétitions.
    \end{itemize}
    \item \textbf{Temps 4 : Relecture / Auto-correction (15-20 min) - Grille de vérification}
    \begin{itemize}
        \item [ ] \textbf{Externe :} En-têtes complets ? Objet gras ? Date ? Lieu ? Signature ? P.J. listées ?
        \item [ ] \textbf{Interne :} Identité claire ? Objet rappelé ? Arguments logiques ? Pétition explicite ? Politesse adaptée au destinataire ?
        \item [ ] \textbf{Langue :} Orthographe (accords, homophones) ? Ponctuation ? Registre soutenu ? Pas d'anglicismes syntaxiques ?
    \end{itemize}
\end{enumerate}
\cle{QQOQCP} \cle{Plan détaillé} \cle{Pétition} \cle{Auto-correction}
\end{methode}

\begin{exoresolu}[Rédaction complète : Demande de convention de stage (Sujet type Bac)]
\textbf{Énoncé :} Vous êtes \textbf{Kamga Kevin}, élève en \textbf{Terminale TEE (Électrotechnique)} au \textbf{Lycée Technique de Bafoussam}, matricule \textbf{LT/BFS/2025/0145}. Dans le cadre du programme officiel, vous devez effectuer un \textbf{stage de fin d'études de 8 semaines (1er juillet - 26 août 2025)}. Vous ciblez l'entreprise \textbf{CAMWATER (Direction Régionale de l'Ouest - Bafoussam)}. Rédigez la requête complète adressée au \textbf{Directeur Régional de CAMWATER Ouest}. Inventez les pièces jointes nécessaires.

\tcblower
\textbf{Solution détaillée :}

\textbf{ÉTAPE 1 : ANALYSE QQOQCP (Brouillon)}
\begin{itemize}
    \item \textbf{Expéditeur :} Kamga Kevin, TEE, LT Bafoussam, Matr. LT/BFS/2025/0145.
    \item \textbf{Destinataire :} Directeur Régional CAMWATER Ouest, Bafoussam.
    \item \textbf{Objet :} Demande de convention de stage fin d'études (8 semaines, Juil/Août 2025).
    \item \textbf{Motifs :} Obligation curriculaire (programme MINESEC), Spécialité Eau/Énergie (pertinent CAMWATER), Mérite (notes), Disponibilité.
    \item \textbf{Pièces :} CV, Lettre motivation, Attestation scolarité, Programme stage, Convention type (fourni par lycée), Assurance.
\end{itemize}

\textbf{ÉTAPE 2 : PLAN DÉTAILLÉ}
\begin{enumerate}
    \item \textbf{En-têtes + Objet}
    \item \textbf{Intro :} Appel + Identité (Kevin, TEE, LT Bafoussam) + Objet (Stage 8 sem) + Fondement (Programme officiel / Arrêté d'organisation).
    \item \textbf{Corps §1 :} Contexte pédagogique (Stage obligatoire, dates fixes, objectifs : pratique pompage/traitement).
    \item \textbf{Corps §2 :} Adéquation profil/entreprise (Spé Électrotechnique = maintenance pompes/automatismes, Bonnes notes sciences ingénieur).
    \item \textbf{Corps §3 :} Engagement / Disponibilité (Respect règlement intérieur, Encadrement tuteur, Assurance).
    \item \textbf{Conclusion :} Pétition (Convention + Accueil) + Dispo entretien.
    \item \textbf{Politesse + Signature + P.J.}
\end{enumerate}

\textbf{ÉTAPE 3 : RÉDACTION FINALE (Copie propre)}

\begin{flushleft}
\textbf{Kamga Kevin} \\
B.P. 352 Bafoussam - Quartier Banengo \\
Tél : 6 7X XX XX XX \hfill Email : kevin.kamga@email.cm \\
\textbf{Élève Terminale TEE - Matricule : LT/BFS/2025/0145} \\
\vspace{0.5cm}

\textbf{À l'attention de Monsieur le Directeur Régional de CAMWATER Ouest} \\
B.P. 145 Bafoussam \\
\vspace{0.5cm}

\textbf{Fait à Bafoussam, le 15 avril 2025}

\textbf{Objet : Demande de convention de stage de fin d'études (Juillet - Août 2025).}
\end{flushleft}

Monsieur le Directeur Régional,

Je soussigné \textbf{Kamga Kevin}, élève en \textbf{Terminale Technicien Électrotechnique (TEE)} au \textbf{Lycée Technique de Bafoussam}, matricule \textbf{LT/BFS/2025/0145}, ai l'honneur de solliciter votre bienveillance pour l'obtention d'une convention de stage de fin d'études au sein de votre direction, conformément aux dispositions de l'arrêté n° 004/MINESEC/SG/DEST du 12 janvier 2010 fixant l'organisation des stages en milieu professionnel.

\textbf{1. Un impératif curriculaire aux dates fixes.}
Dans le cadre de la validation de mon année terminale, la réalisation d'un stage d'immersion professionnelle de \textbf{huit (8) semaines} est une obligation pédagogique. Ce stage doit se dérouler impérativement \textbf{du 1er juillet au 26 août 2025}. Il vise à me confronter aux réalités techniques de la production et de la distribution d'eau potable, cœur de métier de CAMWATER, et à mettre en application mes acquis en électrotechnique (commande de moteurs, automatismes programmables, maintenance préventive).

\textbf{2. Une adéquation profil / structure technique.}
Mon choix pour votre direction régionale n'est pas fortuit. La spécialité Électrotechnique que je prépare trouve un champ d'application direct dans la gestion des stations de pompage, du traitement de l'eau et des réseaux de distribution automatisés. Mon dossier scolaire (moyenne annuelle 14,5/20 en Sciences de l'Ingénieur, Rang 2/38 au 2\textsuperscript{e} trimestre - Pièce n°2) atteste de ma rigueur et de ma capacité à intégrer rapidement vos équipes techniques.

\textbf{3. Un engagement total et encadré.}
Conscient des exigences de sécurité et de continuité de service propres à votre secteur, je m'engage à respecter scrupuleusement le règlement intérieur de votre établissement et les consignes de mon tuteur de stage. Une attestation d'assurance responsabilité civile et accidents du travail (Pièce n°4) ainsi que le modèle de convention type fourni par mon lycée (Pièce n°3) sont joints au présent dossier.

\textbf{En conséquence, je vous prie de bien vouloir m'accorder une convention de stage pour la période susmentionnée et de bien vouloir désigner un tuteur technique au sein de votre service Maintenance ou Production.}

Je me tiens à votre entière disposition pour un entretien à votre convenance afin de préciser les modalités pratiques de cette collaboration.

En vous remerciant par avance de l'attention que vous voudrez bien réserver à ma candidature, \textbf{je vous prie d'agréer, Monsieur le Directeur Régional, l'assurance de ma considération distinguée.}

\vspace{1cm}
\begin{flushright}
\textbf{Kamga Kevin} \\
\textit{(Signature)}
\end{flushright}

\vspace{0.5cm}
\textbf{Pièces jointes :} \\
1. Curriculum Vitæ détaillé ; 2. Lettre de motivation ; 3. Attestation de scolarité 2024-2025 ; 4. Modèle de convention de stage (Lycée Technique) ; 5. Attestation d'assurance scolaire extrascolaire ; 6. Programme prévisionnel des activités de stage (8 semaines).
\end{exoresolu}

\courssub{Confrontation, amélioration, compte rendu et remédiation}

\begin{methode}[Auto-évaluation et amélioration de sa production (Grille Bac)]
\textbf{Objectif :} Devenir son propre correcteur pour gagner les points de « Qualité de la langue » et « Respect des consignes ».
\begin{enumerate}
    \item \textbf{Phase 1 : Confrontation (Lecture à voix basse / silencieuse active)}
    \begin{itemize}
        \item Lire \textbf{à l'envers} (dernière phrase $\rightarrow$ première) pour casser l'automatisme de lecture et voir les fautes d'orthographe (accords, homophones).
        \item Vérifier la \textbf{cohérence énonciative} : Le « Je » est-il stable ? Le destinataire est-il bien le « Vous » tout du long ? (Pas de basculement « on / nous »).
    \end{itemize}
    \item \textbf{Phase 2 : Vérification par items (Grille de 10 points critiques)}
    \begin{center}
        \begin{tabular}{|c|l|c|}
        \hline
        \textbf{N°} & \textbf{Point de contrôle} & \textbf{Action si erreur} \\
        \hline
        1 & \textbf{Destinataire exact} (Qualité + Service) & Corriger en-tête + Appel + Politesse finale (cohérence trio) \\
        2 & \textbf{Objet} : Nominal, Précis, Gras & Reformuler (pas de phrase verbe) \\
        3 & \textbf{Identité complète} dans l'intro (Nom, Classe, Spé, Matricule, Établissement) & Compléter \\
        4 & \textbf{Pétition (Conclusion)} : Verbe infinitif clair (\emph{m'accorder, m'autoriser, délivrer}) & Réécrire la phrase clé \\
        5 & \textbf{Formule de politesse} : Adaptée au rang (Ministre $\neq$ Chef service) & Choisir dans le tableau de référence \\
        6 & \textbf{Pièces jointes} : Listées, Numérotées, Cohérentes avec le corps & Ajouter / Supprimer / Renuméroter \\
        7 & \textbf{Connecteurs logiques} (Au moins 3 variés : \emph{En outre, Par conséquent, Toutefois}) & Insérer / Varier \\
        8 & \textbf{Lexique spécialisé} (0 terme familier, 0 anglicisme syntaxique) & Substituer (Voir Méthode 3) \\
        9 & \textbf{Accords / Homophones} (étant/étant, a/à, son/sont, ce/se, ou/où) & Corriger au stylo rouge \\
        10 & \textbf{Mise en page} : Aération (saut ligne §), Alignements, Lisibilité signature & Réorganiser visuellement \\
        \hline
    \end{tabular}
    \end{center}
    \item \textbf{Phase 3 : Réécriture ciblée (Remédiation)}
    \begin{itemize}
        \item Ne recopier \textbf{que les segments défectueux} (gain de temps).
        \item Appliquer la correction \textbf{propre} (pas de ratures multiples).
    \end{itemize}
    \item \textbf{Phase 4 : Compte rendu méta-cognitif (Oral ou Écrit court)}
    \begin{itemize}
        \item \emph{« J'ai corrigé l'accord de "pièces jointes" (féminin pluriel). J'ai remplacé "valider" par "accorder". J'ai ajouté le connecteur "Par ailleurs" au §2. »}
        \item Identifier \textbf{1 point fort} et \textbf{1 axe de progression} pour la prochaine fois.
    \end{itemize}
\end{enumerate}
\cle{Auto-correction} \cle{Grille d'évaluation} \cle{Remédiation} \cle{Méta-cognition}
\end{methode}

\begin{exoresolu}[Correction guidée d'une copie élève (Simulation remédiation)]
\textbf{Énoncé :} Voici une production d'élève (brouillon) pour une demande de bourse. Appliquez la \textbf{Méthode 6 (Grille des 10 points)}. Relevez les erreurs, proposez les corrections et rédigez le \textbf{compte rendu de remédiation}.

\emph{« Yaoundé le 10/05/25 \\
Monsieur le Directeur de l'OCDE, BP 123 Paris \\
Objet : Bourse \\
Monsieur, \\
Je m'appelle Mbarga Paul, je suis en Terminale F3 au Lycée de Ngoa-Ekelle. Je veux une bourse pour aller en France l'an prochain. Mes notes sont bonnes (moyenne 14). Mes parents n'ont pas d'argent. Je joins mes bulletins. \\
Donc donnez-moi la bourse. \\
Merci. \\
Paul Mbarga \\
PJ : Bulletins »}

\tcblower
\textbf{Solution détaillée :}

\textbf{APPLICATION DE LA GRILLE DES 10 POINTS}

\begin{center}
\begin{tabular}{|c|l|l|l|}
\hline
\textbf{N°} & \textbf{Point} & \textbf{Constat (Erreur / Manquant)} & \textbf{Correction} \\
\hline
1 & \textbf{Destinataire} & « Directeur de l'OCDE » (Organisation internationale, pas bourse étude Cameroun). Adresse Paris. & \textbf{Monsieur le Ministre des Enseignements Supérieurs (MINESUP)} ou \textbf{Directeur des Bourses MINESUP}, B.P. Yaoundé. \\
2 & \textbf{Objet} & « Bourse » (Trop vague, pas de gras visible). & \textbf{Objet : Demande de bourse d'études supérieures pour l'année académique 2025-2026.} \\
3 & \textbf{Identité Intro} & Manque : Matricule, Spécialité précise (F3 = Électronique ?), « Je m'appelle » (familier). & \emph{Je soussigné Mbarga Paul, élève en Terminale F3 (Électronique) au Lycée de Ngoa-Ekelle, matricule n°..., ai l'honneur de solliciter...} \\
4 & \textbf{Pétition} & « Donnez-moi la bourse » (Impératif, familier, agressif). & \emph{Je vous prie de bien vouloir m'accorder une bourse d'études...} \\
5 & \textbf{Politesse} & « Merci. » (Inacceptable). & \emph{Je vous prie d'agréer, Monsieur le Ministre, l'assurance de ma considération distinguée.} \\
6 & \textbf{Pièces jointes} & « Bulletins » (Vague, pas numérotées, manque : avis d'admission, quittance, certificat nationalité, extrait naissance, photo, enveloppe). & \textbf{Liste numérotée complète (8-10 pièces standard).} \\
7 & \textbf{Connecteurs} & Aucun (Juste « Donc »). & \emph{En effet / Par ailleurs / C'est pourquoi / En conséquence.} \\
8 & \textbf{Lexique} & « Je veux », « Mes parents n'ont pas d'argent », « l'an prochain ». & \emph{Je sollicite / Situation financière précaire / année académique à venir.} \\
9 & \textbf{Orthographe / Syntaxe} & Date « 10/05/25 » (format admin : « Yaoundé, le 10 mai 2025 »). Phrases courtes, hachées. & \textbf{Date en toutes lettres. Phrases complexes subordonnées.} \\
10 & \textbf{Mise en page} & Tout collé, pas de sauts de paragraphes, pas de signature distincte. & \textbf{Structure aérée : En-têtes / Intro / §1 / §2 / Conclu / Politesse / Sign / P.J.} \\
\hline
\end{tabular}
\end{center}

\textbf{VERSION REMISE EN ÉTAT (Extrait corrigé - Corps + Conclusion)}

\begin{quote}
\emph{« Monsieur le Ministre,}

\emph{Je soussigné \textbf{Mbarga Paul}, élève en \textbf{Terminale F3 (Électronique)} au \textbf{Lycée de Ngoa-Ekelle}, matricule \textbf{n° 2025/NE/045}, ai l'honneur de solliciter votre haute bienveillance pour l'obtention d'une \textbf{bourse d'études supérieures} pour l'année académique 2025-2026.}

\emph{\textbf{En effet}, mon admission en classe terminale avec une moyenne annuelle de \textbf{14/20} (Rang 5/42) témoigne de mon sérieux et de ma capacité à réussir dans l'enseignement supérieur. \textbf{Par ailleurs}, la situation financière précaire de ma famille, mon père étant à la retraite et ma mère sans emploi stable, ne me permet pas d'envisager sereinement les frais d'inscription et de vie inhérents à une poursuite d'études, notamment à l'étranger où j'ai été présélectionné à l'Université de \textbf{[Nom Univ]}.}

\emph{\textbf{C'est pourquoi}, je vous prie de bien vouloir \textbf{m'accorder cette bourse} qui conditionne mon avenir professionnel.}

\emph{Je me tiens à votre disposition pour tout complément d'information.}

\emph{En vous remerciant par avance de l'attention portée à ma demande, \textbf{je vous prie d'agréer, Monsieur le Ministre, l'assurance de ma très haute considération.}}

\emph{\textbf{Mbarga Paul} \hfill \textit{(Signature)}}
\end{quote}

\textbf{COMPTE RENDU DE RÉMÉDIATION (Méta-cognition)}

\begin{quote}
\textbf{Analyse de la copie initiale :} La production relevait du \textbf{registre familier/oral} (« Je veux », « donnez-moi », « l'an prochain ») et ignorait les \textbf{conventions administratives} (en-têtes incomplets, objet vague, politesse absente, pièces jointes insuffisantes). La structure argumentative était inexistante (un seul bloc).

\textbf{Actions de correction effectuées :}
\begin{enumerate}
    \item \textbf{Rectification du destinataire} (MINESUP vs OCDE) pour cohérence contextuelle.
    \item \textbf{Reconstruction de l'en-tête} (Expéditeur complet + Destinataire exact + Date longue + Objet nominal précis).
    \item \textbf{Réécriture de l'introduction} avec formules codifiées (\emph{Je soussigné... ai l'honneur de solliciter}).
    \item \textbf{Structuration du corps en 2 paragraphes} distincts : Mérite scolaire (Argument 1) + Situation sociale (Argument 2), reliés par connecteurs (\emph{En effet, Par ailleurs}).
    \item \textbf{Transformation de la pétition} en forme infinitive polie (\emph{m'accorder}).
    \item \textbf{Application de la formule de politesse ministérielle} (Très haute considération).
    \item \textbf{Énumération normée des pièces jointes} (8 items standard).
\end{enumerate}

\textbf{Point fort conservé :} L'honnêteté de la situation sociale et le chiffre de la moyenne (14/20) étaient de bons arguments de fond.

\textbf{Axe de progression pour l'élève :} \textbf{Maîtriser le "Kit de survie administratif"} : 1) Les 3 formules de politesse types (Ministre, DG, Chef Service) par cœur. 2) La liste type des 8 pièces jointes "bourse/stage". 3) La structure Intro (Qui/Quoi/Fondement) / Corps (Arg 1 / Arg 2) / Conclu (Pétition/Politesse). S'entraîner à rédiger sur brouillon en 20 min chrono.
\end{quote}
\end{exoresolu}

\begin{attention}[Les 5 erreurs qui coûtent des points au Bac (Requête)]
\begin{enumerate}
    \item \textbf{Incohérence du "Trio Destinataire" :} Écrire « Monsieur le Directeur » en en-tête, « Monsieur le Proviseur » dans l'appel, et « Monsieur le Ministre » dans la formule de politesse finale. \textbf{Sanction :} Perte de points \textbf{Structure externe} + \textbf{Qualité de langue} (incohérence énonciative). \textbf{Règle :} Identifier \textbf{UNE SEULE FOIS} le destinataire exact sur le brouillon et reporter partout.
    \item \textbf{Pétition floue ou absente :} Terminer par « En espérant... » sans dire \textbf{précisément} ce qu'on demande (\emph{m'accorder, m'autoriser, me délivrer}). \textbf{Sanction :} Zéro point pour la \textbf{Conclusion / Pétition} (cœur de la requête). \textbf{Règle :} La dernière phrase du corps (avant politesse) doit commencer par \emph{« Je vous prie de bien vouloir [Verbe Infinitif] [Objet précis]... »}.
    \item \textbf{Registre familier / Anglicismes syntaxiques :} « Je \textbf{valide} ma demande », « C'est une \textbf{opportunité} », « Je \textbf{gère} le dossier », « \textbf{Stage} de 2 mois » (sans précision). \textbf{Sanction :} Pénalité \textbf{Qualité de langue} (Registre). \textbf{Règle :} Substitution systématique : \emph{Accorder / Occasion / Traiter / Période de stage}.
    \item \textbf{Pièces jointes fantaisistes ou manquantes :} Lister « CV, Lettre » seulement pour une bourse (il faut : extrait naissance, certificat nationalité, quittance, avis admission, photo, enveloppe timbrée...). Ou ne pas lister du tout. \textbf{Sanction :} Perte points \textbf{Structure externe} (Bas de page) + \textbf{Cohérence} (le corps cite des preuves non listées). \textbf{Règle :} Mémoriser les \textbf{listes types} par nature de requête (Bourse, Stage, Diplôme, Autorisation).
    \item \textbf{Mise en page "brouillon" :} Pas de saut de ligne entre paragraphes, en-têtes collés, objet non visible, signature au milieu du texte, date en chiffres (10/05/25). \textbf{Sanction :} Perte points \textbf{Présentation / Structure externe} (souvent 2 à 4 pts sur 20). \textbf{Règle :} \textbf{Aérer} : 1 ligne vide entre chaque bloc (En-têtes / Objet / Appel / §1 / §2 / Pétition / Politesse / Signature / P.J.). Date en \textbf{toutes lettres}.
\end{enumerate}
\end{attention}

\newpage

\coursec{Exercices}

\courssub{Je vérifie mes connaissances}

\begin{exercice}[QCM : Actes de langage et structure de la requête]
Parmi les propositions suivantes, une seule est exacte. Recopiez la lettre correspondante.
\begin{enumerate}
    \item Dans une requête administrative, l'acte illocutoire principal est généralement :
    \begin{itemize}
        \item[A.] Une assertion (décrire un fait).
        \item[B.] Une directive (demander, exiger, solliciter).
        \item[C.] Une commissive (promettre, s'engager).
        \item[D.] Une expressive (remercier, féliciter).
    \end{itemize}
    \item L'entête d'une requête formelle adressée au Proviseur d'un lycée technique doit obligatoirement mentionner :
    \begin{itemize}
        \item[A.] Le numéro de téléphone personnel de l'expéditeur.
        \item[B.] La qualité de l'expéditeur et la qualité du destinataire.
        \item[C.] La date de naissance de l'expéditeur.
        \item[D.] Le matricule de l'élève concerné (si l'expéditeur est un parent).
    \end{itemize}
    \item L'« Objet » d'une requête se caractérise par :
    \begin{itemize}
        \item[A.] Une phrase complète se terminant par un point.
        \item[B.] Une phrase nominale concise, sans verbe conjugué, souvent précédée de « Objet : ».
        \item[C.] Un paragraphe développant les motifs de la demande.
        \item[D.] La formule de politesse finale.
    \end{itemize}
    \item L'effet perlocutoire attendu d'une requête de réclamation est :
    \begin{itemize}
        \item[A.] Que le destinataire lise le courrier.
        \item[B.] Que le destinataire comprenne la colère de l'expéditeur.
        \item[C.] Que le destinataire prenne une décision favorable (réparation, correction, nomination).
        \item[D.] Que le destinataire archive le document.
    \end{itemize}
\end{enumerate}
\end{exercice}

\begin{exercice}[Vrai / Faux justifié : Lexique spécialisé et emprunts]
Pour chaque affirmation, cochez V (Vrai) ou F (Faux) et justifiez brièvement votre réponse en citant la règle ou l'équivalent correct.
\begin{enumerate}
    \item \textbf{V / F} : Dans une requête administrative, l'emploi du « camfranglais » (ex : « Je viens \emph{demander} pour mon fils ») est toléré pour la clarté.
    \item \textbf{V / F} : Le terme « \emph{ci-joint} » s'accorde en genre et en nombre avec le nom qui le précède lorsqu'il est placé après ce nom (ex : « les pièces \emph{ci-jointes} »).
    \item \textbf{V / F} : L'emprunt lexical « \emph{planning} » est un anglicisme qu'il faut remplacer par « \emph{calendrier} » ou « \emph{programme} » dans un écrit formel.
    \item \textbf{V / F} : La formule « \emph{Je vous prie d'agréer, Monsieur le Directeur, l'expression de mes salutations distinguées} » est une formule de clôture adaptée à une requête adressée à un chef d'établissement.
\end{enumerate}
\end{exercice}

\begin{exercice}[Définitions et identification : Textes fonctionnels et iconiques]
\begin{enumerate}
    \item Définissez le terme \cle{texte fonctionnel} en donnant deux caractéristiques principales qui le distinguent du texte littéraire.
    \item Citez trois éléments de la \cle{structure externe} de la requête (mise en page, présentation).
    \item Citez les trois parties constitutives de la \cle{structure interne} de la requête (contenu rédactionnel).
    \item À partir de l'organigramme simplifié ci-dessous d'un Lycée Technique, identifiez le destinataire hiérarchique direct d'une requête déposée par un parent d'élève pour contester une note.
\end{enumerate}
\begin{popfigure}
\begin{center}
\begin{tikzpicture}[node distance=1.2cm and 0.5cm, every node/.style={draw, rounded corners, fill=popblueL, text width=2.5cm, align=center, font=\small}]
    \node (prov) {Proviseur};
    \node (censeur) [below left=of prov] {Censeur / Surveillant Général};
    \node (prof) [below right=of prov] {Professeur Principal};
    \node (eleve) [below=of censeur] {Élève / Parent d'élève};
    \draw[->, thick, popdark] (prov) -- (censeur);
    \draw[->, thick, popdark] (prov) -- (prof);
    \draw[->, thick, popdark] (censeur) -- (eleve);
    \draw[->, thick, popdark] (prof) -- (eleve);
\end{tikzpicture}
\end{center}
\legende{Organigramme simplifié d'un lycée technique : la voie hiérarchique pour une contestation de note passe par le Professeur Principal, puis le Censeur, avant le Proviseur.}
\end{popfigure}
\end{exercice}

\begin{exercice}[Calculs directs et repérage : Taux d'erreurs et structure]
Un élève a rédigé une requête de 180 mots. Lors de la relecture, le professeur a relevé :
\begin{itemize}
    \item 3 erreurs de formules de politesse (entête/fin) ;
    \item 5 fautes de lexique spécialisé (anglicismes/camfranglais) ;
    \item 2 erreurs de structure interne (inversion corps/introduction) ;
    \item 4 fautes d'orthographe grammaticale.
\end{itemize}
\begin{enumerate}
    \item Calculez le pourcentage d'erreurs « structurelles et lexicales » (formules + lexique + structure) par rapport au total des erreurs relevées.
    \item Si le barème pénalise 0,5 point par erreur de structure/lexique et 0,25 point par faute d'orthographe, quelle note sur 20 cet élève obtiendrait-il sur la forme (sur une base de 20 points avant pénalité) ?
    \item Parmi les éléments suivants, entourez ceux qui appartiennent au \textbf{corps de la requête} : \textit{Expéditeur ; Exposé des faits ; Formule d'appel ; Arguments ; Pièces jointes ; Demande précise ; Signature ; Objet}.
\end{enumerate}
\end{exercice}

\courssub{Je m'entraîne}

\begin{exercice}[Rédaction de l'entête et de l'objet]
M. \textsc{Kamga} Joseph, parent d'élève, demeurant à Bafoussam (Quartier Banengo, BP 124), souhaite écrire au Proviseur du Lycée Technique de Bafoussam pour contester l'affectation en internat de son fils, \textsc{Kamga} Paul, élève en Terminale F4, matricule 2024-0451.
\begin{enumerate}
    \item Rédigez l'entête complet de la requête (qualités, adresses, lieu et date).
    \item Rédigez l'objet de la requête en respectant les normes administratives (phrase nominale, concise).
\end{enumerate}
\end{exercice}

\begin{exercice}[Rédaction de l'introduction et de la conclusion]
Reprenez la situation de l'exercice précédent (contestation d'affectation en internat).
\begin{enumerate}
    \item Rédigez l'\textbf{introduction} (formule d'appel + phrase d'amorce situant le contexte : qualité de l'expéditeur, qualité du destinataire, référence à la décision contestée).
    \item Rédigez la \textbf{conclusion} (formule de transition + demande précise formulée avec un verbe au subjonctif ou à l'infinitif + formule de politesse finale complète).
\end{enumerate}
\end{exercice}

\begin{exercice}[Rédaction du corps de la requête : Faits, Arguments, Demande]
Toujours dans le cadre de la contestation d'affectation en internat (Paul Kamga, Terminale F4, affecté en internat alors qu'il est externe depuis la Seconde et habite à 500 m du lycée).
\begin{enumerate}
    \item Rédigez l'\textbf{exposé des faits} (chronologie : inscription, notification d'affectation, constat).
    \item Développez \textbf{deux arguments} solides (ex : proximité du domicile, coût financier, réussite scolaire antérieure en externat, état de santé).
    \item Formulez la \textbf{demande précise} (réintégration en externat / modification de la liste).
\end{enumerate}
\emph{Contrainte :} utilisez un lexique administratif soutenu (« sollicite », « y faire droit », « au vu de », « pièces justificatives »). Évitez le camfranglais.
\end{exercice}

\begin{exercice}[Analyse de texte iconique : Organigramme et destinataire]
Observez l'organigramme simplifié de la \cle{SONARA} (Société Nationale de Raffinage) ci-dessous. Un agent de la maintenance (M. \textsc{Ndongo}, matricule 8842) souhaite déposer une requête pour une révision de sa prime de risque non perçue depuis trois mois.
\begin{popfigure}
\begin{center}
\begin{tikzpicture}[node distance=1cm and 0.8cm, every node/.style={draw, rectangle, rounded corners, fill=poporangeL, text width=2.2cm, align=center, font=\small, inner sep=2pt}]
    \node (dg) {Directeur Général};
    \node (drh) [below left=of dg] {DRH};
    \node (dprod) [below right=of dg] {Directeur Production};
    \node (chefmaint) [below=of drh] {Chef Service Maintenance};
    \node (agent) [below=of chefmaint] {Agent Maintenance (M. Ndongo)};
    \draw[->, thick, popdark] (dg) -- (drh);
    \draw[->, thick, popdark] (dg) -- (dprod);
    \draw[->, thick, popdark] (drh) -- (chefmaint);
    \draw[->, thick, popdark] (dprod) -- (chefmaint);
    \draw[->, thick, popdark] (chefmaint) -- (agent);
\end{tikzpicture}
\end{center}
\legende{Organigramme simplifié de la SONARA : la voie hiérarchique impose de passer par le Chef de Service avant la DRH.}
\end{popfigure}
\begin{enumerate}
    \item Qui est le \textbf{destinataire direct} (supérieur hiérarchique immédiat) de M. Ndongo ?
    \item Qui est le \textbf{destinataire final} (décisionnaire) pour une prime de risque ?
    \item Rédigez \textbf{uniquement l'entête et l'objet} de la requête adressée au destinataire direct, en respectant la forme administrative (qualités, adresses, objet).
\end{enumerate}
\end{exercice}

\courssub{Je me prépare au Bac}

\begin{exercice}[Type Baccalauréat Technique : Rédaction complète de requête (20 pts)]
\textbf{Situation :} Vous êtes M. \textsc{Fouda} André, parent d'élève, commerçant au Marché Central de Bafoussam (BP 356). Votre fille, \textsc{Fouda} Grâce, matricule 2024-1122, élève en Terminale G2 au Lycée Technique de Bafoussam, a été \textbf{erronément affectée en internat} pour l'année scolaire 2024-2025. Or, elle est suivie pour une pathologie chronique (asthme sévère) nécessitant un environnement familial strict et des soins nocturnes quotidiens. De plus, le coût de l'internat (150 000 FCFA/an) pèse lourdement sur vos finances. Vous disposez d'un certificat médical du Dr \textsc{Nguyen} (pneumologue, Hôpital Régional de Bafoussam) et d'une attestation d'hébergement.

\textbf{Travail demandé :} Rédigez la \textbf{requête complète} adressée au Proviseur du Lycée Technique de Bafoussam pour obtenir la réintégration de votre fille au régime \textbf{externe}.

\textbf{Barème indicatif :}
\begin{itemize}
    \item Structure externe (entête, objet, pièces jointes, signature) : 4 pts
    \item Structure interne (introduction ; corps : faits + 2 arguments ; conclusion/demande) : 8 pts
    \item Langue (lexique administratif, syntaxe, orthographe, ponctuation, formules de politesse) : 8 pts
\end{itemize}
\end{exercice}

\begin{exercice}[Type Probatoire : Analyse de requête et remise en ordre (15 pts)]
\textbf{Document support :} Voici les éléments mélangés d'une requête adressée par M. \textsc{Essomba}, agent technique à la Mairie de Douala 3\textsuperscript{e}, au Maire de la Commune, pour un avancement d'échelon.

\begin{enumerate}
    \item \textit{Objet : Demande d'avancement d'échelon.}
    \item \textit{Monsieur le Maire,}
    \item \textit{Je vous prie d'agréer, Monsieur le Maire, l'expression de ma considération distinguée.}
    \item \textit{Fait à Douala, le 15 octobre 2024.}
    \item \textit{Par la présente, je sollicite votre bienveillance pour un avancement à l'échelon supérieur, conformément aux dispositions du Statut Général de la Fonction Publique.}
    \item \textit{Au vu de mon ancienneté de 10 ans à l'échelon actuel et de ma notation administrative « Très Bien » lors des trois dernières évaluations annuelles (2021, 2022, 2023), je estime remplir les conditions requises.}
    \item \textit{M. Essomba Paul, Agent Technique, Matricule 4455, Service Voirie, BP 88 Douala.}
    \item \textit{À l'attention de Monsieur le Maire de la Commune de Douala 3\textsuperscript{e}, BP 123 Douala.}
    \item \textit{Pièces jointes : 1) Dernières fiches de notation ; 2) Copie du dernier arrêté de nomination.}
    \item \textit{J'ai l'honneur de porter à votre connaissance ma situation administrative.}
\end{enumerate}

\textbf{Travail demandé :}
\begin{enumerate}
    \item Remettez ces dix éléments dans l'ordre logique de la structure externe et interne de la requête (numérotez de 1 à 10).
    \item Justifiez l'ordre retenu en identifiant pour chaque élément sa fonction (ex : « Élément 7 : Entête -- Expéditeur », « Élément 6 : Corps -- Arguments »).
    \item Relevez dans l'élément 6 \textbf{deux maladresses de langue} (syntaxe/lexique) et proposez la correction adaptée au style administratif.
\end{enumerate}
\end{exercice}

\begin{exercice}[Remédiation : Correction ciblée et réécriture (15 pts)]
\textbf{Sujet :} La requête ci-dessous, adressée au Directeur d'une entreprise publique par un agent pour un congé annuel, contient \textbf{dix erreurs ciblées} (formules, ponctuation, lexique, structure, anglicismes/camfranglais).

\textit{« Monsieur le Directeur Général,
Objet : Demande de congé annual.
Monsieur,
Je viens par la présente vous demander mon congé de cette année. J'ai pas pris de congé l'année dernière parce que j'étais sur un projet important. Maintenant je veux partir au village voir ma famille pendant 30 jours du 15 juillet au 15 aout. C'est urgent.
Jointe : mon planning.
Salutations.
Dupont Jean, Agent, Matricule 123. »}

\textbf{Travail demandé :}
\begin{enumerate}
    \item Identifiez et numérotez les \textbf{dix erreurs} (relevez le passage fautif et citez la catégorie d'erreur : Formule, Ponctuation, Lexique/Registre, Structure, Orthographe/Grammaire).
    \item Proposez pour chaque erreur la \textbf{correction} attendue.
    \item Rédigez la \textbf{version définitive corrigée} de la requête, prête à être signée et envoyée.
\end{enumerate}
\end{exercice}



\newpage

\coursec{Corrigés des exercices}

\coursec{Les textes fonctionnels et les actes de langage}

\begin{corrige}[Exercice 1 — QCM : Actes de langage et structure de la requête]
\begin{enumerate}
    \item \textbf{B.} La requête vise à obtenir quelque chose : l'acte \cle{illocutoire} principal est une \cle{directive} (demander, solliciter, requérir). L'assertion se contente de décrire, la commissive engage celui qui parle, l'expressive manifeste un sentiment : aucune de ces trois fonctions ne correspond à la finalité d'une requête.
    \item \textbf{B.} L'entête mentionne obligatoirement la \cle{qualité} de l'expéditeur (ex : parent d'élève, agent) et celle du destinataire (ex : Monsieur le Proviseur), avec leurs coordonnées. Le numéro de téléphone personnel et la date de naissance ne sont pas exigés ; le matricule de l'élève figure plutôt dans l'objet ou le corps.
    \item \textbf{B.} L'objet est une \cle{phrase nominale} concise, sans verbe conjugué ni point final : « Objet : demande de réintégration en externat. » Une phrase complète ou un paragraphe alourdirait la mention.
    \item \textbf{C.} L'effet \cle{perlocutoire} visé est la décision favorable du destinataire (réparation, correction, nomination). La simple lecture (A) ou la compréhension (B) relèvent de la réception, pas de l'effet recherché ; l'archivage (D) est neutre.
\end{enumerate}
\textbf{Points de vigilance :} bien distinguer les trois niveaux de l'acte de langage : \emph{locutoire} (ce qui est dit), \emph{illocutoire} (ce que l'on fait en le disant : demander), \emph{perlocutoire} (l'effet obtenu : la décision).
\end{corrige}

\begin{corrige}[Exercice 3 — Définitions et identification : Textes fonctionnels et iconiques]
\begin{enumerate}
    \item Un \cle{texte fonctionnel} est un texte utilitaire produit pour répondre à une situation de communication pratique de la vie courante ou professionnelle. Deux caractéristiques le distinguent du texte littéraire : (a) il vise un \textbf{objectif précis et immédiat} (informer, demander, réclamer, rendre compte) et non le plaisir esthétique ; (b) il obéit à des \textbf{conventions de forme fixes} (structure codifiée, formules rituelles, registre soutenu), alors que le texte littéraire privilégie la création, la subjectivité et la liberté formelle.
    \item Structure externe (trois éléments parmi les suivants) : l'\textbf{entête} (coordonnées de l'expéditeur, qualité et adresse du destinataire, lieu et date), l'\textbf{objet}, la mention des \textbf{pièces jointes}, la \textbf{signature}.
    \item Structure interne : l'\textbf{introduction} (formule d'appel + phrase d'amorce), le \textbf{corps} (exposé des faits, arguments, demande précise), la \textbf{conclusion} (formule de transition + formule de politesse finale).
    \item Pour une contestation de note, la voie hiérarchique impose de s'adresser d'abord au \textbf{Professeur Principal} (qui vérifie la note auprès de l'enseignant concerné) ; à défaut de solution, on remonte au \textbf{Censeur}, puis au \textbf{Proviseur}. La requête formelle, en dernier recours, est adressée au \textbf{Proviseur}, autorité administrative de l'établissement. Respecter la voie hiérarchique est une condition de recevabilité de la requête.
\end{enumerate}
\end{corrige}

\coursec{La structure externe et interne de la requête}

\begin{corrige}[Exercice 4 — Calculs directs et repérage : Taux d'erreurs et structure]
\begin{enumerate}
    \item Erreurs « structurelles et lexicales » : formules ($3$) + lexique ($5$) + structure ($2$) $= 10$ erreurs. Total des erreurs relevées : $3 + 5 + 2 + 4 = 14$.
    \[ \text{Taux} = \frac{10}{14} \times 100 = \frac{1000}{14} \approx 71,4\,\% \]
    Les erreurs de forme (structure et lexique) représentent donc environ \textbf{71,4\,\%} des erreurs, contre $4/14 \approx 28,6\,\%$ pour l'orthographe.
    \item Pénalités de structure/lexique : $10 \times 0,5 = 5$ points. Pénalités d'orthographe : $4 \times 0,25 = 1$ point. Total des pénalités : $5 + 1 = 6$ points.
    \[ \text{Note} = 20 - 6 = \textbf{14/20} \]
    \item Éléments appartenant au \textbf{corps de la requête} : \textbf{Exposé des faits ; Arguments ; Demande précise}. Les autres se répartissent ainsi : structure externe (Expéditeur, Pièces jointes, Signature, Objet) ; introduction (Formule d'appel).
\end{enumerate}
\textbf{Point de vigilance :} ne pas confondre le taux demandé (par rapport au total des erreurs, soit 14) avec un taux par rapport au nombre de mots (180), qui n'intervient pas ici.
\end{corrige}

\begin{corrige}[Exercice 5 — Rédaction de l'entête et de l'objet]
\textbf{1. Entête complet :}
\begin{quote}
Monsieur \textsc{Kamga} Joseph,\\
Parent d'élève,\\
Quartier Banengo, BP 124, Bafoussam\\[0.5em]
\hspace*{6cm}\textbf{À Monsieur le Proviseur}\\
\hspace*{6cm}du Lycée Technique de Bafoussam\\
\hspace*{6cm}Bafoussam\\[0.5em]
\hspace*{6cm}Bafoussam, le \ldots\ldots\ldots\ldots\ldots\ldots
\end{quote}
\textbf{2. Objet :}
\begin{quote}
\textbf{Objet :} Demande de révision de l'affectation en internat de l'élève \textsc{Kamga} Paul, Terminale F4, matricule 2024-0451.
\end{quote}
\textbf{Points de vigilance :} l'expéditeur figure en haut à gauche avec sa qualité (« parent d'élève ») ; le destinataire en haut à droite, précédé de « À » ou « À l'attention de » ; le lieu et la date se placent à droite, sous le destinataire. L'objet est une \cle{phrase nominale} (aucun verbe conjugué, pas de point exclamatif) ; la mention du matricule est admise car elle facilite le traitement du dossier.
\end{corrige}

\begin{corrige}[Exercice 6 — Rédaction de l'introduction et de la conclusion]
\textbf{1. Introduction (formule d'appel + amorce) :}
\begin{quote}
Monsieur le Proviseur,

J'ai l'honneur de m'adresser à votre haute bienveillance en ma qualité de père de l'élève \textsc{Kamga} Paul, inscrit en classe de Terminale F4 sous le matricule 2024-0451, au sujet de la décision d'affectation en internat notifiée à mon fils pour l'année scolaire en cours.
\end{quote}
L'amorce remplit trois fonctions : situer la \textbf{qualité de l'expéditeur} (« en ma qualité de père de l'élève... »), identifier le \textbf{destinataire} par sa fonction (« votre haute bienveillance »), et annoncer la \textbf{référence de la décision contestée}.
\par\textbf{2. Conclusion (transition + demande + formule de politesse) :}
\begin{quote}
En conséquence, je sollicite de votre bienveillance la révision de cette décision et la réintégration de mon fils au régime externe.

Dans l'attente d'une suite que j'ose espérer favorable, je vous prie d'agréer, Monsieur le Proviseur, l'expression de ma haute considération.
\end{quote}
\textbf{Points de vigilance :} la demande peut aussi être formulée avec un subjonctif (« je demande qu'il soit procédé à la révision... ») ou un infinitif (« je sollicite de votre bienveillance la révision... ») ; la formule finale rappelle obligatoirement le titre du destinataire entre virgules, sans jamais écrire « l'expression de mes sentiments » à un supérieur hiérarchique (réserver « sentiments » aux égaux ou aux formules protocolaires précises).
\end{corrige}

\begin{corrige}[Exercice 7 — Rédaction du corps de la requête : Faits, Arguments, Demande]
\textbf{1. Exposé des faits (chronologie) :}
\begin{quote}
Mon fils, \textsc{Kamga} Paul, est inscrit au Lycée Technique de Bafoussam depuis la classe de Seconde, où il a toujours suivi sa scolarité au régime externe. À la publication des listes d'affectation pour l'année scolaire en cours, j'ai constaté, avec surprise, qu'il figure sur la liste des internes, sans qu'aucune demande en ce sens n'ait été formulée de ma part.
\end{quote}
\textbf{2. Deux arguments développés :}
\begin{quote}
D'une part, notre domicile familial est situé à environ 500 mètres de l'établissement, ce qui rend l'internat inutile : mon fils peut se rendre au lycée à pied en quelques minutes, dans des conditions de sécurité optimales.

D'autre part, les résultats scolaires de mon fils, constamment satisfaisants depuis la Seconde au régime externe, démontrent que ce régime est favorable à sa réussite ; un changement brutal de cadre de vie risquerait de perturber ses études en cette année décisive du Baccalauréat.
\end{quote}
\textbf{3. Demande précise :}
\begin{quote}
Au vu de ces éléments et des pièces justificatives ci-jointes, je sollicite de votre bienveillance la réintégration de mon fils au régime externe et la modification en conséquence de la liste d'affectation. Je vous saurais gré de bien vouloir y faire droit.
\end{quote}
\textbf{Points de vigilance :} les faits sont rapportés objectivement, au passé composé ou au présent, sans injure ni dramatisation ; les arguments sont introduits par des connecteurs (« d'une part... d'autre part ») ; la demande emploie le lexique administratif exigé (« sollicite », « au vu de », « y faire droit », « pièces justificatives ») et aucune trace de camfranglais.
\end{corrige}

\coursec{Le lexique spécialisé et l'emprunt}

\begin{corrige}[Exercice 2 — Vrai / Faux justifié : Lexique spécialisé et emprunts]
\begin{enumerate}
    \item \textbf{Faux.} La requête exige le \cle{registre soutenu} et le \cle{lexique administratif} ; le camfranglais est un registre familier à proscrire absolument dans un écrit officiel. Correction attendue : « J'ai l'honneur de solliciter, par la présente, une audience en faveur de mon fils. »
    \item \textbf{Vrai.} Placé \emph{après} le nom, « ci-joint » fonctionne comme un adjectif et s'accorde : « les pièces ci-jointes », « la copie ci-jointe ». En revanche, placé avant le nom ou en tête de phrase, il reste invariable : « Ci-joint les pièces demandées », « Vous trouverez ci-joint copie de... ».
    \item \textbf{Vrai.} « Planning » est un \cle{emprunt} à l'anglais (anglicisme) ; en français soutenu, on lui substitue « calendrier », « programme » ou « échéancier ». Dans un écrit formel, l'emprunt non assimilé doit être évité.
    \item \textbf{Vrai.} C'est la formule de politesse finale classique de la correspondance administrative : verbe de politesse (« je vous prie d'agréer »), rappel du titre du destinataire entre virgules (« Monsieur le Directeur »), et formule de considération (« l'expression de mes salutations distinguées »).
\end{enumerate}
\end{corrige}

\coursec{Rédaction guidée : entête, introduction, corps et conclusion}
\courssub{Exploitation des textes iconiques pour préparer la requête}

\begin{corrige}[Exercice 8 — Analyse de texte iconique : Organigramme et destinataire]
\begin{enumerate}
    \item Le destinataire direct est le \textbf{Chef du Service Maintenance}, supérieur hiérarchique immédiat de M. Ndongo : la requête doit d'abord transiter par lui (visa ou avis), conformément à la voie hiérarchique.
    \item Le destinataire final, décisionnaire en matière de primes, est le \textbf{Directeur des Ressources Humaines (DRH)}, sous l'autorité du Directeur Général.
    \item Entête et objet de la requête adressée au destinataire direct :
\begin{quote}
Monsieur \textsc{Ndongo}, Agent de Maintenance,\\
Matricule 8842, Service Maintenance, SONARA, Limbé\\[0.5em]
\hspace*{6cm}\textbf{À Monsieur le Chef du Service Maintenance}\\
\hspace*{6cm}SONARA, Limbé\\[0.5em]
\hspace*{6cm}Limbé, le \ldots\ldots\ldots\ldots\ldots\ldots\\[0.5em]
\textbf{Objet :} Demande de régularisation de la prime de risque non perçue depuis trois mois.
\end{quote}
\end{enumerate}
\textbf{Point de vigilance :} l'organigramme est un \cle{texte iconique} : il se lit de haut en bas, chaque flèche indiquant un lien de subordination. Sauter un échelon hiérarchique (écrire directement au DRH) expose la requête au rejet pour irrecevabilité.
\end{corrige}

\coursec{Analyse de situation, plan détaillé et rédaction complète}

\begin{corrige}[Exercice 9 — Type Baccalauréat Technique : Rédaction complète de requête]
\textbf{Requête modèle :}
\begin{quote}
Monsieur \textsc{Fouda} André,\\
Commerçant au Marché Central,\\
BP 356, Bafoussam\\[0.5em]
\hspace*{6cm}\textbf{À Monsieur le Proviseur}\\
\hspace*{6cm}du Lycée Technique de Bafoussam\\
\hspace*{6cm}Bafoussam\\[0.5em]
\hspace*{6cm}Bafoussam, le \ldots\ldots\ldots\ldots\ldots\ldots\\[0.5em]
\textbf{Objet :} Demande de réintégration au régime externe de l'élève \textsc{Fouda} Grâce, Terminale G2, matricule 2024-1122.\\[0.3em]
\textbf{Pièces jointes :} 1) Certificat médical ; 2) Attestation d'hébergement.\\[0.8em]
Monsieur le Proviseur,

J'ai l'honneur de m'adresser à votre haute bienveillance en ma qualité de père de l'élève \textsc{Fouda} Grâce, inscrite en classe de Terminale G2 sous le matricule 2024-1122, au sujet de son affectation en internat pour l'année scolaire 2024-2025.

En effet, à la publication des listes d'affectation, j'ai constaté que ma fille figure parmi les élèves internes, alors qu'aucune demande en ce sens n'a été formulée de ma part et qu'elle a toujours étudié au régime externe.

Cette affectation est incompatible avec l'état de santé de mon enfant. D'une part, elle souffre d'asthme sévère, pathologie chronique qui exige un environnement familial strict et des soins nocturnes quotidiens, comme l'atteste le certificat médical du Dr \textsc{Nguyen}, pneumologue à l'Hôpital Régional de Bafoussam, ci-joint. D'autre part, le coût de l'internat, qui s'élève à 150 000 FCFA par an, excède largement mes modestes ressources de commerçant, alors que notre domicile se situe à proximité immédiate de l'établissement.

En conséquence, je sollicite de votre bienveillance la révision de cette décision et la réintégration de ma fille au régime externe, afin de préserver sa santé et sa réussite scolaire en cette année d'examen.

Dans l'attente d'une suite que j'ose espérer favorable, je vous prie d'agréer, Monsieur le Proviseur, l'expression de ma haute considération.\\[0.8em]
\hspace*{8cm}\textsc{Fouda} André\\
\hspace*{8cm}(Signature)
\end{quote}
\textbf{Barème indicatif (20 pts) :}
\begin{itemize}
    \item Structure externe (4 pts) : entête expéditeur/destinataire complet et bien disposé (2 pts) ; lieu et date (0,5 pt) ; objet en phrase nominale (1 pt) ; pièces jointes et signature (0,5 pt).
    \item Structure interne (8 pts) : formule d'appel + amorce (2 pts) ; exposé des faits chronologique et objectif (2 pts) ; deux arguments développés et pertinents (santé, finances) avec connecteurs (2 pts) ; demande précise + conclusion avec formule de politesse complète (2 pts).
    \item Langue (8 pts) : lexique administratif soutenu (« sollicite », « au vu de », « ci-joint »), absence d'anglicismes et de camfranglais, correction syntaxique, orthographe et ponctuation, accords.
\end{itemize}
\textbf{Points de vigilance :} ne jamais tutoyer le destinataire ni employer « je veux » ; citer les pièces jointes dans le corps (« comme l'atteste le certificat ci-joint ») ; formuler la demande avec un infinitif ou un subjonctif ; rappeler le titre du destinataire dans la formule finale.
\end{corrige}

\begin{corrige}[Exercice 10 — Type Probatoire : Analyse de requête et remise en ordre]
\textbf{1. et 2. Ordre logique et fonction de chaque élément :}
\begin{center}
\begin{tabularx}{\linewidth}{|c|c|X|}\hline
\rowcolor{popblueL} \textbf{Ordre} & \textbf{Élément} & \textbf{Fonction} \\ \hline
1 & 7 & Entête : expéditeur (nom, qualité, matricule, adresse) \\ \hline
2 & 8 & Entête : destinataire (qualité, adresse) \\ \hline
3 & 4 & Entête : lieu et date \\ \hline
4 & 1 & Objet (phrase nominale) \\ \hline
5 & 9 & Mention des pièces jointes \\ \hline
6 & 2 & Introduction : formule d'appel \\ \hline
7 & 10 & Introduction : phrase d'amorce (annonce du sujet) \\ \hline
8 & 6 & Corps : arguments (ancienneté, notation) \\ \hline
9 & 5 & Corps/conclusion : demande précise \\ \hline
10 & 3 & Conclusion : formule de politesse finale \\ \hline
\end{tabularx}
\end{center}
\textbf{Justification de l'ordre :} la structure externe (éléments 7, 8, 4, 1, 9) précède toujours la structure interne ; dans celle-ci, la formule d'appel (2) ouvre obligatoirement le texte, l'amorce (10) annonce le sujet avant les arguments (6), et la demande (5) ne peut précéder les arguments qui la fondent ; la formule de politesse (3) clôt l'ensemble.
\par\textbf{3. Maladresses de l'élément 6 et corrections :}
\begin{itemize}
    \item « \emph{je estime} » : faute d'élision. Correction : « \textbf{j'estime} remplir les conditions requises ».
    \item « \emph{ma notation administrative « Très Bien »} » : maladresse de registre et de formulation. Correction : « mon \textbf{appréciation générale "Très Bien" portée sur mes fiches de notation} » ou « mes fiches de notation portant l'appréciation "Très Bien" ». On peut également renforcer l'amorce : « \emph{Au vu de mon ancienneté...} » devient « \textbf{Justifiant d'une ancienneté de dix ans à l'échelon actuel...} ».
\end{itemize}
\textbf{Barème indicatif (15 pts) :} ordre correct (5 pts, 0,5 pt par position) ; fonctions correctement identifiées (5 pts) ; deux maladresses relevées et corrigées (5 pts).
\end{corrige}

\coursec{Confrontation, amélioration des productions, compte rendu et remédiation}

\begin{corrige}[Exercice 11 — Remédiation : Correction ciblée et réécriture]
\textbf{1. et 2. Relevé des dix erreurs et corrections :}
\begin{center}
\begin{tabularx}{\linewidth}{|c|X|X|}\hline
\rowcolor{popblueL} \textbf{N°} & \textbf{Erreur (catégorie)} & \textbf{Correction} \\ \hline
1 & Absence d'entête : expéditeur, destinataire, lieu et date manquants (Structure) & Rédiger un entête complet : qualité et adresse de l'expéditeur en haut à gauche, qualité et adresse du destinataire à droite, lieu et date \\ \hline
2 & « congé annual » (Orthographe) : faute sur l'adjectif & « congé \textbf{annuel} » \\ \hline
3 & « Je viens par la présente vous demander » (Lexique/Registre) : tournure familière & « J'ai l'honneur de solliciter, par la présente, » \\ \hline
4 & « mon congé de cette année » (Lexique) : imprécision & « mon congé annuel au titre de l'année en cours » \\ \hline
5 & « J'ai pas pris » (Grammaire/Registre) : négation familière incomplète & « Je \textbf{n'ai pas} pris » \\ \hline
6 & « parce que j'étais sur un projet important » (Registre) : tournure relâchée & « étant donné que j'étais mobilisé sur un projet prioritaire » \\ \hline
7 & « Maintenant je veux partir » (Lexique/Registre) : « je veux » inadapté à une demande hiérarchique & « Je souhaite à présent bénéficier de » \\ \hline
8 & « 15 aout » (Orthographe) : accent omis & « 15 \textbf{août} » \\ \hline
9 & « Jointe : mon planning. » (Lexique + Grammaire) : anglicisme et formulation fautive & « \textbf{Pièce jointe :} mon \textbf{calendrier} de congé. » \\ \hline
10 & « Salutations. » (Formule) : clôture familière inadaptée & « Je vous prie d'agréer, Monsieur le Directeur Général, l'expression de ma considération distinguée. » \\ \hline
\end{tabularx}
\end{center}
\textbf{3. Version définitive corrigée :}
\begin{quote}
Monsieur \textsc{Dupont} Jean, Agent,\\
Matricule 123,\\
BP \ldots, Yaoundé\\[0.5em]
\hspace*{6cm}\textbf{À Monsieur le Directeur Général}\\
\hspace*{6cm}de \ldots\ldots\ldots\ldots\\
\hspace*{6cm}Yaoundé\\[0.5em]
\hspace*{6cm}Yaoundé, le \ldots\ldots\ldots\ldots\ldots\ldots\\[0.5em]
\textbf{Objet :} Demande de congé annuel.\\[0.3em]
\textbf{Pièce jointe :} calendrier de congé.\\[0.8em]
Monsieur le Directeur Général,

J'ai l'honneur de solliciter, par la présente, l'autorisation de bénéficier de mon congé annuel au titre de l'année en cours.

En effet, je n'ai pas pris de congé l'année dernière, étant donné que j'étais mobilisé sur un projet prioritaire du service. Je souhaite à présent rejoindre ma famille au village pour une durée de trente (30) jours, du 15 juillet au 15 août.

Dans l'attente d'une suite favorable, je vous prie d'agréer, Monsieur le Directeur Général, l'expression de ma considération distinguée.\\[0.8em]
\hspace*{8cm}\textsc{Dupont} Jean\\
\hspace*{8cm}(Signature)
\end{quote}
\textbf{Barème indicatif (15 pts) :} identification des dix erreurs avec catégorie (5 pts, 0,5 pt par erreur) ; corrections proposées (5 pts) ; version définitive conforme (structure externe et interne, registre soutenu, orthographe) (5 pts).
\par\textbf{Points de vigilance :} en remédiation, toujours procéder en trois temps (relever, corriger, réécrire) ; vérifier dans la version finale qu'aucune erreur corrigée ne subsiste ; la formule de clôture doit reprendre exactement le titre du destinataire annoncé dans l'entête.
\end{corrige}

\newpage

\coursec{Fiche bilan}

\begin{popfigure}
\begin{tikzpicture}[mindmap, grow cyclic, every node/.style={concept, font=\small},
  root concept/.append style={concept color=popblue, font=\bfseries},
  level 1/.append style={level distance=3.4cm, sibling angle=60},
  level 1 concept/.append style={font=\scriptsize}]
\node[root concept]{La requête}
  child[concept color=poporange]{ node {Actes de langage : locutoire, illocutoire, perlocutoire} }
  child[concept color=popgreen]{ node {Structure externe : entête, objet, formules, date, signature} }
  child[concept color=poppurple]{ node {Structure interne : introduction, corps, conclusion} }
  child[concept color=poppink]{ node {Lexique spécialisé et emprunts} }
  child[concept color=popgold]{ node {Textes fonctionnels et iconiques : lire avant d'écrire} }
  child[concept color=popblueL]{ node {Méthode : analyse, plan, rédaction, relecture} };
\end{tikzpicture}
\legende{Carte mentale de la leçon : les six savoirs essentiels autour de la production d'une requête.}
\end{popfigure}

\begin{bilan}
\textbf{Définitions clés}
\begin{itemize}
  \item \cle{Requête} : texte fonctionnel par lequel un citoyen (le \cle{requérant}) demande à une autorité administrative un droit, un service ou une faveur.
  \item \cle{Acte locutoire} : le fait de produire un énoncé (ce qui est dit).
  \item \cle{Acte illocutoire} : la valeur de l'énoncé (demander, ordonner, promettre, remercier).
  \item \cle{Acte perlocutoire} : l'effet produit sur le destinataire (convaincre, obtenir, émouvoir).
  \item \cle{Lexique spécialisé} : vocabulaire propre à un domaine (administration, justice, technique).
  \item \cle{Emprunt} : mot pris à une langue étrangère et intégré au français (\emph{budget}, \emph{match}).
\end{itemize}

\textbf{Structure de la requête (à connaître par cœur)}
\begin{center}
\begin{tabularx}{\linewidth}{|l|X|}
\hline
\rowcolor{popblueL} \textbf{Partie} & \textbf{Contenu} \\
\hline
Entête & Nom et coordonnées du requérant (à gauche), destinataire avec sa fonction (à droite), lieu et date, \textbf{objet} précis. \\
\hline
Introduction & Formule de politesse + présentation du requérant et de sa situation. \\
\hline
Corps & Exposé clair et ordonné des faits, des motifs et de la demande précise ; pièces jointes éventuelles. \\
\hline
Conclusion & Remerciements, formule de haute considération, signature. \\
\hline
\end{tabularx}
\end{center}

\textbf{Méthodes express}
\begin{itemize}
  \item Analyser la situation : \emph{qui} demande, \emph{à qui}, \emph{quoi}, \emph{pourquoi}.
  \item Bâtir un plan détaillé avant de rédiger ; une idée par paragraphe.
  \item Employer un ton respectueux, des phrases courtes, le lexique administratif.
  \item Relire : orthographe, accords, cohérence, formules réglementaires.
\end{itemize}
\end{bilan}

\begin{aretenir}[Checklist avant le Bac]
\begin{itemize}
  \item $\square$ Je sais définir la requête et citer ses qualités (clarté, précision, courtoisie).
  \item $\square$ Je distingue les trois actes de langage : locutoire, illocutoire, perlocutoire.
  \item $\square$ Je place correctement l'entête : requérant à gauche, destinataire à droite.
  \item $\square$ Je rédige un objet bref et précis (une seule ligne).
  \item $\square$ Je connais les formules d'introduction et de conclusion réglementaires.
  \item $\square$ J'expose les faits dans un ordre logique (chronologique ou par importance).
  \item $\square$ Je formule une demande explicite et réalisable.
  \item $\square$ J'emploie le lexique spécialisé de l'administration à bon escient.
  \item $\square$ Je reconnais un emprunt et je sais l'intégrer correctement au français.
  \item $\square$ Je sais exploiter un texte iconique (tableau, schéma) pour appuyer ma demande.
  \item $\square$ Je fais un plan détaillé avant toute rédaction.
  \item $\square$ Je relis et j'améliore ma production (forme et fond) avant de rendre.
\end{itemize}
\end{aretenir}

\findecours

\end{document}
